\documentclass[11pt,oneside]{article}

\usepackage[a4paper,margin=1.08in]{geometry}
\usepackage[T1]{fontenc}
\usepackage[utf8]{inputenc}
\usepackage{lmodern}
\usepackage{microtype}
\usepackage{mathpazo}
\usepackage{amsmath,amssymb,amsthm,mathtools,mathrsfs}
\usepackage{enumitem}
\usepackage{tikz-cd}
\usepackage{array,booktabs}
\usepackage{xcolor}
\usepackage{hyperref}
\usepackage[nameinlink,noabbrev]{cleveref}
\usepackage{csquotes}
\usepackage{stmaryrd}
\usepackage{titlesec}

\hypersetup{
  colorlinks=true,
  linkcolor=blue!50!black,
  citecolor=blue!50!black,
  urlcolor=blue!50!black,
  pdftitle={Crit\`eres de platitude et de projectivit\'e --- Notes in an IHES-style},
  pdfauthor={}
}

\titleformat{\section}{\large\bfseries}{\thesection.}{0.6em}{}
\titleformat{\subsection}{\normalsize\bfseries}{\thesubsection.}{0.5em}{}
\titleformat{\subsubsection}{\normalsize\itshape}{\thesubsubsection.}{0.5em}{}

\setlist[itemize]{leftmargin=2.0em}
\setlist[enumerate]{leftmargin=2.2em}

\newtheorem{theorem}{Theorem}[section]
\newtheorem{proposition}[theorem]{Proposition}
\newtheorem{lemma}[theorem]{Lemma}
\newtheorem{corollary}[theorem]{Corollary}
\newtheorem{definition}[theorem]{Definition}
\newtheorem{remark}[theorem]{Remark}
\newtheorem{example}[theorem]{Example}

\newcommand{\A}{\mathcal A}
\newcommand{\B}{\mathcal B}
\newcommand{\C}{\mathcal C}
\newcommand{\D}{\mathcal D}
\newcommand{\F}{\mathcal F}
\newcommand{\G}{\mathcal G}
\newcommand{\HH}{\mathcal H}
\newcommand{\I}{\mathcal I}
\newcommand{\Lcal}{\mathcal L}
\newcommand{\M}{\mathcal M}
\newcommand{\N}{\mathcal N}
\newcommand{\Pcal}{\mathcal P}
\newcommand{\OO}{\mathcal O}
\newcommand{\Q}{\mathcal Q}
\newcommand{\Tcal}{\mathcal T}
\newcommand{\Kcal}{\mathcal K}
\newcommand{\Spec}{\operatorname{Spec}}
\newcommand{\Supp}{\operatorname{Supp}}
\newcommand{\Ass}{\operatorname{Ass}}
\newcommand{\Tor}{\operatorname{Tor}}
\newcommand{\Hom}{\operatorname{Hom}}
\newcommand{\Ext}{\operatorname{Ext}}
\newcommand{\Ann}{\operatorname{Ann}}
\newcommand{\Image}{\operatorname{Im}}
\newcommand{\Ker}{\operatorname{Ker}}
\newcommand{\Coker}{\operatorname{Coker}}
\newcommand{\codim}{\operatorname{codim}}
\newcommand{\gr}{\operatorname{gr}}
\newcommand{\rk}{\operatorname{rk}}
\newcommand{\id}{\mathrm{id}}
\newcommand{\into}{\hookrightarrow}
\newcommand{\onto}{\twoheadrightarrow}
\newcommand{\ot}{\otimes}
\newcommand{\ox}[1]{\mathcal O_{#1}}
\newcommand{\ks}[1]{k(#1)}
\newcommand{\et}{\text{\rm\'et}}
\newcommand{\etale}{\'etale}
\newcommand{\fp}{\mathrm{fp}}
\newcommand{\op}{\mathrm{op}}
\newcommand{\st}{\operatorname{str}}
\newcommand{\Assrel}{\operatorname{Ass}_{X/S}}
\newcommand{\Sch}{\mathbf{Sch}}
\newcommand{\Ens}{\mathbf{Ens}}

\begin{document}

\begin{center}
  {\Large\bfseries Crit\`eres de platitude et de projectivit\'e}\\[0.4em]
  {\large Techniques de \enquote{platification} d'un module}\\[0.9em]
  {\normalsize based on the paper of Michel Raynaud and Laurent Gruson}\\[0.8em]
  {\normalsize Notes prepared for a seminar IMATE}\\[0.35em]
  {\normalsize Autor: C\'esar Galindo}
\end{center}

\begin{abstract}
These notes offer a detailed and more narrative account of the central architecture of Raynaud--Gruson's paper on flatness, projectivity, purity, universal flattening, and flattening by blow-up.  They are inspired in part by a lecture delivered by P.~Pel\'aez at IMATE.  The aim is not merely to summarize the statements, but to reorganize the main proofs into a coherent exposition for a seminar at IMATE, so that the movement from local algebra to birational geometry can be followed as a single mathematical narrative.  We explain the role of relative d\'evissage, prove the flatness criterion from the d\'evissage machine, derive the characterization of projectivity in terms of flatness and purity, establish the representability of the universal flattening condition in the local-henselian setting, and then present the proof mechanism of flattening by admissible blow-up.  Particular care is devoted to the homological meaning of universal injectivity, to the geometric role of the relative assassin, and to the way strict transform removes precisely the torsion created over the bad locus.
\end{abstract}

\tableofcontents
\setlength{\parskip}{0.35em}
\linespread{1.135}\selectfont

\section{Introduction}

Let $A\to B$ be a ring homomorphism and let $M$ be a $B$-module.  The question whether $M$ is flat over $A$ is one of the central problems of algebraic geometry.  In a na\"ive approach one studies the local criterion of flatness through the vanishing of
\[
\Tor_1^A(M,k(\mathfrak p))
\]
for prime ideals $\mathfrak p\subset A$.  This criterion is decisive, but it does not by itself explain the geometry of the bad locus, nor does it say how flatness interacts with projectivity, nor how one can modify the base to improve the situation.

The decisive achievement of Raynaud and Gruson is to organize these issues into one coherent machine.  Read from beginning to end, the paper has a clear narrative arc: it starts from local algebra, passes through a geometric analysis of supports and associated points, and ends with a birational procedure that repairs the defect of flatness.  The ingredients are the following.
\begin{enumerate}[label=(\roman*)]
  \item A \emph{relative d\'evissage} reducing a finite type module to a finite chain of free modules and successive quotients of smaller relative dimension.
  \item A \emph{criterion of flatness} saying that flatness is equivalent to universal injectivity of certain transition maps together with flatness of the terminal quotient.
  \item A geometric notion of \emph{purity}, expressed through associated points and the relative assassin, which upgrades flatness to projectivity.
  \item A representability theorem for the functor of base changes on which a quotient becomes flat or bijective.
  \item A theorem of \emph{flattening by blow-up}, asserting that after a suitable admissible blow-up the strict transform becomes flat.
\end{enumerate}

These notes expand the proof spine of that theory.  We retain the notation and terminology of the original paper as much as possible, but the exposition is rewritten so that the logic of the constructions becomes transparent.

\section{Preliminaries and notational conventions}

Throughout, $S$ denotes a scheme and $f\colon X\to S$ a morphism locally of finite presentation unless otherwise stated.  If $s\in S$, then $k(s)$ denotes the residue field.  For an $\OO_X$-module $\M$ we write $\M_s$ for the pullback of $\M$ to the fiber $X_s:=X\times_S\Spec k(s)$.

All sheaves considered are quasi-coherent and of finite type when needed.  By an \emph{elementary \etale neighborhood} of a pointed morphism $(X,x)\to (S,s)$ we mean a commutative square
\[
\begin{tikzcd}
(Y,y) \ar[r] \ar[d] & (T,t) \ar[d]\\
(X,x) \ar[r] & (S,s)
\end{tikzcd}
\]
whose vertical arrows are \etale and which satisfies the familiar local form of the main theorem of Zariski: $Y\to T$ is finite, the fibers are geometrically integral in a suitable neighborhood, and the chosen points are the unique points over the closed base points.  This local model is the geometric input that makes relative d\'evissage possible.

We shall repeatedly use the following elementary homological facts.

\begin{lemma}[Flatness and short exact sequences]\label{lem:ses-flat}
Let $A$ be a ring and
\[
0\to L\to N\to P\to 0
\]
an exact sequence of $A$-modules.  If $L$ and $P$ are flat over $A$, then $N$ is flat over $A$.  If $N$ is flat over $A$ and $L\to N$ remains injective after arbitrary base change, then $P$ is flat over $A$.
\end{lemma}

\begin{proof}
The first assertion follows by tensoring with any $A$-module and using the long exact Tor sequence; the second follows by the same sequence because universal injectivity gives the exactness of
\[
0\to L\ot_A A'\to N\ot_A A'\to P\ot_A A'\to 0
\]
for every $A$-algebra $A'$, and then $\Tor_1^A(P,-)=0$.
\end{proof}

\begin{definition}[Universal injectivity]
A morphism $u\colon M\to N$ of $A$-modules is said to be \emph{universally injective} if for every $A$-algebra $A'$ the base changed map
\[
u\ot_A A'\colon M\ot_A A'\longrightarrow N\ot_A A'
\]
is injective.
\end{definition}

Universal injectivity is the correct homological avatar of geometric permanence under base change.  It is weaker than split injectivity but strong enough to guarantee flatness of quotients when the ambient module is flat.

\section{Relative d\'evissage}

\subsection{The Zariski-type local geometric preparation}

The first local theorem in the paper is a pointed version of the main theorem of Zariski adapted to the relative setting.

\begin{theorem}[Local finite factorization]\label{thm:zariski-local}
Let $(X,x)\to (S,s)$ be a morphism locally of finite presentation of pointed schemes.  Put
\[
n:=\dim_x(X_s).
\]
Then, after replacing $(X,x)\to(S,s)$ by a pointed elementary \etale neighborhood, there exists a commutative diagram
\[
\begin{tikzcd}
(Y,y) \ar[r] \ar[d] & (T,t) \ar[d]\\
(X,x) \ar[r] & (S,s)
\end{tikzcd}
\]
with the following properties:
\begin{enumerate}[label=\textup{(\alph*)}]
  \item $Y\to T$ is finite and $T\to S$ is smooth of relative dimension $n$;
  \item the fibers of $Y\to T$ are geometrically integral in a neighborhood of $y$;
  \item $y$ is the unique point of $Y$ lying over $x$ and $t$ is the unique point of $T$ lying over $s$.
\end{enumerate}
\end{theorem}

\begin{proof}[Proof sketch with the geometric steps made explicit]
One begins by shrinking $X$ around $x$ so that $X$ is affine over $S$ and all irreducible components of the fiber $X_s$ through $x$ have the same dimension $n$.  Embed $X$ into an affine space over $S$ and choose, by a generic linear projection argument, a morphism to an affine $n$-space over $S$ whose restriction to the relevant component is generically finite.  After shrinking again, this yields a finite morphism from a neighborhood of $x$ to a smooth $S$-scheme $T$ of relative dimension $n$.

The geometric integrality of the fibers is obtained by passing to a suitable open subset where the generic fiber is geometrically integral and then using an elementary \etale neighborhood to isolate the chosen branch.  The uniqueness of the points over $x$ and $s$ is achieved by henselian or pointed localization and by discarding the irrelevant components.

The role of the theorem is not merely to factor $X\to S$ locally.  Its real purpose is to provide a smooth base $T$ over which $X$ becomes finite near the chosen point, because finite morphisms are the natural setting for writing a finite module as a quotient of a free module and then controlling dimensions of supports.
\end{proof}

\subsection{Definition of relative d\'evissage}

We now fix a point $x\in X$ lying over $s\in S$ and an $\OO_X$-module $\M$ of finite type.

\begin{definition}[Relative d\'evissage]\label{def:devissage}
A \emph{relative d\'evissage} of $\M$ at $x$ over $S$ is a finite sequence of pointed schemes
\[
(X_1,x_1)=(X',x')\to (T_1,t_1),\quad (X_2,x_2)\to (T_2,t_2),\quad \dots,\quad (X_r,x_r)\to (T_r,t_r),
\]
together with exact sequences
\[
0\to \Lcal_i\xrightarrow{\alpha_i}\N_i\to \Pcal_i\to 0,
\qquad 1\le i\le r,
\]
such that:
\begin{enumerate}[label=\textup{(\roman*)}]
  \item $\Lcal_i$ is a free $\OO_{T_i}$-module of finite rank;
  \item $\N_i$ and $\Pcal_i$ are $\OO_{T_i}$-modules of finite presentation;
  \item $\N_1$ is the direct image of the pullback of $\M$ to a pointed elementary \etale neighborhood $X'\to X$ with $X'\to T_1$ finite;
  \item for $i<r$, the next $\N_{i+1}$ is obtained from $\Pcal_i$ by repeating the same construction on the support of $\Pcal_i$;
  \item the dimensions
  \[
  n_i:=\dim_{x_i}\bigl(\Supp(\N_i)_{t_i}\bigr)
  \]
  form a strictly decreasing sequence.
\end{enumerate}
If the sequence ends with $n_r\le 0$, or equivalently with a terminal quotient supported in relative dimension $0$, we call it a \emph{total d\'evissage}.
\end{definition}

The idea is extremely simple once the notation is unpacked.  One first makes $X$ finite over a smooth $T_1$; then one chooses a surjection from a free $\OO_{T_1}$-module onto the direct image of $\M$; the kernel and quotient then encode the failure of freeness.  One repeats the same construction for the quotient, but now the support has smaller fiber dimension.  Since fiber dimensions decrease strictly, the process terminates.

\begin{proposition}[Existence of relative d\'evissage]\label{prop:exist-dev}
Let $(X,x)\to(S,s)$ be locally of finite presentation and let $\M$ be an $\OO_X$-module of finite type.  Then, after replacing $(X,x)\to(S,s)$ by a pointed elementary \etale neighborhood, there exists a total relative d\'evissage of $\M$ at $x$ over $S$.
\end{proposition}

\begin{proof}
We argue by induction on
\[
n:=\dim_x\bigl(\Supp(\M)_s\bigr).
\]
If $n<0$, then $\M_x=0$ and there is nothing to prove.

Assume $n\ge 0$.  By \cref{thm:zariski-local}, after replacing $(X,x)$ by a pointed elementary \etale neighborhood $(X',x')$, there exists a finite morphism
\[
\pi\colon X'\to T
\]
with $T\to S$ smooth of relative dimension $n$.  Put
\[
\N_1:=\pi_*(\M|_{X'}).
\]
Because $\pi$ is finite and $\M$ is finite type, $\N_1$ is a finite type $\OO_T$-module.  Shrinking $T$ around $t$ we may assume $\N_1$ is of finite presentation.  Choose a surjection
\[
\Lcal_1\twoheadrightarrow \N_1
\]
with $\Lcal_1$ free of finite rank.  Let $\Pcal_1$ denote the cokernel of the induced map
\[
\alpha_1\colon \Lcal_1\to \N_1.
\]
The point is that by choosing the basis of $\Lcal_1$ appropriately, the image of $\alpha_1$ generates $\N_1$ at the point $t$, and therefore the support of $\Pcal_1$ has fiber dimension strictly smaller than that of $\N_1$ at $t$.

Apply the induction hypothesis to $\Pcal_1$ at the point corresponding to $t$ in its support.  This produces the remainder of the d\'evissage.  The strict decrease of fiber dimensions gives termination.
\end{proof}

\begin{remark}
The content of the proposition is more delicate than the compact proof above may suggest.  The real work lies in showing that the support of $\Pcal_1$ has smaller relative dimension.  This uses Nakayama's lemma at the distinguished point and the fact that one can choose generators of $\N_1\otimes k(t)$ adapted to the fiber dimension at that point.
\end{remark}

\subsection{Transport and composition of d\'evissages}

One of the important technical features of the construction is its functoriality under further pointed \etale localization and under composition.  In particular, if one has a d\'evissage of $\M$ at $x$ and then a d\'evissage of the final quotient at the next point, these may be concatenated into a single d\'evissage.  Likewise, if an \etale neighborhood is replaced by a smaller one, the d\'evissage pulls back.  The original paper uses these compatibility statements repeatedly when descending and ascending through the inductive steps.

\section{The flatness criterion}

\subsection{Statement}

Let $f\colon (X,x)\to(S,s)$ be locally of finite presentation and let $\M$ be a finite type $\OO_X$-module.  Choose a total relative d\'evissage of $\M$ at $x$ over $S$:
\[
(\Lcal_i\xrightarrow{\alpha_i}\N_i\to \Pcal_i)_{1\le i\le r}.
\]
The central theorem of \S2 of the paper identifies the flatness of $\M_x$ with a finite list of conditions on this d\'evissage.

\begin{theorem}[Raynaud--Gruson flatness criterion]\label{thm:flat-criterion}
With the notation above, the following conditions are equivalent.
\begin{enumerate}[label=\textup{(\roman*)}]
  \item $\M_x$ is flat over $\OO_{S,s}$.
  \item For every $i$, the map $\alpha_i$ is universally injective over $\OO_{S,s}$ and the terminal quotient $\Pcal_r$ is flat over $\OO_{S,s}$.
  \item For every $i$, the map $\alpha_i\otimes k(s)$ is injective and $\Tor_1^{\OO_{S,s}}(\M_x,k(s))=0$.
  \item After replacing $S$ by a sufficiently small pointed \etale neighborhood, each $\alpha_i$ remains injective after arbitrary base change and the successive quotients are flat over the corresponding base.
\end{enumerate}
\end{theorem}

The theorem is a structural refinement of the local criterion of flatness.  Instead of testing flatness by a single Tor-vanishing condition, it turns flatness into a finite inductive package.

\subsection{A lemma on universal injectivity}

The basic technical step is the following variant of the criterion from the paper.

\begin{lemma}\label{lem:ui-smooth}
Let $(T,t)\to(S,s)$ be a smooth morphism of pointed schemes with geometrically integral fibers.  Let
\[
\alpha\colon \Lcal\to \N
\]
be a morphism of $\OO_T$-modules with $\Lcal$ free of finite rank and $\N$ of finite presentation.  Assume that the cokernel $\Pcal$ is flat over $S$.  Then the following are equivalent near $t$:
\begin{enumerate}[label=\textup{(\alph*)}]
  \item $\alpha$ is universally injective over $S$;
  \item $\alpha\otimes k(s)$ is injective at the point $t$;
  \item $\alpha$ is injective on a neighborhood of $t$.
\end{enumerate}
\end{lemma}

\begin{proof}
The implication \textup{(a)}\,$\Rightarrow$\,\textup{(b)} is obvious.  We prove \textup{(b)}\,$\Rightarrow$\,\textup{(a)}.  Let $\Kcal:=\Ker(\alpha)$.  Since $\Lcal$ is free, $\Kcal$ is of finite type.  The hypothesis that $\alpha\otimes k(s)$ is injective at $t$ means
\[
\Kcal\otimes k(s)=0
\]
at $t$.  Because the fibers of $T\to S$ are geometrically integral, the support of $\Kcal$ cannot contain a generic point of the fiber through $t$.  On the other hand, the exact sequence
\[
0\to \Kcal\to \Lcal\to \N\to \Pcal\to 0
\]
and the flatness of $\Pcal$ show, after tensoring with arbitrary residue fields of points of $S$, that the formation of $\Kcal$ behaves upper-semicontinuously in fibers.  Thus $\Kcal$ vanishes on a neighborhood of $t$ by Nakayama's lemma and standard support arguments.  Hence \textup{(c)} holds.

Now \textup{(c)}\,$\Rightarrow$\,\textup{(a)} follows because injectivity of a map from a free module is preserved by base change whenever the quotient is $S$-flat: indeed, if
\[
0\to \Lcal\xrightarrow{\alpha}\N\to \Pcal\to 0
\]
is exact and $\Pcal$ is $S$-flat, tensoring with any $\OO_S$-algebra remains exact on the left.  Thus $\alpha$ is universally injective.
\end{proof}

\subsection{Proof of the flatness criterion}

\begin{proof}[Proof of \cref{thm:flat-criterion}]
We proceed by induction on the length $r$ of the d\'evissage.  For $r=1$ the d\'evissage is a single exact sequence
\[
0\to \Lcal_1\xrightarrow{\alpha_1}\N_1\to \Pcal_1\to 0
\]
with $\Lcal_1$ free and $\Pcal_1$ supported in relative dimension $0$.

Assume first that $\M_x$ is flat over $\OO_{S,s}$.  Since $\N_1$ is obtained from $\M$ by finite pushforward along a finite morphism after \etale localization, $\N_1$ is flat over $S$ at the relevant point.  Therefore the exact sequence above shows that $\Pcal_1$ is flat and $\alpha_1$ is universally injective by \cref{lem:ses-flat}.  This proves \textup{(i)}\,$\Rightarrow$\,\textup{(ii)} when $r=1$.

Conversely, if $\alpha_1$ is universally injective and $\Pcal_1$ is flat, then $\N_1$ is flat by \cref{lem:ses-flat}.  Since $\N_1$ represents $\M$ after the local finite factorization, $\M_x$ is flat.  Thus the theorem holds for $r=1$.

Assume now $r>1$.  The first step of the d\'evissage gives
\[
0\to \Lcal_1\xrightarrow{\alpha_1}\N_1\to \Pcal_1\to 0.
\]
If $\M_x$ is flat, then the same argument as above yields that $\alpha_1$ is universally injective and $\Pcal_1$ is flat.  But $\Pcal_1$ carries a d\'evissage of length $r-1$, so by induction all the subsequent maps are universally injective and the terminal quotient is flat.  Hence \textup{(i)}\,$\Rightarrow$\,\textup{(ii)}.

For the converse, assume \textup{(ii)}.  By induction applied to the d\'evissage of $\Pcal_1$, the quotient $\Pcal_1$ is flat over $S$.  Since $\alpha_1$ is universally injective and $\Lcal_1$ is free, the exact sequence above implies that $\N_1$ is flat.  Therefore $\M_x$ is flat.  This proves \textup{(ii)}\,$\Rightarrow$\,\textup{(i)}.

The equivalence with \textup{(iii)} follows from the local criterion of flatness and \cref{lem:ui-smooth}: injectivity after tensoring with $k(s)$ promotes to universal injectivity because the relevant quotients are already flat by the inductive argument.  Condition \textup{(iv)} is merely the geometric reformulation of the same phenomenon after shrinking the pointed \etale neighborhoods.
\end{proof}

\begin{corollary}[Flatness by stratified free presentation]\label{cor:flat-open-near-point}
Let $f\colon X\to S$ be locally of finite presentation and $\M$ a finite type $\OO_X$-module.  If $\M$ is flat at a point $x\in X$, then after shrinking around $x$ there exists a finite filtration
\[
0=M_0\subset M_1\subset\cdots\subset M_r=M
\]
whose successive quotients are cokernels of universally injective maps between free modules over smooth local models.
\end{corollary}

\begin{proof}
Apply \cref{prop:exist-dev} and then \cref{thm:flat-criterion}.
\end{proof}

\begin{remark}
This is the point at which the original paper goes decisively beyond EGA.  EGA gives criterion and openness statements for flatness.  Raynaud--Gruson produce a finite machine that keeps track of the defect of flatness in a way suitable for later birational correction.
\end{remark}

\section{Purity, the relative assassin, and projectivity}

\subsection{Projective modules of countable type and the assassin}

The opening lemmas of \S3 collect several facts on projective modules, some descending to Daniel Lazard's theorem on flat modules as filtered colimits of finite free modules.  We isolate only the statements that enter the geometric argument.

\begin{lemma}\label{lem:proj-loc}
Let $A$ be a ring.
\begin{enumerate}[label=\textup{(\roman*)}]
  \item A countably generated flat $A$-module is a filtered colimit of finite free modules.
  \item If $P$ is projective over $A$ and $u\colon P\to M$ is an $A$-linear map whose base change to every residue field is injective, then $u$ is universally injective.
  \item If $A\to B$ is faithfully flat and $M\ot_A B$ is projective over $A$ through the structure map, then $M$ is projective over $A$ provided the descent datum is effective.
\end{enumerate}
\end{lemma}

\begin{proof}
Part \textup{(i)} is Lazard's theorem.  For part \textup{(ii)}, write $P$ as a direct summand of a free module $F$.  The map $u$ extends to a map $F\to M\oplus Q$ for a complementary summand $Q$.  Fiberwise injectivity implies universal injectivity on the free module by the same argument as in \cref{lem:ui-smooth}; restricting back to the summand gives the claim.  Part \textup{(iii)} is the standard fpqc descent of projective modules.
\end{proof}

\begin{definition}[Relative assassin]
Let $f\colon X\to S$ be locally of finite type and let $\M$ be a finite type $\OO_X$-module.  The \emph{relative assassin} of $\M$ over $S$ is the subset
\[
\Ass_{X/S}(\M):=\bigcup_{s\in S}\Ass(\M_s)\subset X.
\]
When $\M=\OO_X$ we write $\Ass(X/S)$.
\end{definition}

The relative assassin records the associated points of the fibers.  It is therefore the correct geometric object for understanding where vertical embedded components may appear.

\subsection{Purity}

\begin{definition}[Purity along a fiber]\label{def:purity}
Let $f\colon X\to S$ be locally of finite type and $\M$ a finite type $\OO_X$-module.  Fix $s\in S$.  We say that $\M$ is \emph{$S$-pure along the fiber $X_s$} if for every henselian local base change $(S',s')\to(S,s)$ and every associated point
\[
\xi\in \Ass\bigl(\M\times_S S'\bigr)
\]
the closure of $\xi$ meets the special fiber $X_{s'}$.
\end{definition}

Equivalently, associated points are not allowed to drift away from the distinguished fiber.  This is a geometric control condition on the support of $\M$.

\begin{definition}[Relative purity]
We say that $\M$ is \emph{$S$-pure} if it is $S$-pure along every fiber.
\end{definition}

\begin{remark}
Flatness alone does not exclude the appearance of associated points lying above generizations but not specializing to the closed fiber.  Purity is precisely the condition that forbids this phenomenon.  In the paper this is what converts flatness into projectivity.
\end{remark}

\subsection{Flat plus pure implies projective}

We now come to the central theorem of \S3.

\begin{theorem}[Flatness plus purity gives projectivity]\label{thm:flat-pure-projective}
Let $A$ be a ring, let $B$ be an $A$-algebra of finite presentation, and let $M$ be a $B$-module of finite presentation.  Then the following are equivalent:
\begin{enumerate}[label=\textup{(\roman*)}]
  \item $M$ is projective as an $A$-module;
  \item $M$ is flat over $A$ and $A$-pure.
\end{enumerate}
\end{theorem}

\begin{proof}
We first prove \textup{(i)}\,$\Rightarrow$\,\textup{(ii)}.  If $M$ is projective over $A$, then it is certainly flat.  We must prove purity.  Let $(A,\mathfrak m)$ be a henselian local base and let $\xi\in\Ass(M)$ be an associated point in the total space.  Since $M$ is projective, any localization map determined by an element of $A$ acts injectively on $M$ whenever it acts injectively on the special fiber.  In particular, if the closure of $\xi$ failed to meet the special fiber, some element of $A\setminus\mathfrak m$ would annihilate a nonzero submodule supported at $\xi$, contradicting projectivity.  Hence the closure of each associated point meets the special fiber.  Thus $M$ is pure.

For the converse assume that $M$ is flat over $A$ and pure.  The proof is geometric and proceeds in several steps.

\smallskip
\noindent\emph{Step 1: reduction to the affine local-henselian case.}
Projectivity is local on $\Spec A$, so we may assume $A$ local henselian.  Because $B$ is of finite presentation, $X:=\Spec B$ is locally of finite presentation over $S:=\Spec A$.

\smallskip
\noindent\emph{Step 2: find an \etale neighborhood meeting the relative assassin.}
By purity and the geometric criterion of Proposition~\ref{prop:etale-assassin} below, there exists an affine \etale morphism
\[
Y\to X
\]
whose image contains $\Ass(X/S;M)$ and such that the pullback $M_Y$ becomes projective over $A$.  The reason is that over such a neighborhood the relative assassin is controlled by a section of the smooth local model, and then the projective-module lemmas of \S3.1 apply.

\smallskip
\noindent\emph{Step 3: descend projectivity from the \etale neighborhood.}
Since the image of $Y\to X$ contains the relative assassin, the natural map
\[
M\longrightarrow M\ot_B \Gamma(Y,\OO_Y)
\]
is universally injective.  This is one of the crucial consequences of the theory of the relative assassin developed in \S3.2.  The module on the right is projective over $A$ by construction.  Because projectivity descends through faithfully flat morphisms and because universal injectivity excludes hidden torsion in the complement of the image, it follows that $M$ itself is projective over $A$.

More concretely, write $B' = \Gamma(Y,\OO_Y)$.  By \etale descent there exists a finite free $A$-module $F$ and a surjection $F\twoheadrightarrow M\ot_B B'$.  The kernel is also $A$-projective because the quotient is $A$-flat.  Since the comparison map $M\to M\ot_B B'$ is universally injective and since purity guarantees that all associated points are seen on $Y$, the map descends to a splitting over $A$.  Therefore $M$ is a direct summand of a free $A$-module, hence projective.
\end{proof}

The proof above uses the following fundamental geometric criterion, which is the core of the purity theory in the paper.

\begin{proposition}[\etale criterion for purity]\label{prop:etale-assassin}
Let $f\colon X\to S$ be locally of finite presentation and let $\M$ be a finite type $\OO_X$-module.  Then $\M$ is $S$-pure if and only if, after replacing $S$ by a suitable local-henselian neighborhood, there exists an affine \etale morphism
\[
Y\to X
\]
such that the image of $Y$ contains $\Ass_{X/S}(\M)$.
\end{proposition}

\begin{proof}
Assume first that such a $Y\to X$ exists.  Let $\xi\in\Ass_{X/S}(\M)$ and choose $\eta\in Y$ mapping to $\xi$.  Because $Y\to X$ is \etale, the behavior of associated points is preserved.  Since $Y$ is affine over a local-henselian base and the corresponding pulled-back module is projective over the base by construction in the paper, the closure of $\eta$ meets the special fiber.  Therefore the closure of $\xi$ meets the special fiber as well.  Hence $\M$ is pure.

Conversely, assume $\M$ is pure.  One constructs $Y$ by taking the smooth local model from the d\'evissage and then isolating the points of the relative assassin by elementary \etale neighborhoods.  The purity assumption ensures that every associated point specializes to the distinguished fiber; thus one can choose a neighborhood meeting all of them simultaneously.  The image of the resulting affine \etale map contains the entire relative assassin.
\end{proof}

\begin{corollary}[Generic projectivity]\label{cor:generic-proj}
Let $f\colon X\to S$ be locally of finite presentation and let $\M$ be a finite presentation $\OO_X$-module flat over $S$.  Then the locus of points $x\in X$ at which $\M$ is $S$-pure is open, and on that locus $\M$ is projective over the base.
\end{corollary}

\begin{proof}
Openness of the purity locus follows from the openness of the relative assassin and the stability of the condition in \cref{prop:etale-assassin} under shrinking.  Projectivity then follows from \cref{thm:flat-pure-projective}.
\end{proof}

\section{The universal flattening condition}

\subsection{The local-henselian representability theorem}

We now enter \S4 of the paper.  The problem is this.  Suppose
\[
u\colon \M\twoheadrightarrow \N
\]
is a quotient of finite type modules on $X\to S$.  For which base changes $T\to S$ does the pulled back quotient become an isomorphism, or more generally become flat?  The main theorem asserts that this condition is representable by a closed subscheme of the base in the local-henselian setting.

\begin{theorem}[Universal flattening in the local-henselian case]\label{thm:universal-flattening-local}
Let $S$ be a scheme, let $X\to S$ be locally of finite presentation, and let
\[
u\colon \M\twoheadrightarrow \N
\]
be a surjective morphism of finite type $\OO_X$-modules.  Assume that $\M$ is $S$-flat and $S$-pure.  Then the contravariant functor
\[
F\colon (\Sch/S)^\op\to \Ens,
\qquad
T\longmapsto \{\ast\ \text{if } \nu_T\text{ is an isomorphism, }\varnothing\text{ otherwise}\}
\]
is representable by a closed subscheme of $S$.  If $\N$ is of finite presentation, the representing closed subscheme is of finite presentation over $S$.
\end{theorem}

\begin{proof}
The proof is one of the cleanest in the paper once the earlier machinery is available.

\smallskip
\noindent\emph{Step 1: reduction to an affine free model.}
Since the question is local on $S$, we may assume $S=\Spec A$ affine.  Because $\M$ is $S$-flat and $S$-pure, \cref{prop:etale-assassin} provides an affine \etale neighborhood
\[
Y\to X
\]
whose image contains the relative assassin of $\M$ and such that the pullback of $\M$ to $Y$ corresponds to a projective, hence after shrinking free, $A$-module.  Replacing $X$ by $Y$ and using universal injectivity over the assassin, we reduce to the case where
\[
X=\Spec B,\qquad M:=\Gamma(X,\M)
\]
is a free $A$-module.

\smallskip
\noindent\emph{Step 2: translate the condition into an ideal.}
Let $R:=\Ker(\nu)\subset M$.  Since $M$ is free over $A$, choose a basis $e_1,\dots,e_r$ of $M$.  Every element $\rho\in R$ may be written uniquely as
\[
\rho = \sum_{i=1}^r a_i(\rho)e_i,
\qquad a_i(\rho)\in A.
\]
Let $I\subset A$ be the ideal generated by all coefficients $a_i(\rho)$ as $\rho$ runs through $R$ and $1\le i\le r$.

We claim that for an $A$-scheme $T$ with structure sheaf $\OO_T$ the base change $\nu_T$ is an isomorphism if and only if $I\OO_T=0$.  Indeed, since $\nu$ is already surjective, $\nu_T$ is an isomorphism precisely when its kernel vanishes, that is, when $R\ot_A\OO_T=0$.  But because $M$ is free with chosen basis, this happens exactly when all coordinates of all elements of $R$ vanish in $\OO_T$, i.e. when $I\OO_T=0$.

\smallskip
\noindent\emph{Step 3: representability.}
The condition $I\OO_T=0$ is represented by the closed subscheme
\[
S':=\Spec(A/I)\hookrightarrow S.
\]
Hence $F$ is representable by $S'$.  If $\N$ is of finite presentation, then $R$ is finitely generated, so $I$ is finitely generated, and $S'\to S$ is of finite presentation.
\end{proof}

\begin{remark}
The essential point of the proof is that purity reduces the geometry to the affine free case, while flatness ensures that the kernel behaves well under base change.  Once there, the representability statement is nothing but linear algebra.
\end{remark}

\subsection{From isomorphism to flatness}

The paper then improves the previous theorem from the case of an isomorphism condition on a quotient to the actual flattening functor for a finite type sheaf.

\begin{theorem}[Flattening stratifier, local-henselian form]\label{thm:flattening-stratifier-local}
Let $f\colon X\to S$ be locally of finite presentation and $\M$ a finite type $\OO_X$-module.  Fix a point $x\in X$ over $s\in S$.  Then, after replacing $(S,s)$ by a pointed local-henselian neighborhood, there exists a closed subscheme
\[
S_\flat\hookrightarrow S
\]
with the following universal property: for every morphism $T\to S$, the pullback of $\M$ to $X_T$ is flat at all points lying over $x$ if and only if $T\to S$ factors through $S_\flat$.
\end{theorem}

\begin{proof}
Choose a total d\'evissage of $\M$ at $x$.  By \cref{thm:flat-criterion}, flatness is equivalent to universal injectivity of the maps $\alpha_i$ and flatness of the terminal quotient.  Apply \cref{thm:universal-flattening-local} successively to the quotients that measure the cokernels of the $\alpha_i$.  Since the intersection of finitely many closed conditions is closed, we obtain a closed subscheme representing the simultaneous fulfillment of all flatness conditions.
\end{proof}

\subsection{Global form}

The global analogue, also stated in the paper, asserts that if $\M$ is of finite presentation then the subfunctor of the final functor on $S$-schemes consisting of those base changes over which $\M$ becomes flat in dimensions $\ge n$ is representable by a quasi-compact separated monomorphism, in fact by a closed immersion after restricting to the relevant open.  This global form is what feeds into the theorem of flattening by blow-up.

\section{Flattening by blow-up}

\subsection{Strict transform of a module}

Before proving the flattening theorem we recall the strict transform construction.  Let $g\colon S'\to S$ be a blow-up with center disjoint from a dense open $U\subset S$, and let $X':=X\times_S S'$.  For an $\OO_X$-module $\M$ of finite type, its pullback to $X'$ may acquire torsion supported over the exceptional locus.  The \emph{strict transform} of $\M$ is the quotient
\[
\st_g(\M):=g_X^*\M/\Gamma_{X'\setminus g_X^{-1}(f^{-1}(U))}(g_X^*\M),
\]
that is, one removes precisely the subsheaf supported over the bad locus.  Geometrically, this is the analogue for modules of the strict transform of a subscheme: it retains the part of the pullback seen from the good open and discards the vertical excess created by the blow-up.

\begin{definition}[Flat in dimension $\ge n$]
Let $f\colon X\to S$ be of finite presentation and $\M$ a finite type $\OO_X$-module.  We say that $\M$ is \emph{$S$-flat in dimension $\ge n$} if there exists a retrocompact open subset $V\subset X$ such that:
\begin{enumerate}[label=\textup{(\roman*)}]
  \item $\M|_V$ is flat over $S$;
  \item every point of $X\setminus V$ has fiber dimension $<n$ over the base.
\end{enumerate}
\end{definition}

\subsection{Statement of the theorem}

\begin{theorem}[Flattening by admissible blow-up]\label{thm:flattening-blowup}
Let $S$ be quasi-compact and quasi-separated, let $U\subset S$ be a quasi-compact open subset, let $f\colon X\to S$ be of finite presentation, and let $\M$ be a finite type $\OO_X$-module.  Assume that $\M|_{f^{-1}(U)}$ is $U$-flat in dimension $\ge n$.  Then there exists a $U$-admissible blow-up
\[
g\colon S'\to S
\]
such that the strict transform $\st_g(\M)$ is $S'$-flat in dimension $\ge n$.
\end{theorem}

This is the celebrated theorem of platification by blow-up.  It says that flatness is not only a local property one can test; it is also a birational property one can force after a suitable modification of the base.

\subsection{The projective prototype}

In the projective noetherian case the theorem is guided by the Quot scheme.  Let us explain the idea because it illuminates the general proof.

Suppose $X\to S$ is projective and noetherian and let $\M$ be coherent.  Over the open subset where $\M$ is already flat, the family defines a point of the appropriate Quot functor.  Properness of the Quot scheme then allows one to take the schematic closure of this point.  The resulting projective morphism
\[
S'\to S
\]
is an isomorphism over the good open, and the universal quotient on $X\times_S S'$ is flat over $S'$.  The strict transform formalism is a conceptual extraction of exactly this quotient after one removes the part supported over the exceptional locus.

The general theorem is harder because projectivity and noetherian hypotheses must be removed, but the underlying idea remains the same: one seeks a birational modification over which the pullback of $\M$ acquires a universal flat quotient that agrees with $\M$ on the good open.

\subsection{Proof mechanism}

We now explain the proof of \cref{thm:flattening-blowup} in a detailed seminar style.

\begin{proof}[Proof of \cref{thm:flattening-blowup}]
The proof is by induction on $n$ together with noetherian approximation and repeated use of the local flattening stratifier.  We separate the argument into six steps.

\smallskip
\noindent\emph{Step 1: reduction to the noetherian affine situation.}
By standard limit arguments for quasi-compact quasi-separated schemes and finitely presented morphisms, one may approximate $S$, $X$, and $\M$ by noetherian models.  Since admissible blow-ups and strict transforms are compatible with these approximations, it suffices to prove the theorem in the noetherian case.

\smallskip
\noindent\emph{Step 2: induction on the threshold dimension $n$.}
If $n<0$ there is nothing to prove.  Assume the theorem known for smaller values of $n$.  The failure of flatness in dimension $\ge n$ is concentrated on a closed subset of relative dimension exactly $n$.  The induction will show how to kill this defect while allowing lower-dimensional defects to persist temporarily.

\smallskip
\noindent\emph{Step 3: local flattening centers.}
Choose finitely many points of the bad locus of relative dimension $n$ and around each such point choose a d\'evissage.  By \cref{thm:flattening-stratifier-local}, the condition that the relevant quotient become flat is represented by a closed subscheme of the local base.  One glues these local closed conditions into a global center disjoint from $U$.  Blowing up this center produces a $U$-admissible morphism
\[
g\colon S'\to S
\]
that improves the universal injectivity conditions appearing in the d\'evissage near all points of relative dimension $n$.

\smallskip
\noindent\emph{Step 4: strict transform and removal of exceptional torsion.}
Let $\M':=g_X^*\M$ and let $\M^{\mathrm{str}}:=\st_g(\M)$.  By definition there is an exact sequence
\[
0\to \Tcal\to \M'\to \M^{\mathrm{str}}\to 0
\]
where $\Tcal$ is supported over the exceptional locus lying above the complement of $U$.  Thus on the inverse image of $U$, the strict transform agrees with the ordinary pullback.  The purpose of passing from $\M'$ to $\M^{\mathrm{str}}$ is precisely to eliminate the vertical torsion responsible for failure of flatness.

\smallskip
\noindent\emph{Step 5: control of associated points.}
The difficult part is to show that after the blow-up no associated point of $\M^{\mathrm{str}}$ lying over a point of relative dimension $\ge n$ can remain in the bad locus.  This uses the local theorems of \S3 and \S4 applied to the d\'evissages chosen in Step~3.  Roughly speaking, the center of the blow-up is chosen so that the universal flattening conditions become satisfied after pullback, and the only obstruction that remains sits in the torsion sheaf $\Tcal$, which is then removed by strict transform.

\smallskip
\noindent\emph{Step 6: conclude by the flatness criterion.}
At a point $x'\in X'$ of fiber dimension at least $n$, the transformed d\'evissage now has universally injective transition maps and terminal quotient flat over the new base.  Therefore \cref{thm:flat-criterion} applies and shows that $\M^{\mathrm{str}}$ is flat at $x'$.  Since this holds at all points of fiber dimension at least $n$, the strict transform is $S'$-flat in dimension $\ge n$.
\end{proof}

\subsection{The endgame and why strict transform is indispensable}

The final stage of the argument deserves to be isolated, because it contains the most conceptual idea in the theorem.

Suppose one has already performed a blow-up $S'\to S$ for which the flattening conditions hold on a dense open subset of every fiber of relative dimension $n$.  The pullback $g_X^*\M$ may nevertheless fail to be flat because of a subsheaf supported on the exceptional divisor.  This subsheaf is invisible on the good open but becomes fatal for flatness if retained.  The strict transform removes exactly this subsheaf.

The reason this works is the following.  Let
\[
0\to \Tcal\to g_X^*\M\to \M^{\mathrm{str}}\to 0
\]
be the defining exact sequence.  The support of $\Tcal$ is contained in the exceptional locus, hence in a set of fiber dimension $<n$ after the blow-up has already improved the codimension-$n$ defect.  Therefore for points of fiber dimension $\ge n$ the stalks of $\Tcal$ vanish.  At such points the strict transform agrees with the pullback, and the d\'evissage criteria show flatness.  Thus the strict transform is the minimal quotient of the pullback that preserves the good behavior on $U$ while removing the unavoidable exceptional torsion.

\section{Further consequences and geometric meaning}

\subsection{Why purity is the bridge to projectivity}

The theorem \cref{thm:flat-pure-projective} may be read as saying that projectivity over the base is not a mysterious lifting property but the conjunction of two geometric conditions:
\[
\text{projective over }A
\quad\Longleftrightarrow\quad
\text{flat over }A\ \text{and pure over }A.
\]
Flatness controls infinitesimal variation in families.  Purity controls the location of associated points and excludes hidden vertical components.  Together they force the module to behave as though it were pulled back from a vector bundle on the base, and this is exactly what projectivity means in the affine situation.

\subsection{Why d\'evissage is the bridge to flattening}

A purely homological flatness criterion would not be suitable for blow-up arguments because blow-ups are geometric modifications.  The d\'evissage breaks the homological obstruction into finitely many geometric pieces, each encoded by a map from a free module and a smaller-dimensional quotient.  Those maps admit universal flattening conditions representable by closed subschemes.  Therefore one can blow up those closed subschemes and correct the failure of flatness one dimension at a time.

\subsection{Connection with later flattening stratifications}

Modern treatments of flattening stratifications often present the result as a representability statement for the functor of base changes over which a coherent sheaf becomes flat.  Raynaud--Gruson's contribution is stronger in two respects.  First, the result is linked to projectivity and purity rather than treated in isolation.  Second, the theorem of flattening by blow-up produces not merely a stratification but a birational modification that resolves the defect.

\section{Appendix: a worked affine picture}

To close, let us spell out the basic affine mechanism in a toy model.  Let $A$ be a noetherian domain, $B=A[x]/(f)$, and let $M$ be a finite $B$-module.  Suppose that after localizing away from one element $a\in A$ the module $M_a$ is flat over $A_a$.  The flattening theorem predicts that after blowing up the ideal generated by $a$ and finitely many coefficients encoding the kernel of a local free presentation, the strict transform becomes flat over the blow-up.

In this elementary setting one chooses a presentation
\[
A^m\xrightarrow{\Phi} A^r\to M\to 0.
\]
The failure of flatness is controlled by the minors of $\Phi$ and by the torsion of its cokernel.  Blowing up the corresponding Fitting ideal separates the locus where the rank jumps.  Pulling back to the blow-up may introduce exceptional torsion; quotienting by that torsion yields the strict transform.  The theorem says that, after iterating this procedure in the correct universal form, the resulting module is flat.

This affine picture is crude compared with the full theory, but it is faithful to the philosophy of the paper: isolate the defect, blow up the defect, remove the exceptional torsion, and then re-test flatness through the d\'evissage criterion.


\section{Supplementary expanded proofs and proof diagrams}

This section enlarges several places where the main text above deliberately compressed the original arguments.  The aim is not to repeat the whole paper line by line, but to make the logical transitions completely explicit in the principal chain of proofs.

\subsection{Expanded proof of the existence of relative d\'evissage}

We revisit \cref{prop:exist-dev} and spell out the precise mechanism behind the dimension drop.  Let $x\in X$ lie over $s\in S$ and let
\[
n:=\dim_x\bigl(\Supp(\M)_s\bigr).
\]
After the local finite factorization of \cref{thm:zariski-local}, we may assume we are given a finite morphism
\[
\pi\colon X'\to T
\]
with $T\to S$ smooth of relative dimension $n$, where $x'\in X'$ maps to $t\in T$ and is the unique point above $x$.  Put
\[
\N_1:=\pi_*(\M|_{X'}).
\]
Because $\pi$ is finite and $\M$ is finite type, $\N_1$ is finite over $\OO_T$ near $t$.  Choose elements of $\N_1$ whose images form a basis of the fiber $\N_1\otimes k(t)$; equivalently, choose a surjection
\[
\Lcal_1:=\OO_T^{\oplus r}\twoheadrightarrow \N_1
\]
that is surjective at $t$.  Let $\Pcal_1$ denote the cokernel.  By the dimension-drop argument established above, the support of $\Pcal_1$ cannot contain any irreducible component of $\Supp(\N_1)$ through $t$, hence
\[
\dim_t\bigl(\Supp(\Pcal_1)_s\bigr)<n.
\]
This is the exact place where Nakayama's lemma enters the construction.  Repeating the same step for $\Pcal_1$ yields the chain of pointed schemes and exact sequences required by the definition.  The process stops because the integer-valued dimensions drop strictly at each stage.

The original paper packages this construction in a flexible way so that the intermediate schemes can be chosen affine and the columns can be taken to be elementary \etale neighborhoods.  This flexibility is indispensable later when the d\'evissage is transported through other pointed neighborhoods and when the flattening conditions are glued.

\subsection{Proof diagram for the flatness criterion}

It is useful to display the inductive mechanism of \cref{thm:flat-criterion} in a single diagram.  At each stage one has an exact sequence
\[
0\to \Lcal_i\xrightarrow{\alpha_i}\N_i\to \Pcal_i\to 0,
\]
with $\Lcal_i$ free.  Assume inductively that $\Pcal_i$ is flat.  Then by the lemma identifying universal injectivity in the presence of a flat cokernel, the map $\alpha_i$ is universally injective as soon as the sequence is exact.  Conversely, if $\alpha_i$ is universally injective and if both $\Lcal_i$ and $\Pcal_i$ are flat, then $\N_i$ is flat by the standard short exact sequence criterion for flatness.  Since $\N_{i+1}$ is built from $\Pcal_i$, the whole induction can be pictured as
\[
\begin{tikzcd}[column sep=large]
\text{flatness of }\N_i \ar[r,Rightarrow] & \text{universal injectivity of }\alpha_i \text{ and flatness of }\Pcal_i \\
\text{flatness of }\Pcal_i \ar[u,Rightarrow] & \text{flatness of the next quotient }\Pcal_{i+1} \ar[l,Rightarrow] \ar[u,Rightarrow]
\end{tikzcd}
\]
terminating at the last quotient $\Pcal_r$.  The theorem is therefore a finite induction down the d\'evissage tower.

\subsection{Expanded proof of the implication flat $+$ pure $\Rightarrow$ projective}

We now sharpen the proof of \cref{thm:flat-pure-projective}.  The delicate point is not the necessity of purity for projective modules, but the sufficiency.

Assume $A$ is local henselian, $B$ is of finite presentation over $A$, and $M$ is a finite $B$-module, flat and $A$-pure.  By \cref{prop:etale-assassin}, there exists an affine \etale map $Y\to X:=\Spec B$ whose image contains the relative assassin of $M$.  Write $B':=\Gamma(Y,\OO_Y)$ and $M':=M\otimes_B B'$.  The reason for choosing $Y$ this way is that the support-theoretic obstructions to projectivity of $M$ are all visible on $Y$: every associated point of every fiber is already met there.

By the preliminary lemmas on projective modules, together with the smooth local models constructed earlier, one may choose $Y$ so that $M'$ becomes projective over $A$.  The key descent input from the paper is that the canonical map
\[
M\longrightarrow M'
\]
is universally injective.  This follows from the fact that the image of $Y\to X$ contains the relative assassin: if an element of $M$ vanished in $M'$, its annihilator would define an associated point outside the image, impossible by construction.  Universal injectivity then implies that $M$ sits inside an $A$-projective module without creating hidden $A$-torsion.

Now choose a surjection $F\twoheadrightarrow M'$ with $F$ finite free over $A$, and let $K'$ be its kernel.  Since $M'$ is $A$-flat, the module $K'$ is also $A$-flat.  The descent datum on $Y\times_X Y$ preserves this exact sequence; because the comparison map $M\to M'$ is universally injective and the obstruction locus has empty intersection with the relative assassin, the descended sequence over $A$ remains exact.  Thus $M$ is a direct summand of a finite free $A$-module, hence projective.  The moral is that purity exactly excludes the possibility that descent fails because of vertical associated components invisible on the chosen \etale chart.

\subsection{Expanded proof of the representability theorem 4.1.1}

The proof of \cref{thm:universal-flattening-local} becomes almost tautological after the reduction to the free case, so it is worth stating the reduction as a separate lemma.

\begin{lemma}\label{lem:reduce-free}
Under the hypotheses of \cref{thm:universal-flattening-local}, the representability question may be checked after replacing $X$ by an affine \etale neighborhood whose image contains $\Ass_{X/S}(\M)$ and on which $\M$ is free over the base.
\end{lemma}

\begin{proof}
Choose $Y\to X$ as in \cref{prop:etale-assassin}.  Because the image contains the relative assassin, the kernel of $\M\to j_*j^*\M$ has no associated points and therefore vanishes.  Hence the question whether a quotient becomes an isomorphism is unaffected by replacing $X$ by $Y$.  On $Y$, projectivity over the base allows one, after shrinking if necessary, to replace $\M$ by a free module.
\end{proof}

With this lemma in hand, the proof of \cref{thm:universal-flattening-local} is reduced to the following elementary fact: if $M$ is free over $A$ and $R\subset M$ is a submodule, then the functor of $A$-schemes $T$ such that $R\otimes_A \OO_T=0$ is represented by the closed subscheme cut out by the coefficients of generators of $R$ in a basis of $M$.  This is exactly the argument already given in the main body, but phrased so that the purity reduction becomes formally separate from the coefficient-ideal computation.

\subsection{Expanded noetherian projective prototype for flattening by blow-up}

To make the general theorem more concrete, let us state the projective prototype in a clean form.

\begin{proposition}[Projective noetherian prototype]\label{prop:quot-prototype}
Assume $S$ is noetherian, $X\to S$ projective, and $\M$ coherent on $X$.  Let $U\subset S$ be the maximal open on which $\M$ is flat.  Then there exists a projective birational morphism $g\colon S'\to S$, an isomorphism over $U$, such that the strict transform of $\M$ on $X\times_S S'$ is flat over $S'$.
\end{proposition}

\begin{proof}
Over $U$, the sheaf $\M|_{X_U}$ is flat, hence defines a $U$-point of the Quot functor parameterizing quotients of a sufficiently large projective presentation on $X/S$.  By Grothendieck's representability theorem, this Quot functor is represented by a projective $S$-scheme.  Taking the schematic closure of the $U$-point inside the Quot scheme yields a projective morphism $S'\to S$ that is an isomorphism over $U$.  Pulling back the universal quotient to $X\times_S S'$ gives a coherent quotient of the pullback of $\M$ which agrees with $\M$ over $U$ and is flat over $S'$.  By uniqueness of the maximal quotient agreeing with the pullback over the inverse image of $U$, this quotient is precisely the strict transform.
\end{proof}

This proposition contains the essential geometry of the general theorem: one does not modify $X$ but rather the base $S$, and the universal quotient on the modification is nothing other than the strict transform.

\subsection{A more explicit endgame for Theorem~\texorpdfstring{\ref{thm:flattening-blowup}}{}}

The original proof of the full flattening theorem ends by combining induction in the dimension parameter with the removal of exceptional torsion.  The following reformulation isolates the final move.

\begin{lemma}[Final strict-transform step]\label{lem:strict-endgame}
Let $g\colon S'\to S$ be a $U$-admissible blow-up.  Suppose there exists a retrocompact open $V'\subset X':=X\times_S S'$ such that
\begin{enumerate}[label=\textup{(\roman*)}]
  \item the pullback $g_X^*\M$ is flat over $S'$ on $V'$;
  \item every associated point of the torsion subsheaf $\Tcal\subset g_X^*\M$ supported over $X'\setminus V'$ has fiber-dimension $<n$.
\end{enumerate}
Then the strict transform $\st_g(\M)$ is $S'$-flat in dimension $\ge n$.
\end{lemma}

\begin{proof}
By definition there is an exact sequence
\[
0\to \Tcal\to g_X^*\M\to \st_g(\M)\to 0.
\]
At any point $x'\in X'$ of fiber-dimension at least $n$, the second hypothesis implies $\Tcal_{x'}=0$.  Hence $(\st_g(\M))_{x'}\cong (g_X^*\M)_{x'}$, which is flat by the first hypothesis if $x'\in V'$ and by density and the d\'evissage criterion if $x'$ lies in the closure of $V'$ inside a fiber of dimension at least $n$.  Therefore $\st_g(\M)$ is flat at all points of fiber-dimension at least $n$, which is exactly the claim.
\end{proof}

This lemma makes completely explicit why strict transform is not cosmetic but structurally necessary: the pullback may fail to be flat only because of exceptional torsion, and once that torsion is shown to live in smaller relative dimension it can be discarded without affecting the desired range of dimensions.

\subsection{Glossary of the proof strategy}

For quick reference, the proof strategy of the paper may be summarized in the following seven verbs:
\begin{center}
\fbox{factor $\to$ present $\to$ d\'evise $\to$ test $\to$ purify $\to$ represent $\to$ blow up}
\end{center}
More explicitly:
\begin{enumerate}[label=\textup{(\arabic*)}]
  \item Factor locally by a finite map over a smooth base.
  \item Present the pushed-forward module by a free module.
  \item D\'evise by iterating on smaller-dimensional quotients.
  \item Test flatness through universal injectivity and the terminal quotient.
  \item Purify the geometry by controlling associated points via the relative assassin.
  \item Represent the flattening condition by a closed subspace of the base.
  \item Blow up that closed locus and pass to the strict transform.
\end{enumerate}
This is the conceptual architecture underlying the whole paper.


\section{Further technical amplifications}

This section pushes several technical points beyond the level needed for the main narrative.  The goal is to provide a more monograph-like set of proofs around the homological algebra of universal injectivity, the geometry of the relative assassin, and the exact manner in which strict transform interacts with flattening.

\subsection{Universal injectivity and base change}

We begin with a small package of lemmas that are used repeatedly, sometimes implicitly, in the main body of the paper.

\begin{lemma}[Characterizations of universal injectivity]\label{lem:ui-char}
Let $A$ be a ring and $u\colon M\to N$ a morphism of $A$-modules.  The following are equivalent.
\begin{enumerate}[label=\textup{(\roman*)}]
  \item $u$ is universally injective.
  \item For every ideal $I\subset A$, the induced map
  \[
  M/IM\longrightarrow N/IN
  \]
  is injective.
  \item For every finitely presented $A$-module $Q$, the map
  \[
  M\ot_A Q\longrightarrow N\ot_A Q
  \]
  is injective.
\end{enumerate}
\end{lemma}

\begin{proof}
The implication (iii)$\Rightarrow$(i) is tautological after filtering an arbitrary $A$-algebra by finitely generated submodules.  The implication (i)$\Rightarrow$(ii) is obtained by taking $Q=A/I$.  It remains to show (ii)$\Rightarrow$(iii).  Let $Q$ be finitely presented, and choose an exact sequence
\[
A^p\xrightarrow{\psi} A^q\twoheadrightarrow Q\to 0.
\]
Tensoring with $M$ and $N$ gives a commutative diagram with exact rows
\[
\begin{tikzcd}
M^{\oplus p} \ar[r,"\psi_M"] \ar[d,"u^{\oplus p}"] & M^{\oplus q} \ar[r] \ar[d,"u^{\oplus q}"] & M\ot_A Q \ar[r] \ar[d] & 0 \\
N^{\oplus p} \ar[r,"\psi_N"] & N^{\oplus q} \ar[r] & N\ot_A Q \ar[r] & 0.
\end{tikzcd}
\]
If an element of $M\ot_A Q$ maps to zero in $N\ot_A Q$, choose a lift $m\in M^{\oplus q}$ whose image in $N^{\oplus q}$ lies in the image of $\psi_N$.  Modulo the image of $\psi_M$, this defines an obstruction living in a quotient of a finite direct sum of copies of $M$ by a finitely generated submodule.  Resolving that quotient by successive cyclic modules and invoking (ii) for modules of the form $A/I$, one concludes that the obstruction vanishes.  Hence the map on tensor products is injective.
\end{proof}

\begin{lemma}[Flat cokernels from universal injectivity]\label{lem:coker-flat-ui-2}
Let
\[
0\to L\xrightarrow{u} N\to P\to 0
\]
be exact with $L$ and $N$ flat over $A$.  If $u$ is universally injective, then $P$ is flat over $A$.
\end{lemma}

\begin{proof}
For any $A$-module $Q$, tensoring yields
\[
0\to L\ot_A Q\to N\ot_A Q\to P\ot_A Q\to 0,
\]
where exactness on the left follows from universal injectivity and exactness on the right follows from right exactness of tensor product.  Therefore $\Tor_1^A(P,Q)=0$ for all $Q$, so $P$ is flat.
\end{proof}

\begin{lemma}[Local detection of universal injectivity]\label{lem:local-ui}
Let $(A,\mathfrak m)$ be a local ring and let
\[
0\to L\xrightarrow{u} N\to P\to 0
\]
be exact with $L$ flat and $P$ flat.  If
\[
L\ot_A k(\mathfrak m)\longrightarrow N\ot_A k(\mathfrak m)
\]
is injective, then $u$ is universally injective.
\end{lemma}

\begin{proof}
It is enough to check injectivity modulo every ideal $I\subset A$.  Suppose for contradiction that the map $L/IL\to N/IN$ is not injective for some $I$.  Localizing at a prime minimal in the support of its kernel, we reduce to the case where the kernel is annihilated by the maximal ideal.  Tensoring further with $k(\mathfrak m)$ then produces a nonzero kernel on the closed fiber, contradiction.
\end{proof}

These lemmas explain why the d\'evissage criterion is naturally phrased in terms of universal injectivity rather than split exactness: universal injectivity is exactly the amount of rigidity preserved under arbitrary base change.

\subsection{Dimension drop and Nakayama's lemma in the d\'evissage}

The induction in the construction of a relative d\'evissage rests on the following elementary but decisive point.

\begin{lemma}[Dimension drop under a fiberwise adapted presentation]\label{lem:dimension-drop-expanded}
Let $T\to S$ be locally of finite presentation, let $t\in T$ map to $s\in S$, and let $\N$ be a finite type $\OO_T$-module.  Assume that a map
\[
\OO_T^{\oplus r}\longrightarrow \N
\]
is surjective on the stalk at $t$ after restriction to the fiber over $s$.  If $\Pcal$ denotes the cokernel, then no irreducible component of $\Supp(\N_s)$ passing through $t$ is contained in $\Supp(\Pcal_s)$.  In particular,
\[
\dim_t\Supp(\Pcal_s)<\dim_t\Supp(\N_s)
\]
whenever $t\in\Supp(\N_s)$.
\end{lemma}

\begin{proof}
By Nakayama's lemma, surjectivity on the fiber at $t$ implies surjectivity on the stalk at $t$, so $\Pcal_t=0$.  Therefore $t\notin\Supp(\Pcal_s)$.  Any irreducible component of $\Supp(\N_s)$ through $t$ is forced to disappear from the support of $\Pcal_s$, which yields the strict dimension drop.
\end{proof}

This simple observation is the numerical engine behind the finiteness of the d\'evissage: one can only strictly decrease a nonnegative integer a finite number of times.

\subsection{Purity, closures of associated points, and the relative assassin}

The relation between purity and the relative assassin is one of the most subtle geometric ideas in the paper.  We record the support-theoretic heart of the argument.

\begin{proposition}[Purity as closure control]\label{prop:purity-support}
Let $f\colon X\to S$ be of finite presentation and let $\M$ be a finite type $\OO_X$-module, flat over $S$.  Fix $s\in S$.  The following are equivalent in a neighborhood of $X_s$.
\begin{enumerate}[label=\textup{(\roman*)}]
  \item $\M$ is pure along $X_s$.
  \item For every associated point $\xi$ of $\M$ lying over a generization of $s$, the closure of $\{\xi\}$ meets $X_s$.
  \item Every point of the relative assassin $\Ass_{X/S}(\M)$ specializes to a point of $\Supp(\M_s)$.
\end{enumerate}
\end{proposition}

\begin{proof}
The equivalence of (ii) and (iii) is a reformulation of the definition of the relative assassin on henselian neighborhoods.  The implication (i)$\Rightarrow$(ii) is exactly the definition of purity.  Conversely, if (ii) holds, then after henselization at $s$ every associated point of the pulled-back module still has closure meeting the special fiber, because henselization preserves the specialization relation and does not destroy generic points of the support.  Thus the pulled-back module is pure along the henselian special fiber, hence $\M$ is pure along $X_s$.
\end{proof}

\begin{corollary}[\'Etale detection of purity]\label{cor:purity-etale-detect}
Under the hypotheses of \cref{prop:purity-support}, let $j\colon Y\to X$ be an affine \etale morphism whose image contains $\Ass_{X/S}(\M)$ near $X_s$.  Then $\M$ is pure along $X_s$ if and only if $j^*\M$ is pure along the fiber of $Y\to S$ over $s$.
\end{corollary}

\begin{proof}
Purity is preserved by arbitrary base change, so one direction is immediate.  Conversely, if $j^*\M$ is pure and $\xi\in\Ass_{X/S}(\M)$ lies near $X_s$, choose $\eta\in Y$ over $\xi$.  Then $\eta$ belongs to the relative assassin of $j^*\M$, so by purity of $j^*\M$ its closure meets $Y_s$.  Applying $j$, the closure of $\xi$ meets $X_s$.  The proposition yields purity of $\M$.
\end{proof}

This is the precise reason the paper repeatedly insists on finding \etale neighborhoods whose image contains the relative assassin: those neighborhoods are exactly the right support-theoretic probes.

\subsection{A fuller descent step in the theorem flat plus pure implies projective}

\begin{proposition}[Projectivity descended from a separating \etale chart]\label{prop:projectivity-descent}
Let $A\to B$ be of finite presentation and let $M$ be a finite $B$-module, flat over $A$.  Suppose there exists an affine \etale $B$-algebra $B'$ such that:
\begin{enumerate}[label=\textup{(\roman*)}]
  \item $M':=M\ot_B B'$ is projective over $A$;
  \item the image of $\Spec B'\to \Spec B$ contains the relative assassin of $\widetilde M$ over $\Spec A$.
\end{enumerate}
Then $M$ is projective over $A$.
\end{proposition}

\begin{proof}
By the support-theoretic criterion developed earlier in the notes, condition (ii) implies that the canonical map
\[
M\longrightarrow M'
\]
is universally injective.  Choose a finite free $A$-module $F$ and a surjection $F\twoheadrightarrow M'$.  Since $M'$ is projective, this surjection splits, so $M'$ is a direct summand of $F$.

Form the fiber product $P:=F\times_{M'} M$.  Because $M\to M'$ is injective, the projection $P\to M$ is surjective.  Its kernel identifies with a submodule of the projective $A$-module $\Ker(F\to M')$, hence is flat after passage to the separating chart.  The universal injectivity of $M\to M'$ ensures that this kernel has no hidden associated points away from the image of $\Spec B'$.  Since that image already contains the relative assassin, the kernel is controlled everywhere.  Consequently $P$ is flat over $A$.

After base change to $B'$, the sequence
\[
0\to \Ker(P\to M)\to P\to M\to 0
\]
becomes split exact because $M'$ is projective over $A$.  The obstruction to descending this splitting is supported outside the relative assassin, hence vanishes by the same argument used in the paper to descend universal injectivity from assassin-covering charts.  Therefore $M$ is a direct summand of the flat finite type $A$-module $P$.  Over a ring, a finite type flat direct summand of a finite free module is projective.  Thus $M$ is projective over $A$.
\end{proof}

The point of this proposition is not to replace the geometric argument of Raynaud--Gruson, but to separate out its algebraic descent content: purity provides a chart on which projectivity is visible, and the relative assassin guarantees that the chart loses no essential support.

\subsection{Strict transform, exceptional torsion, and maximality}

The strict transform of a module under an admissible blow-up is the exact quotient needed for flattening.  Let us isolate its formal properties.

\begin{definition}[Strict transform of a module]
Let $g\colon S'\to S$ be a $U$-admissible blow-up, let $f\colon X\to S$, and put $X':=X\times_S S'$.  Let $g_X\colon X'\to X$ be the projection and let $U':=g^{-1}(U)$.  For a quasi-coherent $\OO_X$-module $\M$, its strict transform $\st_g(\M)$ is the quotient of $g_X^*\M$ by the subsheaf $\Tcal$ of sections supported in $X'\setminus f^{-1}(U')$.
\end{definition}

\begin{lemma}[Maximal quotient characterization]\label{lem:strict-maximal}
With the notation above, $\st_g(\M)$ is the maximal quotient of $g_X^*\M$ whose restriction to $f^{-1}(U')$ agrees with the pullback of $\M|_{f^{-1}(U)}$.
\end{lemma}

\begin{proof}
By construction, the kernel $\Tcal$ vanishes on $f^{-1}(U')$, so the quotient agrees with the original pullback over that open subset.  Conversely, if $Q$ is any quotient of $g_X^*\M$ with the same property, then its kernel is supported in $X'\setminus f^{-1}(U')$, so it is contained in $\Tcal$ by maximality of the subsheaf of sections supported there.  Hence the quotient factors uniquely through $\st_g(\M)$.
\end{proof}

\begin{lemma}[Flatness in large dimension from support control]\label{lem:strict-flat-codim}
Assume $g_X^*\M$ is flat over $S'$ on an open subset $V'\subset X'$ whose complement has fiber-dimension $<n$, and assume the torsion subsheaf $\Tcal$ defining the strict transform is supported in $X'\setminus V'$.  Then $\st_g(\M)$ is $S'$-flat in dimension $\ge n$.
\end{lemma}

\begin{proof}
On $V'$ the kernel $\Tcal$ vanishes, so the quotient map
\[
g_X^*\M\twoheadrightarrow \st_g(\M)
\]
is an isomorphism.  Therefore $\st_g(\M)$ is flat on $V'$.  Since the complement of $V'$ has fiber-dimension $<n$, this is exactly the definition of $S'$-flatness in dimension $\ge n$.
\end{proof}

These lemmas show very clearly why strict transform is not a cosmetic modification.  The pullback can fail to be flat solely because exceptional torsion appears above the bad locus; strict transform removes exactly that torsion and nothing more.

\subsection{A cleaner proof skeleton for flattening by blow-up}

We may now compress the logic of the full theorem into a proof skeleton that mirrors the original paper while making the induction easier to track.

\begin{proposition}[Proof skeleton for the flattening theorem]\label{prop:blowup-skeleton}
Let $f\colon X\to S$ be of finite presentation, let $\M$ be a finite type $\OO_X$-module, and let $U\subset S$ be a quasi-compact open such that $\M|_{f^{-1}(U)}$ is flat in dimension $\ge n$.  The proof of flattening by $U$-admissible blow-up proceeds through the following stages.
\begin{enumerate}[label=\textup{(\arabic*)}]
  \item Reduce by noetherian approximation to the case where $S$ is noetherian.
  \item Reduce to a projective situation where Quot represents the relevant family of quotients.
  \item Construct a universal quotient on a $U$-admissible blow-up $S'\to S$ which is flat over $S'$ where needed.
  \item Identify that quotient with the strict transform via \cref{lem:strict-maximal}.
  \item Use relative d\'evissage and the flatness criterion to show that the exceptional torsion is supported in fiber-dimension $<n$.
  \item Conclude by \cref{lem:strict-flat-codim}.
\end{enumerate}
\end{proposition}

\begin{proof}
The global reductions in steps (1)--(3) are the long algebro-geometric part of the original proof.  The decisive formal part begins once one has a flat quotient $Q$ of $g_X^*\M$ that agrees with the original sheaf over the inverse image of $U$.  By the maximal quotient characterization, one gets a canonical factorization
\[
g_X^*\M\twoheadrightarrow \st_g(\M)\twoheadrightarrow Q.
\]
Conversely, because $Q$ itself has the defining property of the strict transform over the good open, there is a factorization in the other direction, hence $Q\cong \st_g(\M)$.  It therefore remains only to show that the kernel of $g_X^*\M\to \st_g(\M)$ is supported in fiber-dimension $<n$.  This is done by descending on the dimension parameter through the d\'evissage tower and applying the flatness criterion from Part~I to the successive quotients.  Once that support estimate is established, the preceding lemma gives the conclusion.
\end{proof}

\subsection{Comparative remarks with classical generic flatness}

\begin{remark}
The EGA theorem of generic flatness says, roughly, that after shrinking the base one reaches flatness.  Raynaud--Gruson make three decisive advances beyond this classical statement.
\begin{enumerate}[label=\textup{(\roman*)}]
  \item They express flatness by a finite inductive homological machine, namely relative d\'evissage plus universal injectivity.
  \item They isolate the support-theoretic obstruction to projectivity through purity and the relative assassin.
  \item They show that bad flatness can be removed birationally by admissible blow-up rather than merely circumvented by deleting a closed subset of the base.
\end{enumerate}
Thus the paper turns generic flatness from an openness theorem into a functorial and birational theory.
\end{remark}

\begin{remark}
From a modern point of view, the paper anticipates a recurring strategy in relative moduli theory: first represent a good quotient condition, then pass to a modification of the base where the universal quotient acquires the desired flatness property, and finally identify that quotient with a canonical strict transform.  The flattening theorem is therefore not just a technical result about modules; it is also a prototype for many later compactification constructions.
\end{remark}


\section{Proof expansions and complements}

This section develops the exposition in a more ample and connected way, closer to a seminar monograph.  The emphasis is on the exact logical junctures where the original paper compresses several reductions into a few lines: the passage from fiberwise injectivity to universal injectivity, the descent step in the theorem ``flat plus pure implies projective'', the compatibility of strict transform with successive admissible blow-ups, and the noetherian approximation step that carries the flattening theorem from the noetherian case to the general qcqs setting.

\subsection{A fully written inductive proof of the flatness criterion}

We now rewrite the proof of the flatness criterion in a more literal induction on the length of a total $S$-d\'evissage.  Let
\[
0\to \Lcal_i\xrightarrow{\alpha_i}\N_i\to \Pcal_i\to 0,
\qquad 1\le i\le r,
\]
be such a d\'evissage at the chosen point, with $\Lcal_i$ free and with $\N_{i+1}$ constructed from $\Pcal_i$ for $i<r$.  Write $A:=\OO_{S,s}$ for the local ring of the base point.

\begin{proposition}[Inductive form of the flatness criterion]\label{prop:inductive-flatness-criterion}
With notation as above, the following are equivalent.
\begin{enumerate}[label=\textup{(\roman*)}]
  \item $\N_1$ is $A$-flat.
  \item For every $i$, the map $\alpha_i$ is universally injective and the final quotient $\Pcal_r$ is $A$-flat.
  \item For every $i$, the map $\alpha_i\otimes_A k(s)$ is injective and $\Pcal_r$ is $A$-flat.
\end{enumerate}
\end{proposition}

\begin{proof}
We prove the equivalence (i)$\Leftrightarrow$(ii)$\Leftrightarrow$(iii).

Assume first that $\N_1$ is flat.  In the exact sequence
\[
0\to \Lcal_1\xrightarrow{\alpha_1}\N_1\to \Pcal_1\to 0,
\]
we know that $\Lcal_1$ is free, hence flat, and $\N_1$ is flat by hypothesis.  Tensoring with an arbitrary $A$-algebra $A'$ preserves left exactness because $\alpha_1$ identifies $\Lcal_1\ot_A A'$ with the kernel of $\N_1\ot_A A'\to \Pcal_1\ot_A A'$.  Therefore $\alpha_1$ is universally injective and $\Pcal_1$ is flat by \cref{lem:coker-flat-ui-2}.  Since the tail of the d\'evissage of length $r-1$ is a d\'evissage of $\Pcal_1$, the induction hypothesis implies that all $\alpha_i$ for $i\ge2$ are universally injective and that $\Pcal_r$ is flat.  Thus (i) implies (ii).

Next, (ii) plainly implies (iii), since universal injectivity may be tested after tensoring with $k(s)$.

It remains to prove that (iii) implies (i).  We argue by induction on $r$.  If $r=1$, then we have one exact sequence
\[
0\to \Lcal_1\xrightarrow{\alpha_1}\N_1\to \Pcal_1\to 0.
\]
The injectivity of $\alpha_1\otimes_A k(s)$, together with flatness of $\Lcal_1$ and $\Pcal_1$, implies by \cref{lem:local-ui} that $\alpha_1$ is universally injective.  Hence $\N_1$ is flat by \cref{lem:ses-flat}.

Assume now $r>1$.  By the induction hypothesis applied to the tail d\'evissage of $\Pcal_1$, the hypotheses in (iii) imply that $\Pcal_1$ is flat.  Since $\Lcal_1$ is free, it is flat, and the injectivity of $\alpha_1\otimes_A k(s)$ again upgrades to universal injectivity by \cref{lem:local-ui}.  The short exact sequence then shows that $\N_1$ is flat.  This proves (iii)$\Rightarrow$(i).
\end{proof}

The importance of this rewritten proof is that it isolates exactly two non-formal inputs: first, the d\'evissage must be arranged so that each successive quotient really has smaller relative dimension; second, the passage from injectivity on the closed fiber to universal injectivity uses the geometric preparation from \S1 and not merely abstract module theory.

\subsection{Projective descent from purity through the relative assassin}

The statement ``flat plus pure implies projective'' is proved in the paper by combining purity with a separation argument using the relative assassin.  We now expand the descent step in full algebraic detail.

\begin{proposition}[Descent package for the theorem flat $+$ pure $\Rightarrow$ projective]\label{prop:projective-descent-package}
Let $(A,\mathfrak m)$ be henselian local, let $B$ be an $A$-algebra of finite presentation, and let $M$ be a finite $B$-module which is flat and pure over $A$.  Assume there exists an affine \etale $B$-algebra $B'$ such that:
\begin{enumerate}[label=\textup{(\alph*)}]
  \item the image of $\Spec B'\to \Spec B$ contains the relative assassin of $M$ over $A$;
  \item $M':=M\otimes_B B'$ is projective as an $A$-module.
\end{enumerate}
Then $M$ is projective as an $A$-module.
\end{proposition}

\begin{proof}
Choose a finite free $A$-module $F$ together with a surjection $F\twoheadrightarrow M'$.  Let $K'$ be the kernel.  Since $M'$ is $A$-projective, the exact sequence
\[
0\to K'\to F\to M'\to 0
\]
splits $A$-linearly; in particular $K'$ is $A$-projective and hence $A$-flat.

Consider the canonical map $u\colon M\to M'$.  Because the image of $\Spec B'\to\Spec B$ contains the relative assassin, the map $u$ is universally injective by the support-theoretic criterion developed in Part~I and sharpened in \cref{prop:etale-assassin,prop:purity-support}.  We may therefore regard $M$ as an $A$-pure submodule of the projective $A$-module $M'$.  Since $A$ is henselian local, to prove that $M$ is projective it is enough to prove that $M$ is a direct summand of some finite free $A$-module.

Let $F\twoheadrightarrow M'$ be as above, and form the pullback module
\[
E:=F\times_{M'} M.
\]
Then there is a commutative diagram with exact rows and columns
\[
\begin{tikzcd}[column sep=large,row sep=large]
0 \ar[r] & K' \ar[r] \ar[d,equal] & E \ar[r] \ar[d] & M \ar[r] \ar[d,"u"] & 0 \\
0 \ar[r] & K' \ar[r] & F \ar[r] & M' \ar[r] & 0.
\end{tikzcd}
\]
Because $K'$ is $A$-flat and $u$ is universally injective, the snake lemma after arbitrary base change shows that $E\to F$ is universally injective.  Consequently, for every $A$-algebra $A'$, the map $E\ot_A A'\to F\ot_A A'$ is injective.  Since $F$ is finite free, the quotient $F/E$ is $A$-flat by \cref{lem:coker-flat-ui-2}.  But $F/E\cong M'/u(M)$.  The support of this quotient misses the relative assassin by construction, and therefore purity forces it to vanish.  Indeed, any nonzero quotient of $M'$ supported away from the relative assassin would create an associated point of a fiber invisible on the chosen \etale chart, contradicting hypothesis (a).  Hence $E\to F$ is surjective as well as injective, so $E\cong F$.

The top row now becomes
\[
0\to K'\to F\to M\to 0.
\]
Because $K'$ is $A$-projective, this sequence splits.  Therefore $M$ is a direct summand of the finite free module $F$, and hence $M$ is finite projective over $A$.
\end{proof}

This argument makes transparent the exact use of purity.  Purity does not by itself build a splitting.  What it does is exclude the existence of a residual quotient supported away from the relevant associated points, thereby forcing the projective object on the \etale chart to descend all the way back to the original module.

\subsection{Strict transform and successive blow-ups}

In the proof of the flattening theorem one repeatedly composes admissible blow-ups.  The strict transform must therefore be compatible with composition.  The original paper uses this repeatedly, often implicitly.  We record it as a separate proposition.

\begin{proposition}[Compatibility of strict transform with composition]\label{prop:strict-composition}
Let $g\colon S'\to S$ and $h\colon S''\to S'$ be admissible blow-ups relative to the same open set $U\subset S$, and let $f\colon X\to S$ be of finite presentation.  Write $X':=X\times_S S'$ and $X'':=X\times_S S''$.  For any finite type $\OO_X$-module $\M$, there is a canonical isomorphism
\[
\st_h(\st_g(\M))\cong \st_{g\circ h}(\M).
\]
\end{proposition}

\begin{proof}
Let $j\colon f^{-1}(U)\hookrightarrow X$ be the open immersion, and similarly $j'\colon (f')^{-1}(U)\hookrightarrow X'$ and $j''\colon (f'')^{-1}(U)\hookrightarrow X''$.  By definition,
\[
\st_g(\M)=\Image\bigl(g_X^*\M\longrightarrow j'_*j'^*(g_X^*\M)\bigr),
\]
and similarly for $h$ and for $g\circ h$.

Over the inverse image of $U$, all three blow-ups are isomorphisms.  Hence the restrictions of $g_X^*\M$, $h_{X'}^*\st_g(\M)$, and $(g\circ h)_X^*\M$ to $X''_{U}$ identify canonically with the same module, namely the pullback of $\M|_{X_U}$.  It follows immediately from the universal property of image sheaves inside the corresponding direct images from the good open that both $\st_h(\st_g(\M))$ and $\st_{g\circ h}(\M)$ are the maximal quotient of $(g\circ h)_X^*\M$ agreeing with this common restriction over $U$.

To make this precise, note first that the pullback of the quotient map $g_X^*\M\twoheadrightarrow \st_g(\M)$ to $X''$ yields a quotient
\[
(g\circ h)_X^*\M\twoheadrightarrow h_{X'}^*\st_g(\M).
\]
The kernel is supported over the complement of $U$, because the quotient is an isomorphism over $U$.  Passing to the strict transform with respect to $h$ kills precisely the largest subsheaf of this pullback supported over the exceptional locus of $h$ above $S'\setminus U$.  Therefore the composite quotient
\[
(g\circ h)_X^*\M\twoheadrightarrow \st_h(\st_g(\M))
\]
agrees with $\M$ over $U$ and factors through $\st_{g\circ h}(\M)$ by the defining maximality of the latter.

Conversely, the canonical map $(g\circ h)_X^*\M\twoheadrightarrow \st_{g\circ h}(\M)$ clearly factors through $h_{X'}^*\st_g(\M)$ because the kernel of $g_X^*\M\to \st_g(\M)$ is already supported over $S'\setminus U$, hence its pullback to $X''$ lies in the kernel of the quotient defining $\st_{g\circ h}(\M)$.  Applying the same maximality argument now in the category of quotients of $h_{X'}^*\st_g(\M)$ which are isomorphic over $U$, we obtain a map
\[
\st_{g\circ h}(\M)\longrightarrow \st_h(\st_g(\M)).
\]
The two constructed morphisms are inverse because they are both the identity over the schematically dense open $U$ and both source and target are quotients of the same pullback by subsheaves supported in the closed complement of $U$.
\end{proof}

This compatibility result clarifies why the inductive blow-up process can be iterated without ambiguity: the strict transform depends only on the resulting birational modification, not on the way the modification is decomposed into elementary steps.

\subsection{Noetherian approximation in the flattening theorem}

One of the most delicate passages in the paper is the extension from the noetherian case to a general quasi-compact quasi-separated base.  The argument belongs to the EGA style of approximation, but the strict-transform formalism forces a particularly careful bookkeeping.  We record the approximation step in a self-contained form.

\begin{proposition}[Approximation step for flattening]\label{prop:approx-flattening}
Let $S$ be qcqs, let $f\colon X\to S$ be of finite presentation, and let $\M$ be a finite type $\OO_X$-module.  Assume there exists a quasi-compact open $U\subset S$ such that $\M|_{f^{-1}(U)}$ is flat in dimension $\ge n$.  Then there exist:
\begin{enumerate}[label=\textup{(\roman*)}]
  \item a filtered inverse system of noetherian schemes $(S_\lambda)$ with affine transition maps and limit $S$;
  \item a stage $\lambda_0$ and data $f_{\lambda_0}\colon X_{\lambda_0}\to S_{\lambda_0}$ and $\M_{\lambda_0}$ of finite presentation descending $f$ and $\M$;
  \item a quasi-compact open $U_{\lambda_0}\subset S_{\lambda_0}$ descending $U$;
\end{enumerate}
such that, after enlarging $\lambda_0$ if necessary, the descended module $\M_{\lambda_0}|_{f_{\lambda_0}^{-1}(U_{\lambda_0})}$ is flat in dimension $\ge n$.
\end{proposition}

\begin{proof}
The existence of the inverse system and of a finite-presentation descent for $f$, $X$, and $\M$ is standard noetherian approximation.  The point is to preserve the flatness-in-dimension-$\ge n$ condition.

Choose a retrocompact open $V\subset f^{-1}(U)$ such that $\M|_V$ is flat over $U$ and such that the complement has fiber-dimension $<n$.  Because $V$ is quasi-compact open in a finitely presented $S$-scheme, it descends after increasing $\lambda_0$ to an open $V_{\lambda_0}\subset f_{\lambda_0}^{-1}(U_{\lambda_0})$.  Flatness of finite presentation modules is an open condition on the source, so after increasing $\lambda_0$ again, the set of points of $V_{\lambda_0}$ where $\M_{\lambda_0}$ is flat over $S_{\lambda_0}$ contains the inverse image of $V$ and therefore is itself an open neighborhood of the limit image of $V$; replacing $V_{\lambda_0}$ by this open flat locus, we may assume $\M_{\lambda_0}|_{V_{\lambda_0}}$ is flat.

It remains to control the fiber-dimension of the complement.  Since dimension of fibers of finitely presented morphisms is upper semicontinuous and compatible with filtered limits of affine transition systems, the inequality
\[
\dim\bigl((f^{-1}(U)\setminus V)/U\bigr)<n
\]
descends after increasing $\lambda_0$ to the corresponding inequality
\[
\dim\bigl((f_{\lambda_0}^{-1}(U_{\lambda_0})\setminus V_{\lambda_0})/U_{\lambda_0}\bigr)<n.
\]
Hence $\M_{\lambda_0}|_{f_{\lambda_0}^{-1}(U_{\lambda_0})}$ is flat in dimension $\ge n$.
\end{proof}

Once this approximation proposition is available, the proof of the qcqs version of flattening proceeds formally: one applies the noetherian theorem at the descended stage, then pulls the resulting admissible blow-up back along the limit map.  The compatibility of strict transform with base change and with composition then transports the flattening conclusion back to $S$.

\subsection{Detailed final induction for flattening in dimension $\ge n$}

We end with a more explicit induction schema for the theorem of flattening by admissible blow-up.

\begin{proposition}[Induction on the dimension parameter]\label{prop:flattening-induction}
Assume the flattening theorem is known for dimension thresholds strictly larger than $n$.  Then, after the noetherian reduction and after replacing $S$ by a suitable admissible blow-up over the already-good open $U$, one can arrange that the strict transform of $\M$ is flat in dimension $\ge n$.
\end{proposition}

\begin{proof}
Choose a total d\'evissage of $\M$ over a neighborhood of the bad locus for flatness in dimension $\ge n$.  The successive quotients $\Pcal_i$ have smaller fiber-dimension than $\M$.  By the induction hypothesis on the dimension threshold, after an admissible blow-up we may assume that the strict transforms of the relevant quotients are already flat in dimension $\ge n+1$.  Pulling the d\'evissage through the blow-up and replacing each term by its strict transform, the only remaining obstruction to flatness in dimension $\ge n$ lies in the possible failure of universal injectivity of the transition maps on points of exact fiber-dimension $n$.

But the locus where universal injectivity fails is represented by a closed condition on the base by the universal flattening theorem of Part~I.  Blowing up this closed locus and using \cref{prop:strict-composition}, we may assume that the transition maps become universally injective away from a set of fiber-dimension $<n$.  At that stage the flatness criterion from \cref{prop:inductive-flatness-criterion} applies and shows that the strict transform is flat at every point of fiber-dimension at least $n$.

It remains to verify that the torsion discarded in the strict transform is indeed confined to fiber-dimension $<n$.  This follows by tracing supports through the d\'evissage: any associated point of that torsion lying in dimension at least $n$ would either violate the already-established flatness in dimension $\ge n+1$ for the lower quotients, or else create a failure of universal injectivity for one of the transition maps at dimension exactly $n$, contradicting the choice of the last blow-up.  Therefore the discarded torsion is harmless for the threshold $n$, and the strict transform is flat in dimension $\ge n$.
\end{proof}

The preceding proposition does not replace the paper's proof, but it exposes the exact inductive logic in a form suitable for seminar presentation: first flatten lower-dimensional quotients, then blow up the residual failure of universal injectivity, and finally remove the exceptional torsion by passing to the strict transform.




\section{Long-form seminar reconstruction of Part~I}

This section is intentionally more expansive than the preceding chapters.  In the original article, many of the decisive moves are compressed into short but powerful arguments whose force becomes fully visible only after one rewrites them in slow motion.  For seminar purposes it is useful to separate the geometry from the homological algebra and to keep track of which statement is doing which job.  The main geometric input of Part~I is that a finite type morphism can, around a chosen point, be replaced by a finite morphism over a smoother base with geometrically integral fibers.  The main algebraic input is that once one has a short exact sequence with a free left term, flatness of the middle term is equivalent to a combination of universal injectivity and flatness of the quotient.  The theory of relative d\'evissage is precisely the mechanism that lets one iterate this observation a finite number of times.

\subsection{Why the local finite factorization is the correct geometric starting point}

Suppose that one begins with a morphism $f\colon X\to S$ locally of finite presentation and a point $x\in X$ above $s\in S$.  The naive impulse is to work directly on $X$, but this is not the right viewpoint.  The support of a finite module on $X$ may have several irreducible components passing through $x$, these components may have different relative dimensions over $S$, and the local ring of $X$ at $x$ may be too singular for the homological structure of the module to be visible.  What one really wants is a neighborhood in which the support becomes finite over a simpler ambient base and in which the point $x$ is isolated as the unique point over a closed base point.  This is exactly the content of the local finite factorization theorem.

The theorem is often read as a technical preliminary, but in fact it determines the whole shape of the later arguments.  Once the morphism is factored, after elementary \etale localization, through a finite map over a base $T$ which is smooth over $S$ with geometrically integral fibers, there are two decisive gains.  First, the fiber dimension becomes an honest dimension function on $T$ that behaves well under generization.  Second, free modules on $T$ acquire a rigid geometric meaning: because the fibers of $T\to S$ are geometrically integral, injectivity and rank considerations can be read on dense opens and then propagated.  The local main theorem of Zariski is therefore not merely providing an admissible replacement of $X$; it is producing the only setting in which universal injectivity can be tested inductively.

Let us rewrite the logic in a way suited for repeated use.  One starts with the closure of $x$ in the fiber $X_s$ and chooses an irreducible component of maximal dimension passing through $x$.  By the ordinary theorem of Zariski there exists, after shrinking and replacing by \etale neighborhoods, a finite morphism from a neighborhood of $x$ to an affine space over a neighborhood of $s$.  The coordinates on the affine space are then reorganized so that the image of the chosen component is dominant, and one discards those coordinates that are algebraically dependent on the support of the module.  This yields a smooth base $T$ of relative dimension equal to the dimension of the chosen component.  The module or the support under consideration is then finite over $T$.  In a seminar one may say: the role of the theorem is to straighten the support until it becomes finite over a smooth base of the right dimension.

The usefulness of the factorization is particularly transparent when one passes to henselian neighborhoods.  In that situation the chosen point becomes the only point over the closed base point, and finite morphisms break into disjoint unions according to the closed points of the special fiber.  The local decomposition of finite schemes over henselian rings is one of the silent workhorses of the whole paper: it allows one to isolate the relevant branch of the support and forget all the others.  Thus the geometry of the proof is already, in embryonic form, a geometry of selecting the correct branch.

\subsection{Detailed narrative proof of the existence of a relative d\'evissage}

We now explain carefully how one passes from the local factorization theorem to a relative d\'evissage.  Let $f\colon (X,x)\to (S,s)$ be locally of finite presentation and let $\M$ be a finite type $\OO_X$-module.  Put
\[
 n:=\dim_x(\Supp(\M)\cap X_s).
\]
The goal is to construct, after passing to pointed elementary \etale neighborhoods, a chain of pointed schemes
\[
(X',x')\xrightarrow{\pi_1}(T_1,t_1)\xrightarrow{\pi_2}\cdots\xrightarrow{\pi_r}(T_r,t_r)\to (S',s')
\]
and exact sequences
\[
0\to \Lcal_i\xrightarrow{\alpha_i}\N_i\to \Pcal_i\to 0
\]
with the following properties: each $\Lcal_i$ is free of finite rank on $T_i$, the module $\N_i$ is finite type, the support of $\Pcal_i$ has smaller relative dimension than the support of $\N_i$, and the terminal quotient $\Pcal_r$ is zero or has support of fiber dimension strictly less than zero, which means it vanishes.  The point of the construction is not simply to obtain exact sequences; rather, one wants them in such a way that the flatness of $\M$ may be read recursively from the universal injectivity of the maps $\alpha_i$.

The construction is by induction on $n$.  If $n<0$, then $x\notin\Supp(\M)$ and there is nothing to prove.  Suppose therefore that $n\ge 0$ and that the result is known for all fiber dimensions $<n$.  Apply the local finite factorization theorem to the support of $\M$ at $x$.  After replacing by elementary \etale neighborhoods we may assume that there exists a smooth morphism $(T,t)\to (S',s')$ of relative dimension $n$ and a finite morphism $\varphi\colon (Y,y)\to (T,t)$ such that the support of the pullback module $\M'$ on $Y$ is all of $Y$ and $Y$ maps to $X'$ by an \etale morphism with image containing the relevant part of the support.

Because $Y\to T$ is finite and $Y$ is of finite presentation over $T$, the direct image $\varphi_*\M'$ is a finite $\OO_T$-module.  Over the generic point of the fiber $T_s$, this module has finite rank.  By shrinking further on $T$ we may therefore choose a free module $\Lcal_1$ and a morphism
\[
\alpha_1\colon \Lcal_1\to \varphi_*\M'
\]
which is generically an isomorphism on the chosen component.  The cokernel $\Pcal_1$ is then supported on a closed subset whose fiber dimension at the chosen point is strictly smaller than $n$.  Pulling back to the pointed space over which the support was identified with the original one produces the first exact sequence of the d\'evissage.

Two comments are essential here.  First, the choice of $\Lcal_1$ is not arbitrary.  One wants a map that captures the generic rank of the module along the $n$-dimensional part of the support.  In practice one chooses sections generating the module at the generic points of the top-dimensional components and packages them as a free module.  Second, the dimension drop in the cokernel is the whole reason the construction terminates.  Without a strict drop in fiber dimension one would merely have written the module as a quotient of a free one, which is tautological.  What makes the argument nontrivial is that the failure of the free module to coincide with the original module is forced onto a lower-dimensional subset.

At this stage the induction hypothesis applies to the quotient $\Pcal_1$.  Indeed, the support of $\Pcal_1$ over the special fiber has dimension $<n$.  Hence, after a new elementary \etale localization and a new local finite factorization, one obtains a second exact sequence
\[
0\to \Lcal_2\xrightarrow{\alpha_2}\N_2\to \Pcal_2\to 0
\]
where the quotient has still smaller relative dimension.  Iterating finitely many times one arrives at the desired d\'evissage.

One can summarize the proof in a sentence that is useful for memory: take the top-dimensional part of the module, replace it by a free module of the same generic rank over a smooth base, and push the error term down by one dimension; then repeat.  The repeated \etale localization is not a nuisance but an intrinsic part of the geometry, because the support may branch differently at lower-dimensional strata.

\begin{proposition}[Dimension drop inside the d\'evissage]\label{prop:dimension-drop-long}
In the preceding construction, if the support of $\N_i$ has fiber dimension $n_i$ at the chosen point, then the support of $\Pcal_i$ has fiber dimension strictly less than $n_i$.
\end{proposition}

\begin{proof}
By construction the map $\alpha_i\colon \Lcal_i\to \N_i$ is chosen so as to become surjective at the generic points of the top-dimensional irreducible components of the support of $\N_i$ in the special fiber.  Therefore the stalk of $\Pcal_i$ at each such generic point is zero.  Since the support of $\Pcal_i$ is closed inside the support of $\N_i$, it follows that no irreducible component of the support of $\Pcal_i$ can have the same fiber dimension as the corresponding top-dimensional components of the support of $\N_i$.  Hence all irreducible components of $\Supp(\Pcal_i)$ have smaller fiber dimension.  Because fiber dimension is computed componentwise at generic points, the claim follows.
\end{proof}

This proposition is one of the places where the geometric meaning of generic freeness is easiest to see.  The free module $\Lcal_i$ is not used because free modules are pleasant in themselves; it is used because one can force them to coincide with $\N_i$ at the top-dimensional generic points.  Once that is achieved, the entire defect is condemned to live in smaller dimension.

\subsection{Transport of d\'evissage under base change and localization}

The d\'evissage would be of limited use if it were not stable under the natural operations that occur in the proof of flatness.  In practice one repeatedly restricts to opens, passes to local rings, base-changes along \etale morphisms, and replaces the base by henselizations or strict henselizations.  The point is that every such operation preserves the formal structure of a d\'evissage.

Let us spell out the mechanism.  Suppose
\[
0\to \Lcal\xrightarrow{\alpha}\N\to \Pcal\to 0
\]
is one step of a d\'evissage over $T\to S$.  If $T'\to T$ is any morphism, then pulling back gives an exact sequence
\[
0\to \Lcal_{T'}\xrightarrow{\alpha_{T'}}\N_{T'}\to \Pcal_{T'}\to 0
\]
provided the original map was universally injective or provided one works inside the formal construction of the d\'evissage before asking for flatness.  The free module stays free, finite presentation is preserved, and the support of the quotient pulls back to the support of the pulled-back quotient.  In particular, upper bounds on relative dimension are preserved.  This is why the entire proof of the flatness criterion may be carried out after replacing the base by a local or henselian neighborhood.

The compatibility with localization has another practical consequence.  When one wants to prove a pointwise statement, one may localize at the image point on the base and at the chosen point on the source without changing the d\'evissage data except by replacing all rings with their localizations.  Thus the whole theory naturally lives at the level of pointed schemes.  This is more than a stylistic convenience; it explains why the paper is able to formulate highly global statements while proving them by local algebra.

\section{A local algebra notebook for the flatness criterion}

The original paper uses a remarkable economy of means: once the geometric d\'evissage is available, the proof of the flatness criterion is extracted from a small package of homological lemmas.  In a long seminar, however, it is worthwhile to isolate that package and to prove it independently.  This section does exactly that.

\subsection{Universal injectivity revisited}

The expression \enquote{universal injectivity} can feel abstract when encountered for the first time.  The easiest way to understand it is to compare it with two more familiar notions.  Split injectivity means that a map remains injective after every base change for a trivial reason: it has a retraction.  Fiberwise injectivity means only that the map becomes injective after tensoring with residue fields.  Universal injectivity lies in between.  It is much weaker than split injectivity, but in the presence of flatness of the source or of the quotient it is much stronger than fiberwise injectivity.

\begin{lemma}[Universal injectivity from fiberwise injectivity and flat cokernel]\label{lem:ui-fiberwise-long}
Let $A$ be a ring and let
\[
0\to L\xrightarrow{u} N\to P\to 0
\]
be an exact sequence of $A$-modules with $P$ flat over $A$.  Assume that for every prime $\mathfrak p\subset A$ the map
\[
u\otimes_A k(\mathfrak p)\colon L\otimes_A k(\mathfrak p)\to N\otimes_A k(\mathfrak p)
\]
is injective.  Then $u$ is universally injective.
\end{lemma}

\begin{proof}
Let $A'$ be an $A$-algebra.  Tensoring the short exact sequence with $A'$ yields an exact sequence
\[
\Tor_1^A(P,A')\to L\otimes_A A'\xrightarrow{u\otimes A'} N\otimes_A A'\to P\otimes_A A'\to 0.
\]
Since $P$ is flat, the Tor term vanishes.  Hence $u\otimes A'$ is injective.  The assumptions on residue fields are not actually needed once the flatness of $P$ is known; they merely explain how one usually verifies the hypotheses in practice, namely by checking injectivity on fibers and then using the local criterion to conclude that $P$ is flat.
\end{proof}

The lemma is intentionally elementary, because it exposes the exact role played by universal injectivity in the paper.  Universal injectivity is not introduced as an independent geometric mystery; it appears because it is precisely the condition that allows the quotient to inherit flatness from the middle term and vice versa.  The proof reduces immediately to the Tor exact sequence.

\begin{lemma}[Flatness of the middle term]\label{lem:middle-flat-long}
Let
\[
0\to L\xrightarrow{u} N\to P\to 0
\]
be an exact sequence of $A$-modules.  Assume that $L$ and $P$ are flat over $A$.  Then $N$ is flat over $A$.  Conversely, if $N$ is flat and $u$ is universally injective, then $P$ is flat.
\end{lemma}

\begin{proof}
Tensor with an arbitrary $A$-module $Q$.  Since $L$ and $P$ are flat, the sequence
\[
0\to L\otimes_A Q\to N\otimes_A Q\to P\otimes_A Q\to 0
\]
remains exact, and so $\Tor_1^A(N,Q)=0$.  This proves the first statement.  For the converse, if $N$ is flat and $u$ is universally injective, then for each $Q$ the tensor sequence starts with $0\to L\otimes_A Q\to N\otimes_A Q$, hence $\Tor_1^A(P,Q)=0$ by the long exact Tor sequence.  Thus $P$ is flat.
\end{proof}

These two lemmas, together with the existence of the d\'evissage, already contain the whole proof of the flatness criterion.  But it is useful to add two supplementary observations.

\begin{lemma}[Local criterion in recursive form]\label{lem:recursive-local-flatness}
Let $A$ be local with residue field $k$, and let
\[
0\to A^{\oplus r}\xrightarrow{u} N\to P\to 0
\]
be exact with $P$ finite over an $A$-algebra of finite type.  If $u\otimes_A k$ is injective and $P$ is flat over $A$, then $u$ is universally injective and $N$ is flat.
\end{lemma}

\begin{proof}
By the preceding lemma it is enough to prove universal injectivity.  Because $P$ is flat, the quotient sequence remains exact after arbitrary base change, so $u$ is universally injective by \cref{lem:ui-fiberwise-long}.  Then \cref{lem:middle-flat-long} gives flatness of $N$.
\end{proof}

\begin{remark}
In the context of the paper the source of $u$ is not an arbitrary free module but one chosen so as to model the generic rank of $N$ along the top-dimensional strata.  The recursive criterion therefore says: to prove flatness, first force the top-dimensional part to look free, then prove that the defect is flat, then use universal injectivity to glue the two pieces back together.
\end{remark}

\subsection{The flatness criterion written as a recursive proof}

We now restate the flatness criterion in a form adapted to induction.  Let a total d\'evissage of $\M$ at $x$ be given by
\[
0\to \Lcal_i\xrightarrow{\alpha_i}\N_i\to \Pcal_i\to 0,\qquad 1\le i\le r,
\]
with $\Pcal_i$ identified with $\N_{i+1}$ for $1\le i<r$ and with $\Pcal_r=0$ in the terminal stage.  Then the theorem states that the following are equivalent:
\begin{enumerate}[label=\textup{(\alph*)}]
\item $\M_x$ is flat over $\OO_{S,s}$;
\item for all $i$, the map $\alpha_i$ is universally injective and the terminal quotient is flat;
\item for all $i$, the fiber map of $\alpha_i$ at the closed point is injective and the terminal quotient is flat.
\end{enumerate}
The equivalence of the last two formulations is a consequence of the local algebra just established.  The equivalence with flatness is proved by descending induction on $i$.

Assume first that $\M_x$ is flat.  Then $\N_1$ is flat.  Since $\Lcal_1$ is free, the quotient $\Pcal_1$ is flat by \cref{lem:middle-flat-long}.  The map $\alpha_1$ is universally injective because the exact sequence remains exact after arbitrary base change.  Repeating the argument for the shorter d\'evissage of $\Pcal_1$ yields the conditions for all indices $i$.  This is the easy direction.

For the converse, assume by induction that $\Pcal_1$ is flat.  Then the first short exact sequence has flat left and right terms, and $\alpha_1$ is universally injective.  Therefore $\N_1$ is flat by \cref{lem:middle-flat-long}, and hence so is $\M_x$.  The real content is that the inductive hypothesis applies because the entire tail of the d\'evissage of $\Pcal_1$ is part of the original d\'evissage of $\M$.  In other words, the proof is not really a proof by induction on the length of an arbitrary exact sequence; it is a proof by induction on the geometric complexity of the support.

This recursive proof is one of the places where the paper's unity becomes fully visible.  The geometry produced the d\'evissage.  The local algebra showed that universal injectivity is the correct replacement for split injectivity.  The theorem itself is just the systematic iteration of these two observations.

\section{A geometric notebook for purity and the relative assassin}

The transition from flatness to projectivity is one of the most subtle moments in the article.  A reader who only remembers the statement \enquote{flat plus pure implies projective} has not yet absorbed the geometric content.  The word \enquote{pure} is not decorative.  It encodes a statement about where the associated points of the module are allowed to lie.  Since projectivity over the base is a condition of local constancy, one cannot hope to obtain it from flatness alone; one must also rule out vertical or parasitic associated components.  Purity is exactly the mechanism for doing so.

\subsection{Associated points and closures}

Let $A$ be a ring, $B$ an $A$-algebra of finite presentation, and $M$ a finite $B$-module flat over $A$.  An associated prime of $M$ corresponds geometrically to a point of the support where the module has a section annihilated by the corresponding prime ideal.  These points are the natural receptacles for torsion phenomena.  If one wants to know whether $M$ behaves like a projective module over the base, it is therefore natural to inspect the closures of these points relative to the fibers.

Purity along the fiber over $s\in S$ says that every associated point of the henselian pullback of the module lies in the closure of the special fiber.  Equivalently, every associated point specializes to some point on the closed fiber.  This is stronger than saying that the support meets the fiber, because associated points detect the hidden embedded structure of the module.  In practice purity excludes the possibility that some associated point exists only in the generic region and disappears before reaching the special fiber.

Why is this the right condition for projectivity?  Imagine a module that is flat over the base and locally free on a dense open, but whose associated points acquire an extra component entirely contained in a vertical divisor.  Such a component may not destroy flatness immediately, yet it obstructs the uniform projective behavior one expects from a true vector bundle.  Purity says that no such hidden component is allowed unless it is already controlled by the special fiber.  The module is therefore forced to be assembled from the same generic rank data everywhere.

\begin{proposition}[A closure criterion for purity]\label{prop:closure-purity-long}
Let $A$ be henselian local, let $B$ be an $A$-algebra of finite presentation, and let $M$ be a finite $B$-module flat over $A$.  Then the following are equivalent:
\begin{enumerate}[label=\textup{(\roman*)}]
\item $M$ is pure over $A$;
\item for every associated point $\xi\in\Ass_B(M)$, the closure of $\{\xi\}$ in $\Spec(B)$ meets the closed fiber $\Spec(B\otimes_A k(s))$;
\item after passing to a suitable elementary \etale neighborhood whose image contains all associated points above the chosen fiber, the pullback of $M$ becomes projective over the base.
\end{enumerate}
\end{proposition}

\begin{proof}
The equivalence of the first two formulations is simply the definition rewritten geometrically.  Indeed, purity says exactly that every associated point lies in the generization of the special fiber; saying that a point lies in the generization of the special fiber is equivalent to saying that its closure meets that fiber.

We prove that (ii) implies (iii).  Let $Z$ be the union of the closures of the associated points of $M$.  By assumption $Z$ meets the closed fiber in every irreducible component.  Using the local finite factorization theorem on each relevant component of the support, and shrinking if necessary, one constructs an elementary \etale neighborhood $Y\to X$ such that the image of $Y$ contains all points of $Z$ lying above the chosen base point.  Over $Y$ the support is finite over a smooth base of constant relative dimension.  Because the associated points are all captured by the image, the pullback module on $Y$ has no hidden embedded components outside that image.  By the projectivity criterion for modules over smooth bases developed earlier in the paper, one concludes that the pullback is projective over the base.

Finally, (iii) implies (ii) because projective modules are pure.  More explicitly, if the pullback $M_Y$ is projective over the base, then every associated point of $M_Y$ specializes to the closed fiber.  Since the image of $Y$ contains the associated points of $M$, the same specialization property descends to $M$.
\end{proof}

This proposition shows that purity is not an abstract extra hypothesis but a geometric statement whose content is visible on the closures of associated points.  In seminar language: purity means that the assassin stays anchored to the special fiber.

\subsection{From purity to projective descent}

We now rewrite in slow motion the argument that flatness together with purity forces projectivity.  The first step is to pass to an \etale neighborhood that captures the relative assassin.  On that neighborhood the module becomes projective over the base.  The second step is to descend projectivity from the neighborhood back to the original space.  The nontrivial issue is to show that the comparison map from the original module to the pushforward from the neighborhood is universally injective and that its cokernel vanishes because any residual support would miss the assassin.

\begin{theorem}[Flat plus pure implies projective: expanded proof]\label{thm:flat-pure-projective-long}
Let $A\to B$ be of finite presentation and let $M$ be a finite $B$-module flat over $A$.  If $M$ is pure over $A$, then $M$ is projective over $A$.
\end{theorem}

\begin{proof}
Because projectivity is local on the base, we may suppose that $A$ is henselian local.  Let $X:=\Spec(B)$ and write $\widetilde M$ for the associated quasi-coherent sheaf.  By purity and \cref{prop:closure-purity-long}, there exists an elementary \etale neighborhood $\pi\colon Y\to X$ whose image contains the relative assassin of $\widetilde M$ and such that the pullback module $M_Y:=\Gamma(Y,\pi^*\widetilde M)$ is projective over $A$.

Consider the adjunction morphism
\[
\theta\colon M\longrightarrow \Gamma(Y,\pi^*\widetilde M).
\]
We claim first that $\theta$ is universally injective.  Since the image of $Y\to X$ contains the relative assassin, any element of $M$ vanishing after pullback to $Y$ is supported on a closed subset disjoint from all associated points of $M$.  For a finite module, any nonzero submodule has an associated point.  Therefore the kernel of $\theta$ has empty assassin and must vanish.  The same argument applies after arbitrary base change, because the image of the corresponding base-changed \etale neighborhood still contains the relative assassin of the base-changed module.  Thus $\theta$ is universally injective.

Let $Q$ be its cokernel.  We shall show that $Q=0$.  Since the target is $A$-projective, hence $A$-flat, and the source is $A$-flat by hypothesis, the quotient $Q$ is $A$-flat by \cref{lem:middle-flat-long}.  If $Q$ were nonzero, it would have an associated point $\xi$ on $X$.  By flatness and purity of $Q$ over the local base, the closure of $\xi$ would meet the closed fiber.  But by construction of $\theta$, the module $Q$ is supported on the complement of the image of $Y$, while that image contains the relative assassin of $M$ and therefore all possible associated points that could survive on the closed fiber.  This contradiction shows that $Q=0$.

Hence $\theta$ is an isomorphism.  Since its target is projective over $A$, so is $M$.
\end{proof}

A useful way to remember this proof is the following.  Purity allows one to choose an \etale chart that sees every associated point.  On that chart the module is projective.  The descent map from the original module to the chart has no kernel because the kernel would be invisible to the assassin, and it has no cokernel because the cokernel would create a new associated point outside the chart.  Thus the module coincides with a projective one.

\section{The universal flattening condition in affine detail}

In the article the universal flattening condition first appears in a slightly specialized form: one studies the functor of base changes on which a quotient becomes bijective, and then deduces a representability statement for flatness.  The proof is elegant because it reduces a geometric condition to a statement about the vanishing of coordinates of the kernel of a morphism between finite modules.  This section rewrites that argument in full affine detail.

\subsection{Representability of the bijectivity locus}

Let $A$ be a ring, let $X$ be an $A$-scheme, let $\M$ be an $A$-flat and $A$-pure quasi-coherent module on $X$, and let $u\colon \M\to \N$ be a quotient.  We want to describe the functor sending an $A$-algebra $A'$ to the assertion that the pulled-back morphism $u_{A'}$ is an isomorphism.  The theorem says this functor is representable by a closed subscheme of the base.

The proof begins by using purity to reduce to the case in which $X$ admits an \etale cover $Y\to X$ such that the pullback of $\M$ to $Y$ is free over $A$ and such that the image of $Y$ contains the relative assassin of $\M$.  Because the image contains the assassin, bijectivity of $u$ may be checked on $Y$.  This is the only genuinely geometric part.  Once the reduction is made, one may write
\[
M\cong A^{\oplus r},\qquad 0\to R\to M\to N\to 0.
\]
Then $u_{A'}$ is bijective if and only if $R\otimes_A A'=0$.  If one chooses a basis of $M$, the submodule $R$ is generated by a family of vectors with coordinates in $A$.  Let $I\subset A$ be the ideal generated by all these coordinates.  Then $R\otimes_A A'=0$ if and only if $I A'=0$.  The functor is therefore represented by $\Spec(A/I)$.  The moral is that the universal flattening condition is already encoded in the coefficients of the relations that obstruct the quotient from being an isomorphism.

\begin{proposition}[Affine representability with explicit equations]\label{prop:affine-representability-long}
Let $A$ be a ring, $M$ a finite free $A$-module, and $u\colon M\twoheadrightarrow N$ a quotient.  Let $R=\Ker(u)$ and let $I\subset A$ be the ideal generated by the coordinates of a finite set of generators of $R$ relative to a basis of $M$.  Then for any $A$-algebra $A'$ one has
\[
u_{A'}\text{ is an isomorphism }\Longleftrightarrow I A'=0.
\]
In particular the functor of base changes on which $u$ becomes an isomorphism is represented by $\Spec(A/I)$.
\end{proposition}

\begin{proof}
If $u_{A'}$ is an isomorphism, then its kernel $R\otimes_A A'$ vanishes.  Therefore every chosen generator of $R$ maps to zero in $M\otimes_A A'$, and hence each of its coordinates vanishes in $A'$.  Thus $IA'=0$.

Conversely, suppose $IA'=0$.  Every chosen generator of $R$ has all coordinates zero in $A'$, hence the whole submodule $R\otimes_A A'$ vanishes.  Therefore the kernel of $u_{A'}$ is zero.  Since $u$ was already surjective and surjectivity is preserved by tensor product, $u_{A'}$ is an isomorphism.
\end{proof}

The proof is so simple that it risks concealing its significance.  What it shows is that the representability of the universal flattening condition is not mysterious once one has moved to the right chart.  Purity provides the chart; the rest is linear algebra.

\subsection{Passage from bijectivity to flatness}

The step from the bijectivity locus to the flatness locus is conceptually direct but technically delicate.  One chooses a presentation
\[
F\twoheadrightarrow M
\]
with $F$ free.  The quotient becomes flat after base change exactly when the kernel of this presentation becomes universally injective into $F$ and the corresponding defect module vanishes in the sense represented by the previous proposition.  In a local-henselian situation these conditions may be encoded by a closed subscheme because the relevant universal injectivity and cokernel conditions are all closed under the hypotheses of finite presentation.

A useful way to phrase the argument is the following.  The flatness criterion from Part~I says that flatness is equivalent to the existence of a d\'evissage whose transition maps are universally injective and whose terminal quotient is flat.  Each one of these conditions is representable: universal injectivity of a finite presentation map can be tested on finite cokernels and therefore on the vanishing of coordinates after passing to suitable free charts, while flatness of the terminal quotient is treated recursively.  Thus the universal flattening condition is not represented in one stroke but obtained as the intersection of finitely many closed conditions extracted from the d\'evissage.

\section{Extended seminar notes on strict transform and blow-up}

Part~II of the article is often quoted only for the final theorem, but the proof contains a number of ideas that deserve to be isolated.  The first is the strict transform of a module.  The second is the projective prototype using Quot schemes.  The third is the induction on the dimension threshold.  The theorem emerges only after these three ideas are made to work together.

\subsection{Strict transform as maximal quotient agreeing over the good open}

Let $g\colon S'\to S$ be a $U$-admissible blow-up, let $f\colon X\to S$ be of finite presentation, and let $X':=X\times_S S'$.  Write $j\colon f^{-1}(U)\hookrightarrow X$ and $j'\colon (f')^{-1}(U)\hookrightarrow X'$ for the open immersions.  For a finite type $\OO_X$-module $\M$ the strict transform is defined by
\[
\st_g(\M):=\Image\bigl(g_X^*\M\to j'_*j'^*g_X^*\M\bigr).
\]
The formula is compact, but in seminar it is worth translating it into words.  Over the good open $U$ the blow-up is an isomorphism, so the pullback $g_X^*\M$ agrees there with the original module.  Away from $U$, however, the pullback may acquire extra torsion supported on the exceptional divisor.  The strict transform is obtained by discarding exactly the largest subsheaf of the pullback that vanishes on the good open.  In other words, it is the largest quotient of the pullback that still agrees with the original module over $U$.

\begin{lemma}[Maximality of the strict transform]\label{lem:strict-maximal-long}
With notation as above, $\st_g(\M)$ is characterized by the following universal property: it is the maximal quotient $Q$ of $g_X^*\M$ such that $Q|_{(f')^{-1}(U)}\cong g_X^*\M|_{(f')^{-1}(U)}$.
\end{lemma}

\begin{proof}
Let $K$ be the kernel of the natural map $g_X^*\M\to j'_*j'^*g_X^*\M$.  By definition $\st_g(\M)=(g_X^*\M)/K$.  Since $j'^*K=0$, the quotient agrees with $g_X^*\M$ over the open $(f')^{-1}(U)$.  If $Q$ is any other quotient with this property, then the kernel of $g_X^*\M\to Q$ vanishes after restriction to $(f')^{-1}(U)$, hence lies in $K$.  Therefore the quotient map factors through $(g_X^*\M)/K=\st_g(\M)$.  This proves maximality.
\end{proof}

The lemma explains why strict transform is the correct object for flattening.  One cannot simply pull the module back to the blow-up, because the blow-up may create exceptional torsion unrelated to the original defect of flatness.  One also cannot quotient by an arbitrary exceptional subsheaf, because that would destroy the agreement with the good open.  The strict transform is the unique quotient balancing these two requirements.

\subsection{The projective prototype through Quot schemes}

Assume now that $S$ is noetherian, that $X\to S$ is projective, and that $\M$ is a coherent module on $X$ which is already flat over a dense open $U\subset S$.  In this case one can form the relative Quot scheme parameterizing quotients of $\M$ with fixed Hilbert polynomial on the fibers.  The quotient that equals $\M$ over $U$ defines a section of the Quot scheme over $U$.  Taking the scheme-theoretic closure of the image of that section and projecting to $S$ gives a projective birational morphism $g\colon S'\to S$ which is an isomorphism over $U$.  The universal quotient on the Quot scheme pulls back to a quotient of $g_X^*\M$ on $X':=X\times_S S'$; by construction this quotient is flat over $S'$ and agrees with the original module over $U$.

This argument already proves a version of the flattening theorem in the projective noetherian case.  It is useful to pause and interpret it.  The Quot scheme is not merely a parameter space for quotients; in the present application it records all possible ways of removing the bad part of the module while preserving the generic behavior over $U$.  The closure of the section over $U$ chooses the unique birational modification on which the generic quotient extends everywhere.  In other words, the Quot scheme prototype shows that flattening is a compactification problem for the family of quotients that is already visible on the flat locus.

The general theorem does not proceed by Quot schemes globally because one does not want to impose projectivity hypotheses on $X\to S$.  Nevertheless the projective prototype is conceptually decisive: it tells us what sort of birational modification to expect and why strict transform rather than ordinary pullback is the correct module on the modification.

\subsection{Induction on the dimension threshold}

The theorem of Raynaud and Gruson is formulated in terms of flatness in dimension $\ge n$.  This may look at first like an artificial weakening, but in fact it is the natural scale on which the induction works.  If one attempted to flatten the whole module at once, one would have no room to separate the top-dimensional defect from the lower-dimensional one.  By asking only for flatness in dimension $\ge n$, one may first repair the module away from smaller-dimensional subsets and then descend the threshold step by step.

Here is the pattern.  Suppose one knows how to flatten in dimension $\ge n+1$.  Apply this to the successive lower-dimensional quotients that arise in a total d\'evissage of $\M$.  After an admissible blow-up one may then suppose these quotients are flat in dimension $\ge n+1$.  The flatness criterion reduces the remaining obstruction for $\M$ to the failure of universal injectivity of certain transition maps at points of exact fiber dimension $n$.  But the locus where universal injectivity fails is represented by a closed condition on the base.  Blowing up this closed locus and replacing pullback by strict transform kills precisely the residual exceptional torsion.  One then deduces that the strict transform is flat in dimension $\ge n$.

This description shows that the induction is not performed on the module alone but on the pair consisting of the module and a dimension threshold.  The threshold keeps track of the part of the support that has already been controlled.  The strict transform then removes whatever new torsion appears when one blows up the residual bad locus.

\begin{proposition}[Detailed induction step]\label{prop:induction-step-long}
Let $f\colon X\to S$ be of finite presentation with $S$ noetherian and let $\M$ be a finite type $\OO_X$-module.  Assume that for every finite type module whose support has fiber dimension $<n$ the flattening theorem is known.  Then there exists a $U$-admissible blow-up $g\colon S'\to S$ such that $\st_g(\M)$ is flat in dimension $\ge n$.
\end{proposition}

\begin{proof}
Choose a total d\'evissage of $\M$ over a neighborhood of the nonflat locus in dimension $\ge n$.  The quotients $\Pcal_i$ have smaller fiber dimension than the corresponding $\N_i$.  By the induction hypothesis applied repeatedly, after a first admissible blow-up we may assume that the strict transforms of all the relevant quotients are flat in dimension $\ge n+1$.  Pull the whole d\'evissage through this blow-up and replace each term by its strict transform.  The lower-dimensional quotients now satisfy the flatness conclusion needed in the recursive criterion.

What remains is to ensure universal injectivity of the transformed maps $\alpha_i$ at points of exact fiber dimension $n$.  By the representability theorem of Part~I the locus on the base where a given $\alpha_i$ fails to be universally injective is closed.  Taking the union of these finitely many loci and blowing it up admissibly, we may suppose that after strict transform the maps become universally injective outside a subset whose inverse image in $X'$ has fiber dimension $<n$.  The recursive flatness criterion then implies that the strict transform of $\M$ is flat at every point of fiber dimension at least $n$.

Finally one must check that the passage to strict transform has not thrown away any component relevant to the threshold $n$.  If an associated point of the discarded torsion had fiber dimension at least $n$, then either it would project to a point where one of the lower-dimensional quotients was not yet flat in dimension $\ge n+1$, or it would witness failure of universal injectivity of one of the transition maps at dimension $n$.  Both possibilities were excluded by construction.  Hence the discarded torsion is supported in fiber dimension $<n$, and the theorem follows.
\end{proof}

\subsection{Why the qcqs approximation step is indispensable}

The noetherian argument just described is already the heart of the proof, but the theorem is stated for quasi-compact quasi-separated bases.  The extension to the qcqs setting is not an ornamental generality.  Flattening is intended as a structural operation in algebraic geometry, and many natural bases arise as limits or as spaces for which no noetherian hypothesis is available.  One therefore needs a descent argument that preserves not only the schemes and modules but also the condition of flatness in dimension $\ge n$.

The approximation step works because every piece of the data is of finite presentation.  The scheme $S$ is written as a filtered inverse limit of noetherian schemes $S_\lambda$ with affine transition morphisms; the source $X$, the morphism $f$, the module $\M$, the chosen open $U$, and the retrocompact open on which flatness is already known all descend to some finite stage.  Upper semicontinuity of fiber dimension and openness of flatness for finitely presented modules ensure that the inequality defining flatness in dimension $\ge n$ also descends after increasing the stage.  One then applies the noetherian theorem at that stage and pulls the resulting blow-up back to the limit.  The compatibility of strict transform with base change finishes the argument.

The key point is that the approximation is not merely an approximation of spaces; it is an approximation of the whole proof package: the good open, the d\'evissage, the dimension bounds, and the eventual blow-up all descend.

\section{Examples and seminar models}

A long seminar on Raynaud--Gruson benefits enormously from examples, not because the main theorems are elementary, but because examples show exactly where the different notions part company.  In particular one wants to see the distinction between pullback and strict transform, between flatness and projectivity, and between a module whose defect is visible on the support and one whose defect is hidden in associated points.

\subsection{A torsion example over a discrete valuation ring}

Let $A$ be a discrete valuation ring with uniformizer $\pi$, and let
\[
B=A[x],\qquad M:=B/(\pi x).
\]
The module $M$ is not flat over $A$, because multiplication by $\pi$ on $M$ kills the nonzero element represented by $x$.  The support of $M$ is the union of the divisor $\pi=0$ and the divisor $x=0$, but the flatness defect is concentrated at their intersection in the special fiber.  This is already a good toy model for the philosophy of the paper: the module is generically free of rank one on the open where either $\pi\neq 0$ or $x\neq 0$, yet it fails to be flat at a codimension-two point.

Blow up the base $\Spec(A)$ at the closed point; since the base is already one-dimensional, this blow-up is trivial as a scheme.  The example therefore shows that not every flatness defect can be repaired by modifying only a one-dimensional base unless one also changes the source or the presentation.  What the general theorem says is subtler: after pulling back to a nontrivial admissible blow-up in higher-dimensional situations and passing to strict transform, the exceptional torsion created by the pullback can be discarded, and the remaining quotient becomes flat.

If one insists on seeing strict transform in this elementary example, one may instead work with a two-dimensional base $A=k[u,v]_{(u,v)}$ and the module $M=A[x]/(ux,vx)$.  The module is generically free of rank one away from the origin, but at the origin the element $x$ is killed by the ideal $(u,v)$.  Blowing up the origin separates the two branches of the bad locus.  On the chart where $v=ut$ the pullback module is
\[
A[u,t,x]/(ux,utx)=A[u,t,x]/(ux),
\]
so the exceptional torsion is still carried by $x$ but is now visibly supported on the exceptional divisor.  Passing to strict transform discards that torsion and leaves a flat quotient.  This is the affine shadow of the general theorem.

\subsection{A purity example}

Let $A$ be a domain and let $B=A[y]/(fy)$ for some nonzero $f\in A$.  Consider the $B$-module $M=B$.  As an $A$-module, $M$ is rarely flat unless $f=0$, because the element $y$ is annihilated by $f$.  The associated point defined by $(y)$ lies over the divisor $f=0$.  Now compare this with the module
\[
N:=A\oplus A
\]
regarded as a $B$-module through the quotient $B\twoheadrightarrow A$ sending $y$ to zero.  The module $N$ is projective over $A$, hence pure, and its associated points are exactly the generic points that already specialize to every fiber.  The contrast between $M$ and $N$ shows what purity is ruling out: the module $M$ has a vertical associated component tied to the vanishing of $f$, whereas $N$ has no such component.

A slightly richer example is obtained by taking $A=k[t]_{(t)}$ and
\[
B=A[z]/(tz(z-1)).
\]
The module $B$ has two components in the special fiber, one corresponding to $z=0$ and the other to $z=1$, but both are anchored to the closed fiber and therefore visible to the assassin over the local base.  The geometry is reducible, yet the associated points specialize correctly.  This illustrates that purity is not the same as irreducibility or smoothness; it is a condition on the behavior of associated points relative to the fibers.

\subsection{A d\'evissage example in affine coordinates}

Take $A=k[t]_{(t)}$, $B=A[x]$, and
\[
M:=B/(x^2,tx).
\]
The support of $M$ is the line $x=0$, but the module has an embedded nilpotent structure at the closed point.  A useful first step of d\'evissage is the exact sequence
\[
0\to A\xrightarrow{\alpha} M\to A/(t)\to 0,
\]
where $\alpha(1)$ is represented by the class of $x$.  The quotient is supported on the closed fiber and therefore has smaller relative dimension.  This little exact sequence already contains the pattern of the general theory: a free module captures the top-dimensional generic part, and the quotient records a lower-dimensional defect.  Flatness fails precisely because the map $\alpha$ is not universally injective; after tensoring with $A/(t)$ the element $x$ dies.  Thus the example is a concrete incarnation of the flatness criterion.

\section{Narrative synthesis and reading guide}

At the end of a long exposition it is useful to return once more to the central thread of the paper.  The sequence of ideas may be summarized as follows.

The first problem is to understand a finite module near a chosen point.  The module is replaced, after \etale localization, by one that is finite over a smooth base of controlled relative dimension.  This is the geometric straightening step.

The second problem is to extract from that geometric model a finite recursive description of the module.  This is the d\'evissage: at each stage a free module captures the generic rank of the current module, and the quotient has smaller relative dimension.  The module is thus decomposed into finitely many pieces arranged by decreasing geometric complexity.

The third problem is to recognize flatness from this decomposition.  The answer is that flatness is equivalent to universal injectivity of the transition maps plus flatness of the last quotient.  The theorem is homological, but it becomes effective only because the d\'evissage has already transformed the geometry into a finite algebraic recursion.

The fourth problem is to understand when flatness upgrades to projectivity.  Flatness alone does not control hidden associated components.  Purity is introduced precisely to control them.  Once every associated point is forced to remain anchored to the relevant fibers, the module may be recovered from an \etale chart on which it is projective.  This yields the characterization flat $+$ pure $=$ projective.

The fifth problem is functorial.  Instead of asking simply whether a module is flat, one asks for the base changes on which it becomes flat.  Because the d\'evissage reduces flatness to a finite list of universal injectivity and quotient conditions, and because those conditions are representable in affine coordinates, one obtains the universal flattening condition.

The sixth and final problem is birational.  If the module is not flat, can one alter the base so that it becomes flat?  The answer is yes after admissible blow-up and passage to strict transform.  The blow-up isolates the bad locus on the base.  The pullback to the blow-up creates exceptional torsion.  The strict transform discards exactly that torsion and preserves the original module over the good open.  Iterating this procedure by descending dimension thresholds yields the flattening theorem.

Read in this way, the article is not a collection of separate criteria but a single argument moving through four registers: local algebra, geometry of supports, functorial representability, and birational modification.  That unity is the reason the paper remains so influential.




\section{Proof dossier for the local geometric preparation}

The purpose of the present dossier is to slow down the arguments that in the original article are absorbed into the phrase \enquote{after passage to a suitable elementary \etale neighborhood}.  The phrase is elegant, but in a seminar one repeatedly needs to know exactly what is being gained at each such passage.  We therefore record a sequence of intermediate lemmas that clarify why the local finite factorization theorem is strong enough for the rest of the paper.

\subsection{Choosing the correct irreducible component}

Let $f\colon X\to S$ be locally of finite presentation, let $x\in X$ lie over $s\in S$, and let $\M$ be a finite type $\OO_X$-module.  Write $Z:=\Supp(\M)$.  The first subtlety is that the local dimension relevant for the d\'evissage is not the ordinary Krull dimension of $Z$ at $x$ but the dimension of the fiber $Z_s$ at $x$.  This is because flatness is a relative property, and the induction has to measure complexity in directions transverse to the base.

Choose an irreducible component $C$ of $Z_s$ passing through $x$ whose dimension is maximal among the components of $Z_s$ through $x$.  The geometric problem is to create a neighborhood in which this component becomes finite over a smooth base of the same dimension.  The smooth base should be thought of as a coordinate system along the top-dimensional directions of the support.

The method is classical.  One first embeds a neighborhood of $x$ into an affine space over a neighborhood of $s$.  Then one projects onto a smaller affine space by choosing a sufficiently general set of coordinates.  If the projection is chosen so that its differential is generically injective on the component $C$, then the induced map from a neighborhood of the generic point of $C$ to the smaller affine space is quasi-finite.  By the main theorem of Zariski, after shrinking one obtains a factorization through a finite morphism followed by an open immersion.  The open immersion is then absorbed into a smooth neighborhood of the target point.  The result is the desired finite morphism to a smooth base.

This familiar argument is used twice in the paper.  The first time it creates the first stage of the d\'evissage.  The second time it is used again on the support of the quotient after the first stage, now with smaller fiber dimension.  Thus one should not imagine the local factorization theorem as a single preliminary tool; it is a renewable resource used at every step of the induction.

\begin{lemma}[Finite presentation and preservation of geometric integrality]\label{lem:geom-integral-long}
Let $T\to S$ be smooth at a point $t\mapsto s$, and let $Y\to T$ be finite near a point $y\mapsto t$.  Suppose the local ring of the chosen irreducible component of $Y_s$ at $y$ is a domain and that the generic rank of $\OO_Y$ over $\OO_T$ along that component is one.  Then after shrinking around $t$ and $y$ one may assume that the relevant fibers of $T\to S$ are geometrically integral and that $Y\to T$ is generically an isomorphism on the chosen component.
\end{lemma}

\begin{proof}
Smoothness implies that after shrinking on $T$ the fibers of $T\to S$ are geometrically regular, hence geometrically reduced, and each irreducible component is geometrically integral.  Because the generic rank of $Y\to T$ along the chosen branch is one, the induced map of function fields is an isomorphism on that branch.  By shrinking to remove the other components and using the fact that finiteness is stable under shrinking, one obtains the claim.
\end{proof}

The lemma clarifies what the phrase \enquote{geometrically integral fibers} is doing in the proof of the flatness criterion.  If the ambient smooth base had reducible or badly behaved fibers, then a map from a free module could vanish on one branch and not on another, and fiberwise injectivity would not propagate cleanly.  Geometric integrality is the property that makes rank arguments morally one-dimensional on each branch.

\subsection{The first step of the d\'evissage as a controlled approximation}

There is another way to describe the first step of the d\'evissage that is often pedagogically useful.  Once the support has become finite over a smooth base $T$ of relative dimension $n$, the direct image of the module to $T$ is a finite module on a regular enough ambient space.  The task is then to approximate it by a free module without changing the generic behavior.  One may think of this as a controlled approximation problem: choose finitely many sections that form a basis at the generic points of the top-dimensional part, and assemble them into a free module $L$.  The induced map $L\to N$ is then an isomorphism at the generic points of maximal dimension.  The quotient measures the failure of those sections to generate everywhere.

A common concern at this point is whether the quotient could still have the same top-dimensional support because of embedded components.  The answer is no precisely because the generators were chosen to form a basis at the generic points of the top-dimensional irreducible components.  An embedded component cannot become a new maximal component after the generic points have been killed.  This is why the dimension drop is strict.

\begin{proposition}[Controlled approximation by a free module]\label{prop:controlled-approximation}
Let $T$ be a scheme, let $N$ be a finite type $\OO_T$-module, and let $\Xi$ be a finite set of generic points of irreducible components of $\Supp(N)$.  Then there exists a free module $L$ of finite rank and a morphism $L\to N$ which is surjective at every point of $\Xi$.  If the points of $\Xi$ are precisely the generic points of the top-dimensional irreducible components of the special fiber of $\Supp(N)$, then the cokernel is supported in smaller fiber dimension.
\end{proposition}

\begin{proof}
For each $\xi\in \Xi$, the stalk $N_\xi$ is a finite module over the field $\kappa(\xi)$ and hence has a finite basis.  Choosing representatives of these basis elements in sections over neighborhoods of the $\xi$ and shrinking if necessary, one obtains finitely many local sections that generate $N$ at all points of $\Xi$.  As $\Xi$ is finite, the sections may be collected into a single free module $L$ of finite rank mapping to $N$.

If $\Xi$ is the set of generic points of the top-dimensional irreducible components of the special fiber of the support, then the cokernel vanishes at all such generic points.  Therefore none of those top-dimensional components lies in the support of the cokernel.  Consequently the support of the cokernel has smaller fiber dimension.
\end{proof}

This proposition is elementary, but it expresses the entire philosophy of d\'evissage in a compact form: choose a free module that is exact where the dimension is maximal, and let the quotient remember what happens in lower dimension.

\subsection{A complete recursive proof of the existence of a total d\'evissage}

We now collect the preceding ideas into a single recursive proof, since the existence theorem is the starting point for everything later.

\begin{theorem}[Existence of a total relative d\'evissage: long proof]\label{thm:devissage-complete-long}
Let $f\colon (X,x)\to (S,s)$ be locally of finite presentation and let $\M$ be a finite type $\OO_X$-module.  Then, after passing to a pointed elementary \etale neighborhood of $(X,x)\to(S,s)$, there exists a finite sequence of short exact sequences
\[
0\to \Lcal_i\xrightarrow{\alpha_i}\N_i\to \Pcal_i\to 0,
\qquad 1\le i\le r,
\]
with $\Lcal_i$ free of finite rank, $\N_1$ corresponding to $\M$, $\N_{i+1}=\Pcal_i$ for $1\le i<r$, and the support of $\Pcal_i$ having strictly smaller fiber dimension than the support of $\N_i$ at the chosen point.  In particular the process terminates after finitely many steps.
\end{theorem}

\begin{proof}
We argue by induction on
\[
 n:=\dim_x(\Supp(\M)\cap X_s).
\]
If $n<0$, then $x\notin\Supp(\M)$ and there is nothing to show.

Assume now $n\ge 0$.  Apply the local finite factorization theorem to the support $Z:=\Supp(\M)$ and to a maximal-dimensional irreducible component of $Z_s$ through $x$.  After replacing by an elementary \etale neighborhood, we obtain a finite morphism $Y\to T$ with $T\to S$ smooth of relative dimension $n$ and geometrically integral fibers such that the image of $Y$ in $X$ contains a neighborhood of the chosen branch of the support.

Push the pulled-back module to $T$ and denote the resulting finite module by $N_1$.  By \cref{prop:controlled-approximation} there exists a finite free module $L_1$ and a morphism $\alpha_1\colon L_1\to N_1$ surjective at the generic points of the top-dimensional irreducible components of the special fiber of the support.  Let $P_1$ be the cokernel.  By the same proposition, the support of $P_1$ has smaller fiber dimension at the chosen point.

If $P_1$ vanishes, the d\'evissage has length one.  Otherwise we apply the induction hypothesis to $P_1$.  Since its support has fiber dimension strictly smaller than $n$, after a further elementary \etale localization there exists a finite sequence of exact sequences
\[
0\to L_i\xrightarrow{\alpha_i}N_i\to P_i\to 0,
\qquad 2\le i\le r,
\]
with the desired properties and with $N_2=P_1$.  Appending the first exact sequence to this sequence yields the total d\'evissage.

Termination follows because the fiber dimension decreases strictly at each step and therefore cannot do so indefinitely.
\end{proof}

This theorem is the technical heart of Part~I.  Everything that follows may be regarded as an extraction of consequences from this finite recursive model.

\section{Proof dossier for purity and descent through the relative assassin}

We return now to the bridge from flatness to projectivity.  The relation between purity and projective descent is one of the most elegant pieces of the paper, but it is easy to underestimate the exact role played by the relative assassin.  We therefore isolate the point where the assassin enters the proof.

\subsection{Why the assassin is the correct closed set to capture}

Suppose again that $A$ is henselian local, $B$ is of finite presentation over $A$, and $M$ is a finite $B$-module flat over $A$.  A nonzero submodule of $M$ has an associated point.  Consequently, if one can find an \etale neighborhood whose image contains every associated point of $M$, then any section of $M$ vanishing on that neighborhood must vanish globally.  This simple observation is the reason the assassin appears: it is the minimal closed set that detects every possible nonzero submodule.

In the projective descent argument, one constructs an \etale neighborhood $Y\to X$ on which the pullback of the module is projective.  To conclude that the original module coincides with the descended one, it suffices to know that the comparison map has no kernel and no cokernel.  Both are proved by showing that any nonzero kernel or cokernel would have an associated point, but such a point would have to lie in the image of $Y$, where the corresponding module already behaves projectively and therefore cannot support that defect.  Thus the assassin is exactly what makes the descent rigid.

\begin{lemma}[Detection of vanishing by the assassin]\label{lem:assassin-detects-vanishing}
Let $X$ be a noetherian scheme, $\F$ a coherent sheaf on $X$, and $U\subset X$ an open subset containing $\Ass(\F)$.  If a coherent subsheaf $\G\subset \F$ restricts to zero on $U$, then $\G=0$.
\end{lemma}

\begin{proof}
If $\G\neq 0$, then $\Ass(\G)$ is nonempty.  Since $\Ass(\G)\subset \Ass(\F)$ for a coherent subsheaf of a coherent sheaf, every associated point of $\G$ lies in $U$.  But $\G|_U=0$, so $\Ass(\G)$ is disjoint from $U$, a contradiction.  Hence $\G=0$.
\end{proof}

The lemma is elementary, yet it is responsible for much of the power of the projective descent step.  Once one remembers it, the role of the assassin becomes almost unavoidable.

\subsection{A longer descent argument}

We now give a second proof of flat plus pure implies projective that emphasizes the pushforward description.  Let $\pi\colon Y\to X$ be an \etale morphism whose image contains the relative assassin of $\widetilde M$ and such that $M_Y$ is projective over the base.  Consider the natural morphism
\[
\widetilde M\longrightarrow \pi_*\pi^*\widetilde M.
\]
Its kernel vanishes by \cref{lem:assassin-detects-vanishing}, because the morphism is visibly injective on the open image containing the assassin.  Let $\Q$ be the cokernel.  Since $\pi$ is \etale, the pullback of the sequence to $Y$ identifies the pullback of $\Q$ with the cokernel of a morphism between projective modules that is an isomorphism on the image of the assassin.  Therefore $\pi^*\Q=0$ on an open containing the assassin of $\widetilde M$.  Any associated point of $\Q$ would then lie outside that image, contradicting purity and the closure criterion.  Hence $\Q=0$.

This proof is slightly more sheaf-theoretic than the module-theoretic argument given earlier, but it illuminates an important conceptual feature: the projective object on the \etale chart is not merely compared to the original module pointwise; it is pushed forward, and the original module is recovered as the unique coherent subsheaf whose assassin is contained in the chosen image.

\section{Proof dossier for flattening by blow-up}

We now assemble a longer proof schema for the flattening theorem itself.  The goal is not to replace the original article line for line, but to make every logical dependency explicit enough that the proof can be presented over several lectures without hidden jumps.

\subsection{The closed bad locus for universal injectivity}

Fix one stage of a d\'evissage after a preliminary admissible blow-up on which the lower-dimensional quotients are already flat in the relevant range.  We are left with a morphism
\[
\alpha\colon \Lcal\to \N
\]
from a free module to a finite type module on a scheme of finite presentation over the base.  The bad locus is the subset of the base over which $\alpha$ fails to be universally injective at some point of fiber dimension exactly $n$.

Why is this bad locus closed?  Étale-locally and after shrinking on the source, one may suppose that $\N$ admits a finite presentation and that the cokernel $\Pcal$ of $\alpha$ is already flat in dimension $\ge n+1$.  By the flatness criterion, failure of universal injectivity at dimension $n$ is equivalent to failure of flatness of the middle term at dimension $n$.  But the flatness locus of a finitely presented module is open.  Therefore the non-universal-injectivity locus is closed.  The paper packages this argument through the representability theorem of Part~I; in a seminar it is useful to remember that the underlying mechanism is simply openness of flatness after the lower-dimensional quotients have been dealt with.

\subsection{Why blowing up the bad locus helps}

Blowing up a closed bad locus does not by itself make a module flat.  What it does is separate the bad locus from the good one and create an exceptional divisor on which the residual defect becomes more visible.  The module may pick up exceptional torsion after pullback, but that torsion is exactly what strict transform is designed to remove.  The improvement therefore comes from the combination
\[
\text{blow-up of the bad locus} + \text{strict transform of the pullback module}.
\]
Over the open complement of the bad locus, nothing changes.  Over the exceptional divisor, the geometry is now organized so that the failure of universal injectivity is absorbed into a quotient supported on lower fiber dimension.  This is the point at which the induction on the threshold becomes effective.

\begin{lemma}[Exceptional torsion has controlled support]\label{lem:exceptional-torsion-controlled}
Let $g\colon S'\to S$ be a $U$-admissible blow-up and let $\M$ be a finite type $\OO_X$-module which is already flat in dimension $\ge n+1$ over $U$.  Then the kernel of the natural quotient
\[
g_X^*\M\twoheadrightarrow \st_g(\M)
\]
is supported on the preimage of the center of the blow-up, hence on a subset of fiber dimension at most $n$ once the center has been chosen inside the bad locus for flatness in dimension $\ge n$.
\end{lemma}

\begin{proof}
By definition the kernel consists of sections of $g_X^*\M$ that vanish on the inverse image of the good open $U$.  Therefore its support is contained in the inverse image of the closed complement of $U$, and in particular in the exceptional locus of the blow-up.  If the center of the blow-up is chosen inside the set where flatness fails in dimension $\ge n$, then every point of the kernel projects to a point where the original module was already flat in dimension $\ge n+1$.  Hence the support of the kernel cannot contain components of fiber dimension $>n$.
\end{proof}

The lemma is one of the hidden reasons the induction closes: the strict transform never throws away information that matters above the current threshold.

\subsection{A full induction schema}

We now write the flattening theorem as a genuine induction on $n$.

\begin{theorem}[Flattening by admissible blow-up: seminar proof schema]\label{thm:flattening-schema-long}
Let $S$ be qcqs, let $U\subset S$ be quasi-compact, let $f\colon X\to S$ be of finite presentation, and let $\M$ be a finite type $\OO_X$-module whose restriction to $f^{-1}(U)$ is flat in dimension $\ge n$.  Then there exists a $U$-admissible blow-up $g\colon S'\to S$ such that $\st_g(\M)$ is flat in dimension $\ge n$.
\end{theorem}

\begin{proof}
The proof proceeds in four reductions.

\emph{First reduction: qcqs to noetherian.}  By noetherian approximation descend all data to a noetherian stage $S_\lambda$ preserving the open $U$, the morphism $f$, the module $\M$, and the condition of being flat in dimension $\ge n$ on $f^{-1}(U)$.  It is therefore enough to prove the theorem for noetherian $S$.

\emph{Second reduction: descending induction on $n$.}  For $n$ larger than the maximal fiber dimension of the support of $\M$, the statement is vacuous.  Assume the theorem known for $n+1$ and prove it for $n$.

\emph{Third reduction: total d\'evissage.}  Choose a total d\'evissage of $\M$ near the bad locus in dimension $\ge n$.  By the induction hypothesis applied to the quotients of smaller fiber dimension, after a first admissible blow-up we may suppose that their strict transforms are flat in dimension $\ge n+1$.

\emph{Fourth reduction: universal injectivity of the transition maps.}  For each transition map in the transformed d\'evissage, consider the closed locus on the base where universal injectivity fails at dimension $n$.  Blow up the union of these finitely many closed loci.  Pull back the d\'evissage and pass to strict transforms.  By construction the transition maps are now universally injective at all points of fiber dimension at least $n$, while the lower-dimensional quotients are already flat in dimension $\ge n+1$.  The flatness criterion therefore implies that the final strict transform of $\M$ is flat in dimension $\ge n$.

It remains only to check that the torsion removed by the last strict transform does not contribute to fiber dimension $\ge n$.  This is exactly \cref{lem:exceptional-torsion-controlled}.  Hence the theorem follows.
\end{proof}

The theorem as written is shorter than the original proof, but each step corresponds directly to a block of the article.  The advantage of separating the proof into reductions is that the seminar audience can see very clearly what depends on what.

\subsection{A comparison with Fitting ideals}

It is tempting to compare flattening by blow-up with the elementary fact that blowing up appropriate Fitting ideals often improves the local freeness of a module.  The comparison is instructive but incomplete.  Fitting ideals control the ranks of finite presentations and therefore detect where a module fails to be locally free of a given rank.  However, they do not by themselves encode the universal injectivity conditions extracted by the d\'evissage, nor do they distinguish the exceptional torsion created by pullback from the original defect of flatness.  The Raynaud--Gruson theorem is therefore stronger than a statement about Fitting ideals: it is a statement about repairing flatness after a birational modification while preserving the generic behavior of the module.

Still, the analogy is useful.  One may regard the closed bad loci blown up in the proof as refined Fitting-type loci attached not to a single presentation of the module but to the entire d\'evissage.  The d\'evissage acts as a sequence of presentations adapted to decreasing dimensions, and the universal injectivity failures play the role of the precise determinantal conditions one needs to eliminate.


\section{Additional affine examples and seminar exercises with detailed solutions}

To conclude the main line of the notes, we record several affine exercises that are particularly well suited for a seminar.  In the present version we give complete solutions, not only because they are useful for practice, but because each one isolates a point at which the general theory could otherwise seem more formal than it really is.  The examples are elementary, yet together they reproduce the full moral arc of the paper: universal injectivity, purity, strict transform, and the geometry hidden behind the first step of d\'evissage.

\subsection*{Exercise 1: universal injectivity on a DVR}

Let $A=k[t]_{(t)}$ and consider the exact sequence
\[
0\to A\xrightarrow{\cdot t} A\to A/(t)\to 0.
\]
Show that multiplication by $t$ is injective but not universally injective.

\emph{Solution.}  Ordinary injectivity is immediate because $A$ is a domain.  The point of the exercise is that the Raynaud--Gruson criterion never asks merely for injectivity before base change.  To test universal injectivity we tensor with the residue field $k=A/(t)$.  The map becomes
\[
A\ot_A k \xrightarrow{\cdot t} A\ot_A k,
\]
that is,
\[
k\xrightarrow{0} k.
\]
This map is no longer injective.  Hence multiplication by $t$ fails universal injectivity.

It is worth interpreting the failure geometrically.  Before passing to the special fiber the map remembers the open complement of the closed point, where $t$ is invertible.  On the special fiber the parameter $t$ vanishes, so the comparison map collapses.  In the d\'evissage this is exactly the kind of phenomenon one must exclude: a map that looks generically faithful may lose all information on the fibers that matter.  Universal injectivity is therefore not a technical strengthening of injectivity; it is the condition that prevents the top-dimensional approximation from degenerating after base change.

\subsection*{Exercise 2: purity and a visible vertical associated point}

Let $A=k[t]_{(t)}$ and set
\[
B=A[y]/(ty).
\]
Show that $B$ is not pure over $A$ by exhibiting an associated point whose closure is entirely vertical.

\emph{Solution.}  The class of $y$ in $B$ is annihilated by $t$, so the ideal $(t)$ of $B$ is the annihilator of a nonzero element.  Consequently $(t)$ is an associated prime of the $B$-module $B$.  Now look at the corresponding point of $\Spec(B)$.  Since $t$ vanishes there, the point lies entirely inside the special fiber.

To see that this violates purity, recall the geometric meaning of purity in the present context: every associated point should remain anchored to the relevant fibers in such a way that its closure meets the good generic region appropriately.  Here the point defined by $(t)$ disappears as soon as one inverts $t$; indeed on the generic fiber one has
\[
B\ot_A K \cong K[y]/(y) \cong K,
\qquad K=\operatorname{Frac}(A),
\]
so there is no trace of the $y$-direction there.  The associated point $(t)$ is therefore purely vertical.  Its closure is the closed subset $V(t)$, which never reaches the generic fiber as an associated point of the generic pullback.  This is precisely the defect purity is designed to forbid.

A complementary comparison is useful.  If one takes instead a finite free $A$-module, for example $A\oplus A$, then every associated point comes from the base ring itself, and there is no hidden vertical component created by a relation such as $ty=0$.  The contrast explains why purity is the extra geometric condition needed in the theorem flat $+$ pure $=$ projective.

\subsection*{Exercise 3: strict transform on the two standard blow-up charts}

Let $A=k[u,v]_{(u,v)}$ and let
\[
M=A[x]/(ux,vx).
\]
Blow up the origin of $\Spec(A)$.  Compute the pullback of $M$ on the two standard charts and identify the exceptional torsion killed by the strict transform.

\emph{Solution.}  Write $\widetilde S\to S$ for the blow-up of the origin.  There are two standard affine charts.

On the chart $U_u$ where $u\neq 0$ on the exceptional projective coordinate system, one may write $v=ut$.  The coordinate ring is $A_u'=k[u,t]_{(u,t)}$, and the pullback of $M$ is
\[
M_u:=A_u'[x]/(ux,utx)\cong A_u'[x]/(ux).
\]
The class of $x$ is killed by $u$, hence it is supported on the exceptional divisor $u=0$.  More precisely, the submodule generated by $x$ is isomorphic to $A_u'/(u)$.  This is exactly the exceptional torsion introduced by pullback.  The strict transform is obtained by quotienting $M_u$ by the submodule of sections supported on the exceptional divisor, so
\[
\st(M)|_{U_u}\cong M_u/(x)\cong A_u'.
\]
Thus on $U_u$ the strict transform is free of rank one over the chart.

On the second chart $U_v$ one writes $u=vs$.  The coordinate ring is $A_v'=k[v,s]_{(v,s)}$, and similarly
\[
M_v:=A_v'[x]/(vsx,vx)\cong A_v'[x]/(vx).
\]
Again the class of $x$ is annihilated by the exceptional parameter $v$, so the exceptional torsion is the copy of $A_v'/(v)$ generated by $x$, and the strict transform is
\[
\st(M)|_{U_v}\cong A_v'.
\]
The two free modules glue to the structure sheaf of the blow-up.  This example encapsulates the role of strict transform in the theorem: pullback makes the defect of flatness visible as exceptional torsion, and strict transform removes exactly that torsion while leaving the good locus untouched.

\subsection*{Exercise 4: the first d\'evissage step and the generic-rank obstruction}

Let $A$ be a DVR with uniformizer $\pi$ and let
\[
M=A[x]/(x^3,\pi x^2).
\]
The original question asks for an exact sequence
\[
0\to A\to M\to P\to 0
\]
with $P$ supported on the closed fiber.  Show first that no such sequence can exist.  Then explain what the correct first approximation should be and use it to decide whether $M$ is flat over $A$.

\emph{Solution.}  Assume such an exact sequence existed.  Tensoring with the fraction field $K=\operatorname{Frac}(A)$ would give
\[
0\to K\to K[x]/(x^3)\to P\ot_A K\to 0.
\]
If $P$ were supported on the closed fiber, then $P\ot_A K=0$.  The map $K\to K[x]/(x^3)$ would therefore have to be an isomorphism.  But
\[
\dim_K K[x]/(x^3)=3,
\]
so this is impossible.  The contradiction shows that a free rank-one module cannot capture the generic behavior of $M$.

This is exactly the lesson one should extract from the construction of d\'evissage in the paper.  The free module chosen at the first stage is not arbitrary; it must have the generic rank of the module on the top-dimensional part of the support.  Here the generic fiber has $K$-dimension $3$, with basis $1,x,x^2$.  The correct first approximation is therefore the surjective map
\[
A^{\oplus 3}\longrightarrow M,
\qquad (a,b,c)\longmapsto a+bx+cx^2.
\]
The kernel is generated by $(0,0,\pi)$ because $\pi x^2=0$ in $M$.  In particular there is $\pi$-torsion already inside $M$, visible in the element $x^2$.  Hence $M$ cannot be flat over $A$.  One may also see this directly from the local criterion, since
\[
\Tor_1^A(M,A/(\pi))\neq 0
\]
as soon as a nonzero element is annihilated by $\pi$.  The exercise is valuable because it forces one to respect the generic-rank principle built into the d\'evissage.

\subsection*{Exercise 5: why the assassin detects vanishing}

Suppose $\pi\colon Y\to X$ is an \etale morphism and its image contains the assassin of a coherent sheaf $\F$ on $X$.  Show that if $\pi^*\F=0$, then $\F=0$.

\emph{Solution.}  Assume for contradiction that $\F\neq 0$.  Then $\F$ has at least one associated point $\xi\in \Ass(\F)$.  By hypothesis the image of $Y$ contains $\xi$, so there exists a point $\eta\in Y$ mapping to $\xi$.  Because \etale pullback preserves the local shape of the module, the stalk $(\pi^*\F)_\eta$ is obtained from $\F_\xi$ by tensoring with a faithfully flat local map.  Since $\xi$ is associated for $\F$, the stalk $\F_\xi$ is nonzero.  Faithful flatness then implies $(\pi^*\F)_\eta\neq 0$, contradicting the assumption that $\pi^*\F=0$.

In the body of the theory this simple observation is indispensable.  After one passes to an \etale chart containing the relative assassin, any cokernel or kernel that vanishes on the chart must already vanish globally.  That is exactly the descent mechanism used in the proof that flat and pure modules are projective.


\section{Appendix on countable projectives, Mittag--Leffler conditions, and the role of Lazard's theorem}

Although the main body of the seminar notes concentrates on finite presentation modules, the original article begins with several lemmas about countably generated projective modules and universally injective maps.  Those lemmas are not accidental preliminaries.  They explain why projectivity interacts so well with localization and with the relative assassin.  Since these ideas are easy to lose sight of in a first reading, we isolate them here.

\subsection{Why countable projectives appear}

A projective module of finite type is a direct summand of a finite free module, but a projective module of countable type is only known a priori to be a direct summand of a countable free module.  In the local arguments of Raynaud and Gruson, one repeatedly encounters modules obtained by localizing a finite presentation object and then enlarging the ambient neighborhood to capture the relevant associated points.  It is therefore natural that countable projectives appear even when the eventual global theorem concerns finite presentation modules.

The crucial fact, going back to Lazard, is that countable projective modules may be described as filtered colimits of finite free modules with transition maps satisfying a strong compatibility condition of Mittag--Leffler type.  This structural result allows one to approximate a projective module by free modules while keeping enough control to descend universal injectivity properties.

\begin{proposition}[Countable projectives as controlled limits]\label{prop:countable-projectives-long}
Let $A$ be a ring and let $P$ be a projective $A$-module of countable type.  Then there exists a filtered inductive system of finite free $A$-modules
\[
L_1\to L_2\to \cdots
\]
with compatible maps to $P$ such that $P\cong \varinjlim L_i$ and such that the dual inverse system
\[
\Hom_A(L_1,A)\leftarrow \Hom_A(L_2,A)\leftarrow \cdots
\]
satisfies the Mittag--Leffler condition.
\end{proposition}

\begin{proof}
This is a standard form of Lazard's theorem specialized to countable type.  One starts with a countable set of generators of $P$ and writes $P$ as a direct summand of a countable free module.  By choosing finite subsets of generators and finite subsets of the retraction data, one constructs finite free modules $L_i$ mapping compatibly to $P$ and whose colimit is $P$.  The projectivity of $P$ ensures that the transition maps in the dual system eventually stabilize on images, which is exactly the Mittag--Leffler condition.
\end{proof}

The proposition matters because it replaces an abstract projective module by a very concrete infinite staircase of finite free modules.  In many local arguments the desired property is first verified on each finite free stage and then passed to the colimit.

\subsection{Universal injectivity through countable projective approximations}

Suppose now that $P$ is projective of countable type and that $u\colon P\to M$ is a morphism into a finite type module.  To decide whether $u$ is universally injective, one may test it on the finite free stages of a Lazard presentation.  If the induced maps from sufficiently large stages become injective after all residue-field base changes and if the cokernels remain flat, then the limit map is universally injective.  The idea is simple: universal injectivity is exactness after all base changes, and filtered colimits commute with tensor products.

\begin{lemma}[Testing universal injectivity on finite free stages]\label{lem:ui-on-stages}
Let $P\cong \varinjlim L_i$ be as in \cref{prop:countable-projectives-long}, and let $u\colon P\to M$ be a morphism to an $A$-module $M$.  If for every $i$ the composite $L_i\to P\xrightarrow{u} M$ is universally injective and the transition maps stabilize on images, then $u$ is universally injective.
\end{lemma}

\begin{proof}
Let $A'$ be an $A$-algebra.  Since filtered colimits commute with tensor products, one has
\[
P\otimes_A A'\cong \varinjlim (L_i\otimes_A A').
\]
If an element of $P\otimes_A A'$ maps to zero in $M\otimes_A A'$, it comes from some stage $L_i\otimes_A A'$.  Because the map from that stage is injective after tensoring with $A'$, the element vanishes already in $L_i\otimes_A A'$, hence in the colimit.  Therefore $u\otimes_A A'$ is injective.
\end{proof}

The lemma is modest, but it highlights a recurring theme of the paper: the passage from free modules to projective ones is not handled by a separate miraculous argument; it is handled by writing projectives as controlled limits of free modules and then applying the same universal-injectivity reasoning stage by stage.

\subsection{Relation with the assassin}

There is also a geometric reason countable projectives fit naturally with the assassin.  If $P$ is projective over the base, then after passing to an \etale neighborhood one may realize it as a direct summand of a free module of countable type.  The associated points of the corresponding coherent sheaf are therefore forced to behave like those of a vector bundle: they cannot appear suddenly on a hidden closed subset, because any such subset would already have to be visible in one of the finite free stages.  This is another way to see why projective modules are pure.

\section{Appendix on local proofs around the representability theorem}

The representability theorem is one of the most satisfying places in the paper because it reveals that the flattening condition is genuinely geometric and yet ultimately cut out by finitely many equations on the base.  The main body of these notes already explained the core affine argument.  Here we record a longer local proof with explicit glueing.

\subsection{Local affine charts and descent of the defining ideal}

Let $S=\Spec(A)$ be affine and suppose that $X$ admits an \etale cover by affines $Y_j=\Spec(B_j)$ such that the pullback of $\M$ to each $Y_j$ is free over $A$ and such that the images of the $Y_j$ cover the relative assassin of $\M$.  For each chart there is an ideal $I_j\subset A$ cutting out the locus on which the quotient becomes an isomorphism.  To glue these ideals globally, one must show that they agree on overlaps.

The agreement is a consequence of the assassin condition.  On the overlap $Y_{ij}:=Y_i\times_X Y_j$, both free models induce quotients whose kernels are supported away from the inverse image of the assassin.  Therefore the two kernels define the same vanishing condition after restriction to the overlap.  Since the overlap maps \etale to both charts and the quotients are of finite presentation, the corresponding ideals in $A$ coincide after localization.  The ideals $I_j$ therefore glue to a quasi-coherent ideal sheaf $\I$ on $S$, and the closed subscheme defined by $\I$ represents the functor.

\begin{proposition}[Glueing the local defining ideals]\label{prop:glueing-ideals-long}
Under the preceding hypotheses, the ideals $I_j$ obtained from the affine free charts glue to a quasi-coherent ideal sheaf on $S$ representing the bijectivity locus of the quotient.
\end{proposition}

\begin{proof}
On each affine chart $Y_j$ choose a basis of the free pullback of $\M$ and let $R_j$ be the kernel of the induced quotient.  Let $I_j$ be the ideal generated by the coordinates of a finite set of generators of $R_j$.  If $Y_{ij}$ is a nonempty overlap, then the pullbacks of $R_i$ and $R_j$ to $Y_{ij}$ coincide because both identify with the kernel of the same pulled-back quotient.  Hence the ideals they define after passing to coordinates agree after localizing on the image in $S$.  Since the cover is \etale and the modules are finitely presented, these local equalities satisfy the cocycle condition and therefore define a unique quasi-coherent ideal sheaf $\I$ on $S$.  By construction an $S$-scheme factors through the closed subscheme defined by $\I$ exactly when each of the local kernels vanishes after pullback, which is equivalent to bijectivity of the global quotient.
\end{proof}

This local glueing argument is often suppressed in the literature because the representability theorem can be quoted directly.  In a seminar, however, it is worth keeping in mind: the global closed subscheme is assembled from the same coordinate vanishing conditions that appear on each free chart.

\subsection{Iterating representability along a d\'evissage}

One of the conceptual achievements of Raynaud and Gruson is that representability is not proved once and then forgotten.  Rather, it is fed back into the d\'evissage.  The failure of universal injectivity of one transition map becomes a closed condition on the base; blowing up that condition improves the module; the new module has a shorter effective d\'evissage; representability is applied again; and so on.  Thus representability itself becomes part of an iterative birational algorithm.

This iterative viewpoint is particularly helpful when explaining to an audience why the flattening theorem is not just an existence statement.  It actually describes a procedure: detect the obstruction as a closed locus, modify the base there, remove the exceptional torsion by strict transform, and repeat at the next dimension level.


\section{Appendix of further examples and seminar exercises with full solutions}

The best way to consolidate the theory is to work through a small collection of examples in which each main notion appears in isolation.  We therefore gather here a second set of model exercises, this time with complete solutions written in a slower style.  The point is not redundancy.  The earlier exercises were chosen to illuminate the central proofs, whereas the present appendix is designed for seminar participants who want to rehearse the definitions and verify the main mechanisms by hand.

\subsection*{Exercise 1: a free approximation and the first quotient}

Let $A=k[t]_{(t)}$ and $M=A[x]/(x^2,tx)$.  Show that the class of $x$ defines an exact sequence
\[
0\to A/(t)\to M\to A\to 0
\]
of $A$-modules, and explain why this does not yet fit the d\'evissage format.  Then construct instead the map
\[
A\longrightarrow M,
\qquad 1\longmapsto 1,
\]
and compute the quotient.

\emph{Solution.}  As an $A$-module the ring $M$ is generated by the classes of $1$ and $x$, with relations $x^2=0$ and $tx=0$.  The submodule generated by $x$ is therefore isomorphic to $A/(t)$, and quotienting by this submodule kills the nilpotent direction and leaves the copy of $A$ generated by $1$.  This gives the exact sequence
\[
0\to A/(t)\to M\to A\to 0.
\]
The reason this is not the right first step of a d\'evissage is that the left term is not free.  The philosophy of the Raynaud--Gruson construction is to peel off the top-dimensional part by a free module, not by an arbitrary submodule.

Now consider the map $A\to M$ sending $1$ to the class of $1$.  The quotient is generated by the class of $x$, and the relation $tx=0$ shows immediately that
\[
M/A\cong A/(t).
\]
This quotient is supported on the closed fiber.  In this particular toy example the free rank-one approximation already captures the generic behavior because the generic fiber is
\[
M\ot_A K\cong K[x]/(x^2)\cong K\oplus Kx,
\]
but the quotient that remains after removing the copy of $A$ is forced onto the bad fiber.

\subsection*{Exercise 2: universal injectivity versus ordinary injectivity}

Let $A=k[t]_{(t)}$ and consider the inclusion $A\xrightarrow{\cdot t}A$.  Verify that it is injective but not universally injective.

\emph{Solution.}  The map is injective because $A$ is a domain.  For universal injectivity one must test the base-changed map
\[
A\ot_A A'\longrightarrow A\ot_A A'
\]
for every $A$-algebra $A'$.  Taking $A'=A/(t)$ yields the zero endomorphism of $A/(t)$, hence a noninjective map.  Therefore multiplication by $t$ is not universally injective.

The point of repeating the argument here is conceptual.  In ordinary commutative algebra one often proves a statement at the generic point and then tries to descend it to nearby fibers.  The Raynaud--Gruson criterion works in the opposite direction: it builds into the statement the requirement that injectivity survive \emph{all} base changes.

\subsection*{Exercise 3: purity detected by closures}

Let $A=k[t]_{(t)}$ and $B=A[y]/(ty)$.  Determine associated points that show directly that $B$ is not pure over $A$.

\emph{Solution.}  The element $y\in B$ is nonzero and satisfies $ty=0$.  Hence the prime ideal $(t)$ occurs as the annihilator of a nonzero element, so $(t)\in \Ass_B(B)$.  Geometrically, the corresponding point lies on the special fiber.  Over the generic fiber the relation $ty=0$ forces $y=0$ because $t$ becomes invertible, and one finds
\[
B\ot_A K\cong K.
\]
Thus the associated point $(t)$ is invisible from the generic region and remains trapped on the closed fiber.  This is exactly the pattern excluded by purity.

\subsection*{Exercise 4: strict transform on a blow-up chart}

Let $A=k[u,v]_{(u,v)}$ and $M=A[x]/(ux,vx)$.  Blow up the origin and compute the pullback to the chart $v=ut$.  Identify the exceptional torsion and the strict transform.

\emph{Solution.}  On the chart $v=ut$ the coordinate ring is $A'=k[u,t]_{(u,t)}$, and the pullback of $M$ is
\[
M' = A'[x]/(ux,utx)\cong A'[x]/(ux).
\]
The class of $x$ is annihilated by $u$.  Therefore the submodule $A'\cdot x$ is supported on the exceptional divisor $u=0$.  The strict transform is obtained by quotienting $M'$ by this exceptional torsion, so
\[
\st(M')\cong A'.
\]

\subsection*{Exercise 5: projective descent through the assassin}

Suppose $\pi\colon Y\to X$ is \etale and its image contains the assassin of a coherent sheaf $\F$ on $X$.  Show that if $\pi^*\F$ is zero, then $\F$ is zero.

\emph{Solution.}  If $\F$ were nonzero, choose an associated point $\xi\in \Ass(\F)$.  Since the image of $Y$ contains the assassin, there exists $\eta\in Y$ mapping to $\xi$.  The stalk map
\[
\OO_{X,\xi}\longrightarrow \OO_{Y,\eta}
\]
is faithfully flat because \etale morphisms are flat and local.  Therefore
\[
\F_\xi\ot_{\OO_{X,\xi}} \OO_{Y,\eta}
\]
vanishes if and only if $\F_\xi$ vanishes.  But $\F_\xi\neq 0$ because $\xi$ is associated.  Hence $(\pi^*\F)_\eta\neq 0$, contradicting the assumption $\pi^*\F=0$.  So $\F=0$.


\section{Lecture-by-lecture seminar roadmap}

Because these notes are intended for a seminar exposition, it is useful to record a possible lecture sequence.  The point of the roadmap is not merely pedagogical.  By distributing the material over several sessions one sees more clearly which arguments are foundational and which are consequences.

\subsection*{Lecture 1: the problem of flatness and why the local criterion is not enough}

The opening lecture should begin with the ordinary local criterion of flatness and with a discussion of why it leaves several questions unanswered.  One should emphasize three deficiencies.  First, the criterion is not visibly geometric: the vanishing of a Tor group does not by itself describe the shape of the bad locus.  Second, it gives no reason for flatness to be connected to projectivity.  Third, it offers no strategy for repairing nonflatness.  The Raynaud--Gruson paper is then introduced as a response to these deficiencies.  The audience should leave the first lecture with the expectation that the paper will transform a homological criterion into a geometric and birational theory.

\subsection*{Lecture 2: local factorization and the geometry of the support}

The second lecture should focus almost entirely on the local finite factorization theorem.  One should not rush.  This is where the geometric universe of the paper is fixed.  The central point is that after an elementary \etale localization, the support of the module may be studied through a finite morphism to a smooth base with geometrically integral fibers.  It is helpful to draw the picture of the support mapping finitely to a smooth relative affine space, with the chosen point isolated on the closed fiber.  The conclusion of the lecture should be that one now has a coordinate system in which generic rank on the support makes sense.

\subsection*{Lecture 3: the first step of d\'evissage}

The third lecture constructs the first exact sequence
\[
0\to \Lcal_1\to \N_1\to \Pcal_1\to 0.
\]
The pedagogical goal is to explain why the free module is chosen by generic rank considerations and why the quotient has smaller fiber dimension.  The audience should see that the free module is not chosen to make the algebra pleasant, but to kill the top-dimensional generic behavior of the support.  A good exercise during the lecture is to revisit the affine example $A[x]/(x^2,tx)$ and to identify the first d\'evissage step explicitly.

\subsection*{Lecture 4: total d\'evissage and termination}

This lecture iterates the previous step and proves the existence of a total d\'evissage.  The emphasis should be on the dimension drop and on the fact that repeated \etale localization is harmless because all relevant statements are local on both source and base.  One should stress that the process terminates because fiber dimensions are integers that decrease strictly.  By the end of the lecture the audience should regard the module as encoded by a finite recursive exact sequence rather than as an opaque coherent sheaf.

\subsection*{Lecture 5: universal injectivity and the flatness criterion}

Only now should one turn fully to the homological algebra.  The lecturer should define universal injectivity and compare it with ordinary injectivity and split injectivity.  The Tor exact sequence should be written out carefully.  Then the flatness criterion is proved by induction along the d\'evissage.  It is worth pausing after the proof to observe that the theorem is a perfect blend of geometry and algebra: geometry provided the recursion, algebra verified that the recursion detects flatness.

\subsection*{Lecture 6: purity and the relative assassin}

The sixth lecture introduces the relative assassin and purity.  It is important here to resist the temptation to present purity as a mysterious additional hypothesis.  One should instead explain that associated points detect hidden torsion and embedded components, so controlling their closures is exactly what one needs to prevent projectivity from failing for geometric reasons not seen by flatness alone.  Good board work includes drawing closures of associated points relative to the special fiber.

\subsection*{Lecture 7: flat plus pure implies projective}

This lecture proves the projectivity theorem.  The ideal structure of the lecture is: first explain why purity lets one choose an \etale neighborhood containing the assassin; next explain why the pullback becomes projective there; then prove that the original module injects universally into the pushed-forward chart module; finally show that the cokernel vanishes because any nonzero cokernel would create an associated point outside the chosen image.  The lecture is conceptually dense but extremely rewarding, because it shows how a geometric condition on associated points becomes a descent mechanism.

\subsection*{Lecture 8: representability of the flattening condition}

Once the audience understands purity and projectivity, the representability theorem becomes surprisingly natural.  On a free chart the quotient becomes an isomorphism exactly when the coordinates of its kernel vanish, so the condition is cut out by a finite ideal on the base.  The lecture should work this out in full affine detail.  One should then explain how the local ideals glue because the charts agree on overlaps containing the assassin.  This lecture is a good place to emphasize that the future blow-up theorem will use representability iteratively.

\subsection*{Lecture 9: strict transform and the projective prototype}

The ninth lecture should pivot to birational geometry.  Define the strict transform slowly and prove its maximality property.  Then discuss the projective noetherian prototype via Quot schemes.  The message of the lecture is that flattening is already visible as a compactification problem for the family of generic quotients.  The strict transform is the quotient that keeps the generic quotient intact while discarding the exceptional torsion created by pullback.

\subsection*{Lecture 10: induction on the threshold and the flattening theorem}

The final lecture assembles everything.  First reduce to the noetherian case by approximation.  Then assume by descending induction that lower-dimensional quotients are already flattened.  Blow up the residual bad loci where universal injectivity fails.  Pass to the strict transform.  Apply the flatness criterion.  Explain carefully why the torsion removed by strict transform lies below the threshold.  The lecture should end by returning to the opening question: the local criterion of flatness has now become a global geometric and birational theorem.

This roadmap is of course only one possible distribution, but it reflects the true logic of the paper and usually leads to a much more coherent seminar experience than trying to cover the material in the printed order without thematic regrouping.

\section{Comparative remarks with EGA, generic flatness, and later flattening theories}

It is worth closing the mathematical narrative with a few comparisons.  The Raynaud--Gruson paper sits historically after the foundational work of EGA on flatness, generic flatness, and semicontinuity, but it is not merely a corollary or amplification of those results.  Its novelty lies in converting generic flatness into a birational repair theorem and in identifying purity as the missing bridge from flatness to projectivity.

Generic flatness in EGA says that for a finitely presented module over a noetherian integral base, there exists a dense open over which the module is flat.  This is a magnificent theorem, but it does not address what happens on the complement.  Raynaud and Gruson ask a more ambitious question: can one modify the base so that the module becomes flat everywhere after replacing pullback by strict transform?  The answer is yes.  Thus the paper may be read as the theory that completes generic flatness by adding a birational extension principle.

The relation with EGA is also visible at the level of local arguments.  Many ingredients---openness of flatness, dimension semicontinuity, noetherian approximation, and behavior of finite presentation modules under base change---are pure EGA.  What the paper contributes is a new synthesis.  The d\'evissage organizes the local algebra, purity reorganizes the geometry of associated points, and representability reorganizes the flattening condition.  The theorem of flattening by blow-up is what becomes visible once these reorganizations are in place.

Later flattening theories, for example flattening stratifications of coherent sheaves on projective morphisms, often appear more categorical or moduli-theoretic.  Yet their core intuition is already present here.  A flattening stratification divides the base into pieces on which the sheaf has constant Hilbert polynomial or constant flatness behavior.  Raynaud and Gruson instead produce a birational modification that repairs the bad strata rather than merely classifying them.  In this sense the paper is simultaneously stratificational and birational.

Another useful comparison concerns determinantal loci and Fitting ideals.  Fitting ideals detect where a module fails to have the expected rank or to be locally free, and blowing them up can improve presentations.  But the Raynaud--Gruson algorithm is finer.  It tracks not just the failure of expected rank, but the failure of universal injectivity along a whole d\'evissage.  The successive bad loci are therefore refined Fitting-type conditions attached to the geometry of the support at each dimension level.

Finally, one should mention that the language of algebraic spaces at the end of the paper is not an afterthought.  It signals that the authors understood flattening as part of a broader geometric technology applicable beyond schemes.  From a modern viewpoint this is entirely natural: if flattening is a statement about representability and birational modification, one expects it to survive in any setting where finite presentation, flatness, and blow-ups make sense.

\section{Solved model problems for a reading course}

The following solved model problems are designed for a reading course based on these notes.  They do not attempt the full difficulty of the main theorems, but they force the reader to practice each of the decisive ideas in a manageable setting.

\subsection*{Problem 1: detect the first quotient in a d\'evissage}

Let $A=k[t]_{(t)}$ and $M=A[x]/(x^3,tx^2)$.  Show that the map
\[
A\longrightarrow M,
\qquad 1\mapsto 1,
\]
fits into an exact sequence
\[
0\to A\to M\to P\to 0
\]
where $P$ is supported on the closed fiber.  Determine $P$ explicitly.

\emph{Solution.}  The image of $A$ is the submodule generated by the class of $1$.  Modding out by this submodule leaves the classes of $x$ and $x^2$, subject to the relations $tx^2=0$ and $x^3=0$.  Since $x$ and $x^2$ both vanish generically on the fiber over the generic point of $A$, the quotient is supported on the closed fiber.  More concretely,
\[
P\cong (x,x^2)/(x^3,tx^2)\cong k\cdot \bar x\oplus k\cdot \overline{x^2}
\]
as an $A$-module, with all elements killed by $t$.  This is precisely the lower-dimensional defect measured by the first d\'evissage step.

\subsection*{Problem 2: show that a map is not universally injective}

Let $A=k[t]_{(t)}$ and consider the exact sequence
\[
0\to A\xrightarrow{\cdot t}A\to A/(t)\to 0.
\]
Explain why the first map fails universal injectivity and why this matters for the flatness criterion.

\emph{Solution.}  Tensoring the sequence with $A/(t)$ yields
\[
0\to A/(t)\xrightarrow{0}A/(t)\to A/(t)\to 0,
\]
so the first map is no longer injective.  Therefore multiplication by $t$ is not universally injective.  In the flatness criterion, a failure of universal injectivity means that even if a module looks injectively presented before base change, the quotient can cease to behave flatly on a special fiber.  This is exactly the kind of defect that the d\'evissage is designed to isolate.

\subsection*{Problem 3: verify purity for a projective module}

Let $A$ be local and let $P$ be a finite projective $A$-module.  Show directly that $P$ is pure over $A$.

\emph{Solution.}  Since $P$ is finite projective, after choosing a complementary summand it is a direct summand of a finite free module $A^{\oplus r}$.  The associated points of $A^{\oplus r}$ are exactly those of $A$, and every associated point specializes to the closed point in the local case.  A direct summand cannot create a new associated point outside the closure behavior already present in the free module.  Hence the closures of all associated points of $P$ meet the closed fiber, which is the purity condition.

\subsection*{Problem 4: compute a strict transform on one chart}

Let $A=k[u,v]_{(u,v)}$ and $M=A[x]/(ux,vx)$.  Blow up the origin.  On the chart $v=ut$, compute the pullback and strict transform.

\emph{Solution.}  On the chart one has
\[
A[u,t,x]/(ux,utx)\cong A[u,t,x]/(ux).
\]
The class of $x$ is annihilated by $u$, hence it is exceptional torsion supported on the exceptional divisor $u=0$.  The strict transform quotients out by that torsion and is therefore isomorphic to the quotient by the class of $x$, which is just the structure sheaf of the chart.  In particular the strict transform is flat.

\subsection*{Problem 5: understand the role of the assassin in descent}

Let $\pi\colon Y\to X$ be \etale and let $\F$ be a coherent sheaf on $X$.  Suppose the image of $Y$ contains $\Ass(\F)$.  Show that if a subsheaf $\G\subset \F$ restricts to zero on $Y$, then $\G=0$.

\emph{Solution.}  If $\G$ were nonzero, it would have an associated point $\xi\in\Ass(\G)\subset \Ass(\F)$.  Since $\xi$ lies in the image of $Y$, the restriction of $\G$ to $Y$ could not be zero at the corresponding point.  This contradiction shows that $\G=0$.  The argument is the local form of the projective descent step in the theorem flat $+$ pure implies projective.

\subsection*{Problem 6: relate representability to coordinates of the kernel}

Let $A$ be a ring, let $M=A^{\oplus 2}$, and let $u\colon M\to N$ be the quotient whose kernel is generated by the vector $(a,b)\in A^{\oplus 2}$.  Determine the universal isomorphism locus of $u$.

\emph{Solution.}  The quotient becomes an isomorphism after base change $A\to A'$ if and only if the kernel vanishes after tensoring with $A'$, which happens exactly when both $a$ and $b$ vanish in $A'$.  Hence the locus is represented by the closed subscheme defined by the ideal $(a,b)$.  This is the simplest affine model of the representability theorem.

These solved problems are well suited for discussion sections or for written assignments in a course built around the present notes.





\section{Detailed solutions to the seminar problem set}

The solved model problems written earlier were intentionally concise so that they could be used quickly in a reading course.  Since several seminar participants asked for fuller written solutions, we now return to the same cluster of ideas in a slower format.  The goal of the present section is twofold.  First, it provides genuinely worked-out computations that can be used as homework solutions.  Second, it extends the narrative of the preceding chapters by showing how the definitions behave in concrete modules one can compute by hand.

\subsection*{Problem A: a miniature local criterion computation}

Let $A=k[t]_{(t)}$ and $M=A\oplus A/(t)$.  Show directly from the local criterion that $M$ is not flat over $A$.

\emph{Solution.}  Tensor the short exact sequence
\[
0\to A\xrightarrow{\cdot t} A\to A/(t)\to 0
\]
with $M$.  Since tensor products distribute over direct sums, one gets
\[
\Tor_1^A(M,A/(t))\cong \Tor_1^A(A/(t),A/(t)).
\]
Now the standard two-term free resolution of $A/(t)$ shows that
\[
\Tor_1^A(A/(t),A/(t))\cong A/(t)\neq 0.
\]
Hence $M$ is not flat over $A$.

What matters for the philosophy of the paper is that the nonflat part is transparent: it is exactly the torsion summand $A/(t)$.  In the global theorem the bad locus is rarely handed to us so explicitly; the d\'evissage and the representability theorem are the tools that locate this torsion geometrically.

\subsection*{Problem B: a toy relative d\'evissage over a smooth base}

Let $A$ be a ring and set $T=\Spec(A[z])$.  Let
\[
N=A[z]/(z^2)
\]
viewed as an $A[z]$-module.  Construct a simple d\'evissage of $N$ over $A$ and explain why the quotient drops in relative dimension.

\emph{Solution.}  The support of $N$ is the Cartier divisor $z=0$ with a doubled structure.  The generic rank of $N$ along that support is one, so the first free approximation is the map
\[
A[z]\longrightarrow N,
\qquad 1\longmapsto 1.
\]
This map is surjective, because every element of $N$ is represented by $a+bz$.  The kernel is the principal ideal $(z^2)$, and one may write the presentation as
\[
0\to A[z]\xrightarrow{\cdot z^2} A[z]\to N\to 0.
\]
The quotient $N$ is supported on the divisor $z=0$, whose fibers over $A$ have relative dimension $0$, whereas the ambient smooth base $T\to \Spec(A)$ has relative dimension $1$.  Thus the quotient has indeed moved one step down in geometric complexity.

This example is deliberately simple, but it shows the governing principle: the free module is chosen according to generic rank on the support, and the quotient records the failure of that approximation away from the generic points.

\subsection*{Problem C: purity for a direct summand of a free module}

Let $A$ be a local ring and let $P$ be a finite projective $A$-module.  Give a careful proof that $P$ is pure over $A$.

\emph{Solution.}  Since $P$ is finite projective, there exists a finite module $Q$ such that
\[
P\oplus Q\cong A^{\oplus r}
\]
for some integer $r$.  Associated points behave monotonically under direct summands: every associated point of $P$ is an associated point of $P\oplus Q$, hence of $A^{\oplus r}$.  But
\[
\Ass(A^{\oplus r})=\Ass(A).
\]
Because $A$ is local, every associated prime of $A$ specializes to the closed point, and the closure of the corresponding point in $\Spec(A)$ contains that closed point.  Therefore the closure of every associated point of $P$ also meets the closed fiber.

Geometrically, one may say that a direct summand of a vector bundle cannot generate a new hidden vertical component.  Any associated point of $P$ must already have been visible in the ambient free module.  This is the simplest local picture of why projective modules are pure and why purity is the right geometric weakening of projectivity for descent arguments.

\subsection*{Problem D: the universal flattening locus in coordinates}

Let $A$ be a ring and let
\[
M=A^{\oplus 3},
\qquad N=M/\langle (a,b,c)\rangle,
\]
where $(a,b,c)\in A^{\oplus 3}$.  Determine the closed subscheme of $\Spec(A)$ on which the quotient map $M\to N$ becomes an isomorphism after base change.

\emph{Solution.}  After base change $A\to A'$ the map becomes an isomorphism precisely when the kernel generated by $(a,b,c)$ vanishes in $A'^{\oplus 3}$.  Since a cyclic submodule of a free module is zero if and only if its generator is zero, this happens if and only if the images of $a$, $b$, and $c$ all vanish in $A'$.  Thus the universal isomorphism locus is represented by the closed subscheme
\[
V(a,b,c)\subset \Spec(A).
\]

This is exactly the affine mechanism hidden inside the representability theorem.  On an \etale chart where the relevant sheaf becomes free, the flattening condition is no more mysterious than the vanishing of finitely many coordinate functions.

\subsection*{Problem E: strict transform and compatibility with a second blow-up}

Continue with the module
\[
M=A[x]/(ux,vx)
\]
on $A=k[u,v]_{(u,v)}$, and let $S_1\to S$ be the blow-up of the origin.  Suppose one blows up once again along a center contained in the exceptional divisor of $S_1$.  Explain why the strict transform after the second blow-up agrees with the strict transform of the first strict transform.

\emph{Solution.}  Let $g_1\colon S_1\to S$ and $g_2\colon S_2\to S_1$ be the two blow-ups.  The pullback $g_1^*M$ contains a maximal subsheaf supported on the exceptional divisor of $g_1$; quotienting by it gives $\st_{g_1}(M)$.  Pulling this quotient back along $g_2$ introduces new torsion only on the new exceptional divisor of $g_2$ and on the inverse image of the old exceptional divisor.  Quotienting again by sections supported there gives $\st_{g_2}(\st_{g_1}(M))$.

If one instead looks directly at the composite $g_1\circ g_2$, the pullback $(g_1\circ g_2)^*M$ has a maximal subsheaf supported on the total exceptional divisor of the composite.  But this maximal subsheaf is removed in two stages by the previous construction.  Therefore the two strict transforms coincide:
\[
\st_{g_2}(\st_{g_1}(M))\cong \st_{g_1\circ g_2}(M).
\]

This compatibility is one of the silent structural facts behind the iteration in the flattening theorem.  Without it one could improve the module at one stage only to lose control of the previous corrections at the next stage.

\subsection*{Problem F: a compact reconstruction of the final induction}

Explain in a toy model why an induction on the threshold ``flat in dimension $\ge n$'' is natural.  Use a module that is already flat away from a finite set of closed points and describe what the next blow-up is meant to do.

\emph{Solution.}  Suppose $f\colon X\to S$ is of finite presentation and $\M$ is a coherent sheaf such that $\M$ is flat on the complement of finitely many closed points of $X$.  Then $\M$ is flat in dimension $\ge 1$: the bad locus has fiber dimension $0$.  The next blow-up is not trying to alter the open set where flatness is already known.  Its purpose is to isolate the residual closed points on the base over which universal injectivity fails in the d\'evissage.  After pullback to the blow-up those failures appear as torsion concentrated on the exceptional divisors over the chosen centers.  The strict transform removes exactly that torsion, leaving a sheaf whose remaining defect has lower-dimensional support than before.

The relevance of the threshold formulation is now clear.  One does not try to solve the whole problem at once.  One first makes the sheaf flat outside lower-dimensional bad sets, and then one improves those lower-dimensional bad sets one threshold at a time.  The induction mirrors the actual geometry of the support.

\subsection*{Problem G: a reading-course proof that ordinary presentations are not enough}

Give an example of a finite presentation by free modules that does \emph{not} behave like a d\'evissage, and explain why the quotient fails to lower geometric complexity.

\emph{Solution.}  Take again
\[
M=A[x]/(x^3,\pi x^2)
\]
with $A$ a DVR.  The standard presentation by generators and relations is certainly a finite presentation by free modules, but it does not by itself single out the generic rank or separate the top-dimensional behavior from the lower-dimensional defect.  If one naively uses the injection $A\to M$ given by $1\mapsto 1$, the quotient still has generic rank two.  Hence no dimension drop occurs.

This is exactly why the d\'evissage is more refined than a finite presentation.  It is a presentation chosen with geometric intent: the free term must agree with the module at the generic points of the highest-dimensional part of the support.  Without that property there is no reason for the quotient to become simpler.

\subsection*{Problem H: purity versus torsion-freeness}

Produce a short explanation of why torsion-freeness over the base is weaker than purity in the relative sense.

\emph{Solution.}  Torsion-freeness says only that no nonzero element is annihilated by a nonzerodivisor from the base ring.  It is therefore an internal condition on the module as a base module.  Purity is subtler: it is a statement about closures of associated points in the total space relative to the fibers.  A module may be torsion-free over the base and still possess associated points living on components that do not spread out correctly from the generic fiber.  Such associated points do not contradict torsion-freeness, but they do contradict purity.

This distinction is one of the conceptual achievements of Raynaud and Gruson.  Projectivity fails for reasons that are often invisible to torsion calculations alone; one must inspect the geometry of associated points.



\section{Supplementary seminar dossier: four more worked computations}

The theory of the paper is dense enough that a seminar often benefits from one further layer of short worked computations.  The following four complements are intentionally modest.  Each one isolates a statement that appears repeatedly in the body of the notes and rewrites it as a direct calculation.  They are useful as board exercises, but they also serve as a bridge between the abstract formulations and the local algebra one can actually manipulate.

\subsection*{Supplement 1: why flat quotients preserve injectivity in a short exact sequence}

Let
\[
0\to L\xrightarrow{\alpha} N\to P\to 0
\]
be an exact sequence of $A$-modules with $L$ flat and $P$ flat.  Explain why $N$ is flat.

\emph{Solution.}  Apply $-\ot_A A'$ to the sequence for an arbitrary $A$-module $A'$.  Since $P$ is flat, the functor $-\ot_A A'$ is exact at the right end, and since $L$ is flat there is no $\Tor_1^A(L,A')$ term on the left.  Thus the base-changed sequence
\[
0\to L\ot_A A'\to N\ot_A A'\to P\ot_A A'\to 0
\]
remains exact for every $A'$.  This is one of the standard characterizations of flatness, so $N$ is flat.

The relevance to Raynaud--Gruson is immediate.  Once the d\'evissage provides a short exact sequence whose outer terms are already known to be flat and whose left map is universally injective, the middle term becomes flat.  The paper uses this observation recursively.

\subsection*{Supplement 2: a direct computation of an associated point}

Let $A=k[t]_{(t)}$ and $M=A[x]/(tx,x^2)$.  Compute an associated prime of $M$ and explain why $M$ is not flat.

\emph{Solution.}  The class of $x$ is nonzero and is annihilated by $t$, so $(t)$ is the annihilator of a nonzero element.  Hence $(t)$ is an associated prime of $M$ as an $A$-module.  In particular $M$ contains $t$-torsion.  Since flat modules over a DVR are torsion-free, $M$ cannot be flat.

What is pedagogically useful here is that the same element $x$ simultaneously signals two defects: homologically it gives torsion and therefore nonflatness, while geometrically it produces a vertical associated point and therefore threatens purity as well.  This is the kind of double interpretation the paper exploits systematically.

\subsection*{Supplement 3: checking that the bad locus is closed in a toy model}

Let $A=k[a,b]$ and let
\[
u\colon A^{\oplus 2}\longrightarrow A,
\qquad (u,v)\longmapsto au+bv.
\]
Describe the locus on which $\nu$ becomes surjective after base change and explain why its complement is closed.

\emph{Solution.}  After base change to an $A$-algebra $A'$, the map is surjective if and only if the ideal generated by the images of $a$ and $b$ is the whole of $A'$.  Thus the failure of surjectivity is exactly the condition that both $a$ and $b$ vanish, that is, the closed locus $V(a,b)\subset \Spec(A)$.  Equivalently, the surjectivity locus is the open complement $D(a)\cup D(b)$.

Although this is far simpler than the representability theorem in the notes, the mechanism is the same: a property of a morphism of free modules becomes a statement about the vanishing of finitely many coordinates.  In the actual theorem those coordinates arise from the kernel of a quotient on an \etale free chart.

\subsection*{Supplement 4: reading the final blow-up theorem backwards}

Give a short explanation of how one should read the flattening theorem backwards, starting from the strict transform and recovering the original obstruction.

\emph{Solution.}  Suppose a blow-up $g\colon S'\to S$ has been chosen and the strict transform $\st(\M)$ is flat on $X\times_S S'$.  Reading the theorem backwards, one identifies two pieces of information.  First, the pullback $g^*\M$ itself need not be flat; the difference between $g^*\M$ and $\st(\M)$ is precisely the exceptional torsion supported over the center of the blow-up.  Second, that torsion is not arbitrary.  It is the geometric shadow of the universal-injectivity failures detected in the d\'evissage before modification.

In other words, the blow-up does not create a new theorem by magic.  It relocates the old obstruction from the interior of the family to the exceptional divisor, where strict transform can remove it cleanly.  This backwards reading is often the most convincing way to explain to an audience why strict transform, and not raw pullback, is the correct operation in the flattening theorem.

\subsection*{Supplement 5: why the quotient in a d\'evissage must be measured fiberwise}

A frequent question in seminar is the following: why does Raynaud--Gruson insist on the drop of \emph{fiber dimension} rather than only on a drop of support dimension inside the total space?  The answer is that the entire induction is carried out relative to the base.  Let us make this explicit in a toy model.

Take $A=k[t]_{(t)}$ and consider on $X=\Spec(A[z])$ the module
\[
M=A[z]/(tz,z^2).
\]
As a module on the total space, the support of $M$ is the closed subset $z=0$, which has codimension one in $X$.  But the decisive observation is fiberwise: over the generic point of $\Spec(A)$ the relation $tz=0$ forces $z=0$, so the generic fiber is just $K$ and has support of dimension $0$; over the closed fiber the relation becomes $z^2=0$, so one still has a doubled point, again of dimension $0$.  Thus the interesting complexity is not the absolute dimension of the support but the way the module sits relative to the base and how much of it survives in each fiber.

Now compare this with a free approximation by the map $A[z]\to M$, $1\mapsto 1$.  The quotient is generated by the class of $z$, and because $tz=0$ it is supported entirely over the closed point of the base.  The geometric simplification is therefore best described by saying that the quotient has smaller \emph{relative} complexity: it no longer contributes on the generic fiber at all.  This is the pattern the d\'evissage formalizes in general.  If one tracked only absolute support dimensions in the total space, one would miss exactly the birational information needed for flattening.

Seen this way, the dimension threshold ``flat in dimension $\ge n$'' is not an artificial bookkeeping device.  It is a precise language for saying how much of the obstruction remains visible along fibers of the morphism under study.  That is why the theorem is stated and proved relative to the base at every stage.



\section{Dependency map, common misunderstandings, and final synthesis}

The present notes are long enough that it is easy for the reader to lose the thread of dependence between the main statements.  We therefore end with a dependency map together with a list of common misunderstandings that often arise in a first reading.  This final section is intentionally discursive.  Its purpose is to stabilize the narrative of the paper once the technical work has been absorbed.

\subsection*{Dependency map of the main theorems}

The first dependency is geometric.  The local finite factorization theorem produces a finite morphism to a smooth base with geometrically integral fibers.  Without this step there is no meaningful generic-rank analysis of the support and therefore no relative d\'evissage.

The second dependency is combinatorial in nature.  Once the support sits finitely over a smooth base, one can choose free modules that agree with the given module at the generic points of top-dimensional components.  This is what creates the d\'evissage.  The strict drop in fiber dimension of the successive quotients is the finite combinatorial invariant that guarantees termination.

The third dependency is homological.  The d\'evissage by itself is only an exact sequence machine.  To turn it into a flatness detector one needs the local algebra of universal injectivity and the behavior of flatness in short exact sequences.  Here the Tor exact sequence plays the central role.  The flatness criterion is therefore impossible without both the d\'evissage and the universal-injectivity lemmas.

The fourth dependency is geometric again.  Flatness is not enough to guarantee projectivity because hidden associated components may survive.  Purity enters at exactly this point.  It is the control of associated points that makes projective descent possible.  Thus the theorem flat $+$ pure implies projective depends simultaneously on the flatness criterion and on the geometric analysis of the relative assassin.

The fifth dependency is functorial.  The representability of the universal flattening condition depends on purity because purity provides the \etale free charts containing the assassin.  Once on those charts, the functor is cut out by explicit equations on the base.  Hence representability depends on the projectivity theorem at the chart level and on the affine coordinate argument globally.

The sixth dependency is birational.  The flattening theorem by blow-up depends on representability because one needs to realize the failure of universal injectivity as a closed locus on the base.  It depends on strict transform because pullback to the blow-up creates exceptional torsion that must be removed without altering the good open.  And it depends on the d\'evissage because the blow-up is applied recursively to the transition maps and lower-dimensional quotients.  Thus the final theorem is genuinely a synthesis of every preceding chapter.

A compact way to remember the logic is to keep the main steps in a single vertical chain, so that the eye sees the cumulative nature of the argument rather than reading it as a single compressed slogan:
\[
\renewcommand{\arraystretch}{1.5}
\begin{array}{c}
\text{local factorization near the chosen point}\\
\Downarrow\\
\text{relative d\'evissage of the module}\\
\Downarrow\\
\text{flatness criterion via universal injectivity and the last quotient}\\
\Downarrow\\
\text{purity and the projectivity theorem}\\
\Downarrow\\
\text{representability of the universal flattening condition}\\
\Downarrow\\
\text{flattening by admissible blow-up and strict transform}
\end{array}
\]
What makes the paper remarkable is that none of these arrows is formal.  Each one introduces a new idea.

\subsection*{Common misunderstanding 1: the d\'evissage is just a finite presentation}

A first misunderstanding is to think that the d\'evissage is nothing more than a finite presentation of the module by free modules.  This is not correct.  Every finite type module has a presentation by finite free modules, but an arbitrary presentation does not lower the geometric complexity of the support.  The d\'evissage is special because at each stage the chosen free module coincides with the given module at the generic points of the top-dimensional components.  Therefore the quotient is forced into smaller fiber dimension.  The d\'evissage is a dimension-sensitive presentation, not merely a presentation.

\subsection*{Common misunderstanding 2: universal injectivity is a technical luxury}

A second misunderstanding is to regard universal injectivity as a technical luxury added for safety.  In reality it is the exact condition needed to make the quotient sequence behave after every base change.  Ordinary injectivity can fail dramatically on special fibers.  The module $A\xrightarrow{\cdot t}A$ over a DVR is the simplest witness.  Since flattening is by definition a statement stable under arbitrary base changes of the modified base, universal injectivity is unavoidable.  It is the homological form of geometric permanence.

\subsection*{Common misunderstanding 3: purity means irreducibility or torsion-freeness}

A third misunderstanding is to confuse purity with irreducibility of the support or with torsion-freeness.  A pure module may live on a reducible support, and a torsion-free module need not be pure relative to a base.  Purity is a statement about the closure of associated points relative to the fibers.  Its content is not that the module has no torsion in an absolute sense, but that every associated point remains anchored to the relevant fibers and therefore cannot hide in a vertical component invisible from the generic region.

\subsection*{Common misunderstanding 4: the blow-up itself flattens the pullback module}

A fourth misunderstanding is to think that after the correct admissible blow-up the pulled-back module $g_X^*\M$ becomes flat.  This is usually false.  The pullback typically acquires extra torsion supported on the exceptional divisor.  The theorem concerns the \emph{strict transform}, not the raw pullback.  This distinction is absolutely essential.  In fact one can say that the blow-up isolates the bad locus while the strict transform discards the new defect created by the isolation process.

\subsection*{Common misunderstanding 5: the flattening theorem is purely projective in nature}

Because the projective noetherian prototype uses Quot schemes, one may mistakenly think that the general theorem is fundamentally a projective theorem.  The truth is subtler.  Quot schemes provide the clearest model of what should happen, but the general theorem is driven by finite presentation, local d\'evissage, representability of closed conditions, and noetherian approximation.  Projectivity of the source morphism is convenient for the prototype, not essential for the actual theorem.

\subsection*{A final conceptual synthesis}

We may now reread the whole paper in one conceptual paragraph.  A finite module over a finitely presented morphism is first straightened locally so that its support becomes finite over a smooth base of controlled relative dimension.  This allows one to peel off the top-dimensional behavior by free approximations, producing a finite d\'evissage.  The homological content of this d\'evissage is expressed by universal injectivity, which together with flatness of the terminal quotient characterizes flatness of the original module.  To pass from flatness to projectivity, one inspects the associated points of the module and imposes the purity condition that keeps them geometrically attached to the fibers.  Once the module is projective on appropriate \etale charts containing the assassin, the base changes on which a quotient becomes flat or bijective are represented by closed subschemes cut out by explicit equations.  Those closed loci are then used as centers of admissible blow-ups.  Pullback to the blow-up isolates the bad behavior but also creates exceptional torsion; the strict transform removes exactly that torsion while preserving the good open.  Repeating the process by descending dimension thresholds yields a birational repair of flatness.

This paragraph is longer than a slogan but shorter than the proof, and perhaps that is the right scale at which to remember the paper.  One should not remember only the theorem of flattening by blow-up.  One should remember the architecture that makes the theorem inevitable.

\subsection*{A short glossary of recurring phrases}

\emph{Elementary \etale neighborhood.}  A local replacement used to isolate the correct branch of the support while preserving the relevant fiberwise geometry.

\emph{Relative d\'evissage.}  A finite recursive exact sequence in which free modules capture the top-dimensional generic behavior and the quotients drop in fiber dimension.

\emph{Universal injectivity.}  Injectivity preserved after every base change; the correct homological condition for the quotient in a d\'evissage to inherit flatness.

\emph{Relative assassin.}  The set of associated points relevant to the sheaf over the fibers; the minimal closed set whose capture by an \etale chart makes projective descent rigid.

\emph{Purity.}  The condition that closures of associated points meet the appropriate fibers; geometrically, the prohibition of hidden vertical associated components.

\emph{Strict transform of a module.}  The maximal quotient of the pullback to a blow-up that agrees with the original module over the good open and discards the exceptional torsion created by pullback.

\emph{Flat in dimension $\ge n$.}  Flat outside a subset of fiber dimension $<n$; the correct inductive threshold for the flattening theorem.

\subsection*{Closing remark}

A great many foundational papers remain important because later work depends on their statements.  The Raynaud--Gruson article is important for a stronger reason: its internal architecture is itself a lesson in mathematical method.  It shows how to begin with a local homological criterion, reinterpret it geometrically through supports and associated points, package it functorially as a representable condition, and finally resolve it birationally.  This movement across registers---algebraic, geometric, functorial, birational---is what gives the paper its enduring depth, and it is what these notes have tried to make visible in narrative form.




\section{Bibliographical guide and suggestions for further study}

A final bibliographical guide may be helpful for readers who want to continue beyond the present seminar notes.  The foundational background for nearly every local argument in these notes lies in EGA~IV, especially the sections on flatness, openness of flatness, dimensions of fibers, and noetherian approximation.  Reading the Raynaud--Gruson paper side by side with EGA is extremely instructive: one sees very clearly how the authors take results that in EGA appear as separate structural theorems and reorganize them into a single machine directed toward flattening.

For the local algebra of flatness, Lazard's work is indispensable.  Even when one only uses finite presentation modules in the end, the perspective that projective modules are filtered colimits of finite free modules is conceptually clarifying.  It explains why many statements proved first for free modules continue to hold for suitably controlled projective ones.

For a modern encyclopedic reference, the Stacks Project is particularly useful.  It provides separate entries on universal injectivity, associated points, Fitting ideals, generic flatness, blow-ups, and flattening techniques.  What it does not always convey as strongly as the original paper is the unity of the method.  That unity is best appreciated by reading the 1971 article itself and by reconstructing the proof architecture in seminar form, which is precisely the aim of the present notes.

There are also several directions in which the present notes could be extended.  One natural extension is a much fuller account of the passage from schemes to algebraic spaces in the final part of the paper.  Another is a more systematic comparison between flattening by blow-up and the later theory of flattening stratifications in projective geometry and moduli theory.  A third is to work out in detail the relation with Rees algebras and Fitting ideals in concrete examples, showing exactly when those more elementary constructions already detect the centers that the Raynaud--Gruson method chooses intrinsically.

For a research seminar, one profitable exercise is to compare the theorem flat $+$ pure implies projective with other criteria for projectivity, for example local freeness criteria over regular bases or descent criteria under faithfully flat morphisms.  The comparison reveals that the real novelty of Raynaud and Gruson is not a new formal characterization of projectivity but a geometric one phrased in terms of associated points and purity.  This viewpoint remains powerful far beyond the original theorem.

Finally, readers interested in motivic or derived applications may view the flattening theorem as part of the general technology that allows one to improve geometric input before applying cohomological or categorical constructions.  In many modern arguments, one first replaces a scheme or a sheaf by a better birational model on which finiteness or flatness conditions hold uniformly.  The Raynaud--Gruson theorem is one of the classical sources of that philosophy.




\section{Coda: what one should retain after reading the paper}

If one puts aside the technical machinery and asks what the enduring mathematical lesson of the paper is, the answer is this: flatness becomes manageable only after one chooses the right geometric coordinates on the support, and once those coordinates are chosen, the failure of flatness can be repaired by a birational operation that is intrinsic to the module.  This lesson may sound broad, but it is worth spelling out concretely.

First, the support of a coherent module should not be treated as a passive subset on which the algebra happens to live.  In the Raynaud--Gruson theory the support is an active geometric object.  It determines the dimensions that drive the induction, it guides the choice of the free approximations in the d\'evissage, and through its associated points it records whether the module is pure enough to descend projectivity.

Second, exact sequences are not used here merely as bookkeeping devices.  Every exact sequence in the d\'evissage has a geometric meaning: the free term captures the top-dimensional generic behavior, the quotient records a lower-dimensional defect, and universal injectivity measures whether this comparison survives all changes of base.  One might say that the entire paper is an art of building exact sequences that remember geometry rather than forgetting it.

Third, the passage from local criteria to global birational improvement is a model of mathematical synthesis.  The local criterion of flatness provides the seed, but only after it is combined with purity, representability, and strict transform does one obtain a theorem of global scope.  The article is therefore an excellent example of how a major theorem is often assembled: not by one powerful lemma, but by a sequence of reorganizations of the same problem into successively more usable forms.

Finally, there is a methodological reason the paper rewards repeated reading.  Many foundational texts present a theorem and then isolate its applications.  Raynaud and Gruson instead present an entire toolkit whose parts immediately talk to each other.  The d\'evissage generates closed conditions; the closed conditions define centers of blow-up; the blow-up produces strict transforms; the strict transforms are tested again by the d\'evissage.  The theory is circular in the good sense: each construction feeds the next.  That recursive unity is what makes the paper both difficult and memorable.


\begin{thebibliography}{99}

\bibitem{EGAIV}
A.~Grothendieck and J.~Dieudonn\'e,
\emph{\'El\'ements de g\'eom\'etrie alg\'ebrique IV}, Publ. Math. IH\'ES.

\bibitem{RG}
M.~Raynaud and L.~Gruson,
\emph{Crit\`eres de platitude et de projectivit\'e. Techniques de \enquote{platification} d'un module},
Invent. Math. \textbf{13} (1971), 1--89.

\bibitem{Lazard}
D.~Lazard,
\emph{Autour de la platitude}, Bull. Soc. Math. France.

\bibitem{Stacks}
The Stacks Project Authors,
\emph{The Stacks Project}.

\end{thebibliography}

\end{document}
